% La gestion de l'énergie — Allemand 1ère A — Cours signé OnBuch+
% Programme officiel MINESEC. Compiler avec : tectonic ALA11-la-gestion-de-l-energie.tex
\newcommand{\DOCMATIERE}{Allemand}
\newcommand{\DOCNIVEAU}{1ère A}
\newcommand{\DOCMODULE}{Chapitre 3 : Environnement, santé et bien-être}
\newcommand{\DOCLECON}{ALA11}
\newcommand{\DOCTITRE}{La gestion de l'énergie}
% OnBuch+ — préambule « cours signés » ENRICHI (compilé avec tectonic / XeLaTeX)
% Système visuel « Pop » d'OnBuch+ : fond crème (jamais blanc), bordures encre
% épaisses, ombres « dures » décalées, titres Archivo Black en capitales.
% Version MATHÉMATIQUES : graphes (pgfplots/tikz), boîtes pédagogiques
% (définition, propriété, méthode, attention, le savais-tu, exercice résolu,
% exercice, TP, bilan), page de garde et signature OnBuch+ ancrée en bas.
%
% Macros attendues dans main.tex AVANT \input{preamble} :
%   \DOCMATIERE  \DOCNIVEAU  \DOCMODULE  \DOCLECON (numéro)  \DOCTITRE
\documentclass[11pt]{article}

\usepackage{fontspec}
\usepackage[ngerman,french]{babel} % français = langue principale ; l'allemand via \all{..} / textebox
\usepackage[a4paper,margin=2cm,top=3.4cm,bottom=2.8cm,headheight=1.4cm,footskip=1.3cm]{geometry}
\usepackage{amsmath,amssymb,amsfonts}
\usepackage{mathtools} % \xmapsto, \coloneqq... souvent utilisés par les modèles
\usepackage{cancel} % simplifications barrées (bilans de forces)
% \abs{z} (module, valeur absolue) et \norm{u} (norme), utilisés par les modèles
\providecommand{\abs}{}\let\abs\relax\DeclarePairedDelimiter{\abs}{\lvert}{\rvert}
\providecommand{\norm}{}\let\norm\relax\DeclarePairedDelimiter{\norm}{\lVert}{\rVert}
\usepackage[table]{xcolor} % [table] : fournit aussi \rowcolors (alternance de lignes)
\usepackage{pagecolor}
\usepackage{tikz}
\usetikzlibrary{arrows.meta,positioning,calc,shapes.geometric,shapes.misc,shapes.arrows,decorations.pathmorphing,decorations.pathreplacing,decorations.markings,patterns,shadows,mindmap,trees,fit,backgrounds,matrix,intersections,angles,quotes}
\usepackage{pgfplots}
\usepackage[version=4]{mhchem} % équations nucléaires \ce{^{235}_{92}U}
\usepackage[european,siunitx]{circuitikz} % schémas électriques
\pgfplotsset{compat=1.18}
\usepgfplotslibrary{fillbetween,groupplots} % aires entre courbes (calcul intégral)
\usepackage{enumitem}
\usepackage{fancyhdr}
% [most] charge toutes les bibliothèques tcolorbox (listings, minted, raster,
% documentation...) et gonfle inutilement la mémoire TeX sur un document long
% (~300 boîtes) : on ne charge que ce qui est réellement utilisé.
\usepackage[skins,breakable]{tcolorbox}
\usepackage{graphicx}
\usepackage{colortbl}
\usepackage{booktabs}
\usepackage{tabularx}
\usepackage{diagbox} % case d'en-tête diagonale des tableaux à double entrée
% Colonnes extensibles centrée / gauche / droite (C, L, R), souvent utilisées par les modèles
\newcolumntype{C}{>{\centering\arraybackslash}X}
\newcolumntype{L}{>{\raggedright\arraybackslash}X}
\newcolumntype{R}{>{\raggedleft\arraybackslash}X}
\usepackage{array}
\usepackage{multicol}
\usepackage{multirow}
\usepackage{siunitx}
% parse-numbers=false : accepte aussi \SI{2,5 \times 10^{-3}}{...}
\sisetup{locale=FR,output-decimal-marker={,},group-minimum-digits=4,parse-numbers=false}
\DeclareSIUnit\ohm{\ensuremath{\symup{\Omega}}} % U+2126 absent de la police
\DeclareSIUnit\division{div} % réglages d'oscilloscope (ms/div, V/div)
\DeclareSIUnit\tr{tr} % tours (vitesse de rotation, ex. \SI{1500}{\tr\per\minute})
\DeclareSIUnit\molar{M} % concentration molaire (ex. \SI{3}{\molar}, \SI{1}{\micro\molar})
\DeclareSIUnit\an{an} % année (le modèle confond parfois avec \year, un compteur TeX) — ex. \SI{5}{\kilo\meter\per\an}
\DeclareSIUnit\FCFA{FCFA} % franc CFA (ex. \SI{500}{\FCFA\per\kilo\gram})
\DeclareSIUnit\degreeBrix{\textdegree Brix} % teneur en sucre d'un sirop/jus (réfractométrie)
\DeclareSIUnit\month{mois} % le modèle utilise parfois \month, un compteur TeX, comme unité de durée
\usepackage{qrcode}
% \degree : commande fréquemment utilisée par les modèles pour un angle ou une
% température en dehors de \SI{}{} (non fournie par les paquets chargés).
\providecommand{\degree}{^{\circ}}
% \qed (fin de démonstration), utilisé par les modèles sans amsthm
\providecommand{\qed}{\hfill$\square$}
% \barre : double barre des valeurs interdites dans un tableau de variations (tkz-tab)
\providecommand{\barre}{\ensuremath{\|}}
% environnement proof (amsthm non chargé) utilisé par les modèles
\ifdefined\proof\else\newenvironment{proof}[1][Démonstration]{\par\noindent\textit{#1.}\ }{\qed\par}\fi
\usepackage{lastpage}
\usepackage{hyperref}
\hypersetup{hidelinks}

% ============================================================
% Palette « Pop » OnBuch+ — mêmes hex que la classe P (plus_ui.dart)
% ============================================================
\definecolor{poporange}{HTML}{F59321}
\definecolor{popdark}{HTML}{E07A0C}
\definecolor{popcream}{HTML}{FFF6E8}
\definecolor{popink}{HTML}{1C1714}
\definecolor{poppurple}{HTML}{7A5AE0}
\definecolor{popgreen}{HTML}{2E9E6A}
\definecolor{poppink}{HTML}{E0457B}
\definecolor{popblue}{HTML}{2F7DE1}
\definecolor{popgold}{HTML}{C98A12}
\definecolor{popmuted}{HTML}{8A7F76}
\definecolor{popline}{HTML}{EADFD2}
\definecolor{popgray}{HTML}{8A7F76}
\colorlet{popgrey}{popgray}
% Alias tolérants (le modèle nomme parfois les couleurs différemment)
\colorlet{popwhite}{white}
\colorlet{popblack}{popink}
\colorlet{popred}{poppink}
\colorlet{popyellow}{popgold}
% Teintes claires (fonds de tableaux/encadrés) — le modèle nomme parfois
% les couleurs avec un suffixe L (ex. popblueL) sans les avoir définies.
\colorlet{poporangeL}{poporange!20}
\colorlet{popdarkL}{popdark!20}
\colorlet{poppurpleL}{poppurple!20}
\colorlet{popgreenL}{popgreen!20}
\colorlet{poppinkL}{poppink!20}
\colorlet{popblueL}{popblue!20}
\colorlet{popgoldL}{popgold!20}
\colorlet{popredL}{popred!20}
\colorlet{popgrayL}{popgray!20}
% Teintes claires pour fonds de tableaux / zones
\colorlet{popblueL}{popblue!10!white}
\colorlet{poporangeL}{poporange!14!white}
\colorlet{popgreenL}{popgreen!12!white}
\colorlet{poppurpleL}{poppurple!12!white}
\colorlet{poppinkL}{poppink!12!white}
\colorlet{popgoldL}{popgold!14!white}

\pagecolor{popcream}

% ============================================================
% Typographie
% ============================================================
\defaultfontfeatures{Ligatures=TeX}
\setmainfont{PlusJakartaSans-Medium.ttf}[Path=../../../fonts/,BoldFont=PlusJakartaSans-Bold.ttf]
% Maths en sans-serif (Fira Math) avec chiffres et lettres droites de Plus
% Jakarta, pour que les formules se fondent dans le texte.
\usepackage[math-style=ISO,bold-style=ISO]{unicode-math}
\setmathfont{FiraMath-Regular.otf}[Path=../../../fonts/]
\setmathfont{PlusJakartaSans-Medium.ttf}[Path=../../../fonts/,range={up/num,up/{Latin,latin},\mathup}]
\usepackage{icomma} % virgule décimale sans espace en mode math
% Caractères mathématiques Unicode absents des polices de texte (Plus Jakarta,
% Archivo Black) : rendus par leur équivalent mathématique, en texte comme en
% formule — sinon ils s'affichent en carré vide (ex. « Arithmétique dans ℤ »).
\usepackage{newunicodechar}
\newunicodechar{ℝ}{\ensuremath{\mathbb{R}}}
\newunicodechar{ℤ}{\ensuremath{\mathbb{Z}}}
\newunicodechar{ℂ}{\ensuremath{\mathbb{C}}}
\newunicodechar{ℕ}{\ensuremath{\mathbb{N}}}
\newunicodechar{ℚ}{\ensuremath{\mathbb{Q}}}
\newunicodechar{𝔻}{\ensuremath{\mathbb{D}}}
\newunicodechar{α}{\ensuremath{\alpha}}
\newunicodechar{β}{\ensuremath{\beta}}
\newunicodechar{γ}{\ensuremath{\gamma}}
\newunicodechar{δ}{\ensuremath{\delta}}
\newunicodechar{ε}{\ensuremath{\varepsilon}}
\newunicodechar{θ}{\ensuremath{\theta}}
\newunicodechar{λ}{\ensuremath{\lambda}}
\newunicodechar{μ}{\ensuremath{\mu}}
\newunicodechar{σ}{\ensuremath{\sigma}}
\newunicodechar{τ}{\ensuremath{\tau}}
\newunicodechar{φ}{\ensuremath{\varphi}}
\newunicodechar{ψ}{\ensuremath{\psi}}
\newunicodechar{ω}{\ensuremath{\omega}}
\newunicodechar{Σ}{\ensuremath{\Sigma}}
% majuscules produites par \MakeUppercase dans les titres (α-aminés -> Α)
\newunicodechar{Α}{\ensuremath{\alpha}}
\newunicodechar{Β}{\ensuremath{\beta}}
\newunicodechar{Γ}{\ensuremath{\Gamma}}
\newunicodechar{Δ}{\ensuremath{\Delta}}
\newunicodechar{Ε}{\ensuremath{\varepsilon}}
\newunicodechar{Θ}{\ensuremath{\Theta}}
\newunicodechar{Λ}{\ensuremath{\Lambda}}
\newunicodechar{Μ}{\ensuremath{\mu}}
\newunicodechar{Τ}{\ensuremath{\tau}}
\newunicodechar{Φ}{\ensuremath{\Phi}}
\newunicodechar{Ψ}{\ensuremath{\Psi}}
\newunicodechar{Ω}{\ensuremath{\Omega}}
\newunicodechar{↦}{\ensuremath{\mapsto}}
\newunicodechar{∈}{\ensuremath{\in}}
\newunicodechar{∉}{\ensuremath{\notin}}
\newunicodechar{∧}{\ensuremath{\wedge}}
\newunicodechar{⇔}{\ensuremath{\Leftrightarrow}}
\newunicodechar{⟺}{\ensuremath{\Longleftrightarrow}}
\newunicodechar{⇒}{\ensuremath{\Rightarrow}}
\newunicodechar{⟹}{\ensuremath{\Longrightarrow}}
\newunicodechar{∘}{\ensuremath{\circ}}
\newunicodechar{∪}{\ensuremath{\cup}}
\newunicodechar{∩}{\ensuremath{\cap}}
\newunicodechar{≡}{\ensuremath{\equiv}}
\newunicodechar{⊂}{\ensuremath{\subset}}
\newunicodechar{⊥}{\ensuremath{\perp}}
\newunicodechar{∃}{\ensuremath{\exists}}
\newunicodechar{∀}{\ensuremath{\forall}}
\newunicodechar{∛}{\ensuremath{\sqrt[3]{\,}}}
\newunicodechar{∜}{\ensuremath{\sqrt[4]{\,}}}
\newunicodechar{⃗}{\ensuremath{{}^{\to}}}
\newunicodechar{ⁿ}{\ifmmode^{n}\else\textsuperscript{\ensuremath{n}}\fi}
\newunicodechar{ˣ}{\ifmmode^{x}\else\textsuperscript{\ensuremath{x}}\fi}
\newunicodechar{ᵇ}{\ifmmode^{b}\else\textsuperscript{\ensuremath{b}}\fi}
\newunicodechar{ʳ}{\ifmmode^{r}\else\textsuperscript{\ensuremath{r}}\fi}
\newunicodechar{ᵘ}{\ifmmode^{u}\else\textsuperscript{\ensuremath{u}}\fi}
\newunicodechar{ʸ}{\ifmmode^{y}\else\textsuperscript{\ensuremath{y}}\fi}
\newunicodechar{ᵃ}{\ifmmode^{a}\else\textsuperscript{\ensuremath{a}}\fi}
\newunicodechar{ᵉ}{\ifmmode^{e}\else\textsuperscript{\ensuremath{e}}\fi}
\newunicodechar{ᵏ}{\ifmmode^{k}\else\textsuperscript{\ensuremath{k}}\fi}
\newunicodechar{ᵖ}{\ifmmode^{p}\else\textsuperscript{\ensuremath{p}}\fi}
\newunicodechar{ᴺ}{\ifmmode^{N}\else\textsuperscript{\ensuremath{N}}\fi}
\newunicodechar{ᶜ}{\ifmmode^{c}\else\textsuperscript{\ensuremath{c}}\fi}
\newunicodechar{ʰ}{\ifmmode^{h}\else\textsuperscript{\ensuremath{h}}\fi}
\newunicodechar{ᵠ}{\ifmmode^{q}\else\textsuperscript{\ensuremath{q}}\fi}
\newunicodechar{ₙ}{\ifmmode_{n}\else\textsubscript{\ensuremath{n}}\fi}
\newunicodechar{ₐ}{\ifmmode_{a}\else\textsubscript{\ensuremath{a}}\fi}
\newunicodechar{ᵢ}{\ifmmode_{i}\else\textsubscript{\ensuremath{i}}\fi}
\newunicodechar{ᵣ}{\ifmmode_{r}\else\textsubscript{\ensuremath{r}}\fi}
\newunicodechar{ₖ}{\ifmmode_{k}\else\textsubscript{\ensuremath{k}}\fi}
\newunicodechar{ᵦ}{\ifmmode_{\beta}\else\textsubscript{\ensuremath{\beta}}\fi}
\newunicodechar{ₓ}{\ifmmode_{x}\else\textsubscript{\ensuremath{x}}\fi}
\newfontfamily\popdisplay{ArchivoBlack-Regular.ttf}[Path=../../../fonts/]
\newfontfamily\poplogo{SpaceGrotesk-Bold.ttf}[Path=../../../fonts/]
\newfontfamily\popsemi{PlusJakartaSans-SemiBold.ttf}[Path=../../../fonts/]
\newfontfamily\popxbold{PlusJakartaSans-ExtraBold.ttf}[Path=../../../fonts/]
\newfontfamily\popxboldls{PlusJakartaSans-ExtraBold.ttf}[Path=../../../fonts/,LetterSpace=3.0]

\setlist[enumerate]{leftmargin=1.7em,itemsep=5pt,topsep=4pt}
\setlist[itemize]{leftmargin=1.7em,itemsep=4pt,topsep=4pt,label={\color{poporange}\textbullet}}
\setlength{\parindent}{0pt}
\setlength{\parskip}{5pt}
\renewcommand{\arraystretch}{1.25}

% pgfplots : style Pop par défaut
\pgfplotsset{
  every axis/.append style={axis line style={popink,line width=1pt},tick style={popink},
    grid style={popline,line width=0.5pt},label style={font=\small},tick label style={font=\footnotesize}},
  popcurve/.style={poporange,line width=1.8pt,no markers,smooth},
}
% Même style disponible dans un \draw TikZ simple (hors pgfplots)
\tikzset{popcurve/.style={poporange,line width=1.8pt}}
\tikzset{popgrid/.style={popline,very thin}}
\colorlet{popcurve}{poporange} % le nom du style est parfois utilisé comme couleur

% ============================================================
% Pill / squiggle
% ============================================================
\newcommand{\poppill}[3]{%
  \tikz[baseline=-0.55ex]\node[draw=popink,line width=1.2pt,rounded corners=2.6pt,%
    fill=#2,inner xsep=6.5pt,inner ysep=3pt]{%
    \popxboldls\fontsize{7}{7}\selectfont\color{#3}\MakeUppercase{#1}};%
}
\newcommand{\popsquiggle}[1]{%
  \tikz[baseline=-0.4ex]{
    \draw[#1,line width=2.6pt,line cap=round]
      plot [smooth] coordinates {(0,0.09) (0.34,0.01) (0.68,0.21) (1.02,0.01) (1.36,0.10)};
  }%
}
% Mot-clé surligné (marqueur orange sous le texte)
\newcommand{\cle}[1]{{\popxbold\color{popdark}#1}}

% ============================================================
% En-tête / pied
% ============================================================
\pagestyle{fancy}
\fancyhf{}
\renewcommand{\headrulewidth}{0pt}
\renewcommand{\footrulewidth}{0pt}
\lhead{\raisebox{-0.15ex}{\poplogo\fontsize{13}{13}\selectfont\color{popink}OnBuch\color{poporange}+}}
\rhead{\poppill{\DOCMATIERE\ \textbullet\ \DOCNIVEAU\ \textbullet\ Leçon \DOCLECON}{white}{popink}}
\lfoot{\popxbold\fontsize{7.6}{7.6}\selectfont\color{popmuted}\MakeUppercase{Cours signé OnBuch+}\ \textbullet\ {\popsemi\DOCTITRE}}
\rfoot{\tikz[baseline=-0.6ex]\node[rounded corners=8.5pt,fill=popink,minimum height=17pt,minimum width=17pt,inner xsep=5pt,inner sep=0]{\popxbold\fontsize{8}{8}\selectfont\color{popcream}\thepage\,/\,\pageref*{LastPage}};}

% ============================================================
% Titres
% ============================================================
\newcommand{\chaptitre}[2]{%
  \begin{tcolorbox}[enhanced,colback=poporange,colframe=popink,boxrule=2.6pt,arc=11pt,
      drop shadow=popink,left=15pt,right=15pt,top=12pt,bottom=11pt,before skip=3pt,after skip=18pt]
    \poppill{#2}{popink}{popcream}\par\vspace{9pt}
    {\raggedright\hyphenpenalty=10000\exhyphenpenalty=10000\popdisplay\fontsize{22}{24}\selectfont\color{popcream}\MakeUppercase{#1}\par}
  \end{tcolorbox}
}
% Section numérotée : pastille orange avec le numéro + titre Archivo
\newcounter{popsec}
\newcommand{\coursec}[1]{%
  \stepcounter{popsec}\setcounter{popsub}{0}%
  \par\vspace{16pt}\noindent
  \begin{minipage}{\linewidth}
  \tikz[baseline=-0.7ex]\node[draw=popink,line width=1.6pt,fill=poporange,rounded corners=4pt,minimum size=22pt,inner sep=0]{\popdisplay\fontsize{12}{12}\selectfont\color{popcream}\thepopsec};%
  \hspace{8pt}{\popdisplay\fontsize{15}{17}\selectfont\color{popink}\MakeUppercase{#1}}\par
  \vspace{4pt}\noindent{\color{poporange}\rule{36pt}{3.6pt}}
  \end{minipage}\par\nopagebreak\vspace{8pt}
}
\newcounter{popsub}
\newcommand{\courssub}[1]{%
  \stepcounter{popsub}%
  \par\vspace{9pt}\noindent{\popxbold\fontsize{11.5}{13}\selectfont\color{popink}{\color{poporange}\thepopsec.\thepopsub}\hspace{6pt}#1}\par\nopagebreak\vspace{4pt}
}

% ============================================================
% Boîtes pédagogiques — même recette : fond blanc, bordure épaisse colorée,
% ombre dure encre, badge plein. Titre optionnel [#1].
% ============================================================
\tcbset{popbase/.style={breakable,enhanced,colback=white,boxrule=2.3pt,arc=9pt,
  shadow={3pt}{-3pt}{0pt}{popink},left=14pt,right=14pt,top=11pt,bottom=11pt,before skip=11pt,after skip=15pt,
  fontupper=\popsemi\fontsize{10.6}{14.5}\selectfont\color{popink},
  fontlower=\popsemi\fontsize{10.6}{14.5}\selectfont\color{popink}}}
% Coche dessinée (le glyphe ✓ n'existe pas dans les polices du document)
\newcommand{\popcheck}[1]{\tikz[baseline=-0.5ex]\draw[#1,line width=1.6pt,line cap=round,line join=round](-0.13,0.0)--(-0.04,-0.09)--(0.13,0.1);}
\providecommand{\coche}{\popcheck{popgreen}}
\providecommand{\resultat}[1]{\par\smallskip\noindent\fcolorbox{popgreen}{popgreenL}{\parbox{\dimexpr\linewidth-2\fboxsep-2\fboxrule\relax}{\textbf{Résultat.} #1}}\par\smallskip}
\newcommand{\popboxhead}[3]{\poppill{#1}{#2}{white}\ifx\relax#3\relax\else\hspace{6pt}{\popxbold\fontsize{10.6}{12}\selectfont\color{popink}#3}\fi\par\vspace{7pt}}

\newtcolorbox{aretenir}[1][]{popbase,colframe=popgreen,before upper={\popboxhead{À retenir}{popgreen}{#1}}}
% Texte en allemand (document de lecture, dialogue, extrait) : cadre
% d'étude. Un titre optionnel (source, auteur) s'affiche dans l'en-tête.
\newtcolorbox{textebox}[1][]{popbase,colframe=popblue,before upper={\popboxhead{Text}{popblue}{#1}\begin{otherlanguage}{ngerman}},after upper={\end{otherlanguage}}}
% Marques ✓ / ✗ (forme correcte / fautive) souvent utilisées par les modèles
\providecommand{\cross}{\textcolor{poppink}{\ensuremath{\times}}}
\providecommand{\xmark}{\cross}\providecommand{\crossmark}{\cross}\providecommand{\cmark}{\ensuremath{\checkmark}}
% Ligne de réponse / justification à compléter (inventée par les modèles)
\providecommand{\justification}[1]{\par\noindent\textit{Justification~:} \dotfill\par}
% \textipa (package tipa non chargé) : la police ne couvre pas l'API -> simple passage
\providecommand{\textipa}[1]{#1}
% Soulignement (ulem non chargé) : \uline, \ul
\providecommand{\uline}[1]{\underline{#1}}\providecommand{\ul}[1]{\underline{#1}}
% \alert (beamer) inventée par les modèles : texte mis en évidence
\providecommand{\alert}[1]{\textbf{\textcolor{poppink}{#1}}}
% variantes ASCII d'ordinaux écrites en macro
\providecommand{\iere}{\textsuperscript{re}}\providecommand{\ieme}{\textsuperscript{e}}\providecommand{\ere}{\textsuperscript{re}}\providecommand{\eme}{\textsuperscript{e}}\providecommand{\er}{\textsuperscript{er}}\providecommand{\ie}{\textsuperscript{e}}
% Ordinaux français écrits en macro par les modèles : 1\ière, 3\ième
\newcommand{\ière}{\textsuperscript{re}}\newcommand{\ième}{\textsuperscript{e}}
% \exemple (inventée par les modèles) : étiquette « Exemple : »
\providecommand{\exemple}{\textbf{Exemple~:}\ }
% \square utilisable aussi hors mode math (cases à cocher)
\AtBeginDocument{\let\squareMath\square\renewcommand{\square}{\ensuremath{\squareMath}}}
% Mot ou phrase en allemand à l'intérieur d'un paragraphe français
\newcommand{\all}[1]{\ifmmode\text{\textit{#1}}\else\begingroup\selectlanguage{ngerman}\itshape #1\endgroup\fi}
\let\esp\all % alias toléré
\newtcolorbox{exemplebox}[1][]{popbase,colframe=poporange,before upper={\popboxhead{Exemple}{poporange}{#1}}}
\newtcolorbox{definition}[1][]{popbase,colframe=popblue,before upper={\popboxhead{Définition}{popblue}{#1}}}
\newtcolorbox{propriete}[1][]{popbase,colframe=poppurple,before upper={\popboxhead{Propriété}{poppurple}{#1}}}
\newtcolorbox{methode}[1][]{popbase,colframe=popgold,before upper={\popboxhead{Méthode}{popgold}{#1}}}
\newtcolorbox{attention}[1][]{popbase,colframe=poppink,before upper={\popboxhead{Attention}{poppink}{#1}}}
\newtcolorbox{experience}[1][]{popbase,colframe=popblue,colback=popblueL,before upper={\popboxhead{Expérience}{popblue}{#1}}}
% « Le savais-tu ? » : carte violette pleine (contexte, culture, Cameroun)
\newtcolorbox{savaistu}[1][]{popbase,colback=poppurple,colframe=popink,
  fontupper=\popsemi\fontsize{10.4}{14.2}\selectfont\color{white},
  before upper={\poppill{Le savais-tu\,?}{white}{poppurple}\ifx\relax#1\relax\else\hspace{6pt}{\popxbold\color{white}#1}\fi\par\vspace{7pt}}}
% Exercice résolu : énoncé en haut, \tcblower puis la solution sur fond crème
\newcounter{popexo}\newcounter{popexr}
\newtcolorbox{exoresolu}[1][]{popbase,colframe=poporange,colbacklower=popcream,
  segmentation style={popink,line width=1.2pt,dashed},
  before upper={\stepcounter{popexr}\popboxhead{Exercice résolu \thepopexr}{poporange}{#1}},
  before lower={\poppill{Solution}{popink}{popcream}\par\vspace{6pt}}}
% Exercice (non résolu en place) — la correction est dans la partie Corrigés
\newtcolorbox{exercice}[1][]{popbase,colframe=popink,
  before upper={\stepcounter{popexo}\popboxhead{Exercice \thepopexo}{popink}{#1}}}
% Corrigé
\newtcolorbox{corrige}[1][]{popbase,colframe=popgreen,colback=popgreenL,
  before upper={\popboxhead{Corrigé}{popgreen}{#1}}}
% Objectifs / prérequis / bilan
\newtcolorbox{objectifs}{popbase,colframe=popink,colback=poporangeL,
  before upper={\popboxhead{Objectifs de la leçon}{popink}{}}}
\newtcolorbox{prerequis}{popbase,colframe=popink,colback=popblueL,
  before upper={\popboxhead{Prérequis}{popblue}{}}}
\newtcolorbox{bilan}{popbase,colframe=popink,colback=white,boxrule=2.8pt,
  before upper={\popboxhead{Fiche bilan}{poporange}{L'essentiel à connaître pour l'examen}}}
% Figure encadrée avec légende
\newtcolorbox{popfigure}[1][]{enhanced,breakable=false,colback=white,colframe=popline,boxrule=1.4pt,arc=8pt,
  left=10pt,right=10pt,top=10pt,bottom=8pt,before skip=10pt,after skip=12pt,halign=center,
  lower separated=false,halign lower=center,
  fontlower=\popsemi\fontsize{9}{11.5}\selectfont\color{popmuted},
  #1}
\newcommand{\legende}[1]{\par\vspace{6pt}{\popsemi\fontsize{9}{11.5}\selectfont\color{popmuted}#1}}

% ============================================================
% Signature OnBuch+
% ============================================================
\newcommand{\poplogoinline}[1]{{\poplogo\fontsize{#1}{#1}\selectfont\color{popink}OnBuch\color{poporange}+}}
\newcommand{\signatureonbuch}{%
  \begin{tcolorbox}[enhanced,colback=popink,colframe=popink,boxrule=2.2pt,arc=10pt,
      drop shadow=poporange,left=14pt,right=14pt,top=10pt,bottom=10pt,before skip=0pt,after skip=0pt]
    \tikz[baseline=-0.7ex]\node[circle,fill=popgreen,draw=popcream,line width=1.3pt,minimum size=22pt,inner sep=0]{\popcheck{white}};%
    \hspace{9pt}\begin{minipage}[c]{0.62\linewidth}
      {\popxbold\fontsize{12}{13}\selectfont\color{popcream}Signé\ \textbullet\ {\poplogo OnBuch\color{poporange}+}}\par\vspace{2pt}
      {\raggedright\popsemi\fontsize{8}{10.5}\selectfont\color{popline}\MakeUppercase{Équipe pédagogique OnBuch+}\\\MakeUppercase{Conforme au programme officiel MINESEC}\par}
    \end{minipage}\hfill
    {\poplogo\fontsize{22}{22}\selectfont\color{popcream}OnBuch\color{poporange}+}
  \end{tcolorbox}%
}

% ============================================================
% Page de garde
% #1 titre · #2 matière · #3 niveau · #4 module · #5 n° leçon · #6 durée ·
% #7 description / objectifs (texte ou itemize)
% ============================================================
\newcommand{\pagegarde}[7]{%
  \thispagestyle{empty}
  \newgeometry{a4paper,margin=1.7cm,top=1.6cm,bottom=1.4cm}%
  \begin{tikzpicture}[remember picture,overlay]
    \fill[poporange!18] ([xshift=-3.2cm,yshift=-3.6cm]current page.north east) circle (4.6cm);
    \fill[poppurple!14] ([xshift=2.4cm,yshift=7.6cm]current page.south west) circle (3.8cm);
  \end{tikzpicture}%
  {\poplogo\fontsize{30}{30}\selectfont\color{popink}OnBuch\color{poporange}+}\hfill
  \raisebox{0.6ex}{\poppill{Cours signé \textbullet\ Premium}{popink}{popcream}}\par
  \vspace{1pt}\popsquiggle{poppurple}\par
  \vspace{8pt}
  \begin{tcolorbox}[enhanced,colback=poporange,colframe=popink,boxrule=2.8pt,arc=13pt,
      drop shadow=popink,left=18pt,right=18pt,top=12pt,bottom=13pt,after skip=10pt]
    \poppill{#2 \textbullet\ #3}{popink}{popcream}\hspace{5pt}\poppill{#4}{white}{popink}\par\vspace{8pt}
    {\popdisplay\fontsize{14}{14}\selectfont\color{popink}LEÇON #5}\par\vspace{4pt}
    {\raggedright\hyphenpenalty=10000\exhyphenpenalty=10000\popdisplay\fontsize{25}{27}\selectfont\color{popcream}\MakeUppercase{#1}\par}
    \vspace{8pt}
    {\popxbold\fontsize{9.5}{9.5}\selectfont\color{popink}\MakeUppercase{Durée conseillée : #6}}
  \end{tcolorbox}
  \begin{tcolorbox}[enhanced,colback=white,colframe=popink,boxrule=2.2pt,arc=10pt,
      drop shadow=popink,left=16pt,right=16pt,top=10pt,bottom=10pt,after skip=8pt]
    \poppill{Dans cette leçon}{poppurple}{white}\par\vspace{5pt}
    \popsemi\fontsize{9.3}{12.6}\selectfont\color{popink}#7
  \end{tcolorbox}
  \noindent\begin{minipage}[t]{0.487\linewidth}
    \begin{tcolorbox}[enhanced,colback=popgreen,colframe=popink,boxrule=2.2pt,arc=10pt,
        drop shadow=popink,left=11pt,right=11pt,top=10pt,bottom=10pt,height=3.1cm,valign=center]
      \tcbox[colback=white,colframe=popink,boxrule=1.3pt,arc=5pt,boxsep=0pt,nobeforeafter,
        left=3pt,right=3pt,top=3pt,bottom=3pt,valign=center]{\qrcode[height=1.9cm]{https://play.google.com/store/apps/details?id=com.luvvix.onbuch&hl=fr}}%
      \hspace{8pt}\begin{minipage}[c]{0.5\linewidth}
        {\popdisplay\fontsize{11}{12.5}\selectfont\color{white}\MakeUppercase{Télécharge OnBuch}}\par\vspace{4pt}
        {\raggedright\popsemi\fontsize{8.4}{11}\selectfont\color{white}Gratuit \textbullet\ Google Play\\Tuteur IA, annales, concours, résultats\par}
      \end{minipage}
    \end{tcolorbox}
  \end{minipage}\hfill
  \begin{minipage}[t]{0.487\linewidth}
    \begin{tcolorbox}[enhanced,colback=poppurple,colframe=popink,boxrule=2.2pt,arc=10pt,
        drop shadow=popink,left=11pt,right=11pt,top=10pt,bottom=10pt,height=3.1cm,valign=center]
      \tikz[baseline=-0.7ex]\node[circle,fill=white,draw=popink,line width=1.3pt,minimum size=1.6cm,inner sep=0]{\popdisplay\fontsize{24}{24}\selectfont\color{poppurple}+};%
      \hspace{8pt}\begin{minipage}[c]{0.6\linewidth}
        {\popdisplay\fontsize{11}{12.5}\selectfont\color{white}\MakeUppercase{Passe à OnBuch+}}\par\vspace{4pt}
        {\raggedright\popsemi\fontsize{8.4}{11}\selectfont\color{white}Tous les cours signés, Tuteur IA illimité, fiches PDF — onglet OnBuch+ de l'app\par}
      \end{minipage}
    \end{tcolorbox}
  \end{minipage}
  \vspace{8pt}
  \popcontenu
  \vfill
  \signatureonbuch
  \restoregeometry
  \newpage
}

% Rangée « Ce cours contient » (page de garde)
\newcommand{\poptile}[3]{% #1 couleur #2 gros texte #3 légende
  \begin{tcolorbox}[enhanced,colback=#1,colframe=popink,boxrule=2pt,arc=9pt,shadow={2.5pt}{-2.5pt}{0pt}{popink},
      left=6pt,right=6pt,top=9pt,bottom=9pt,halign=center,valign=center,height=1.75cm,width=\linewidth,nobeforeafter]
    {\popdisplay\fontsize{12.5}{13}\selectfont\color{white}\MakeUppercase{#2}}\par\vspace{3pt}
    {\centering\popsemi\fontsize{7.8}{9.5}\selectfont\color{white}#3\par}
  \end{tcolorbox}}
\newcommand{\popcontenu}{%
  \par\noindent{\popxboldls\fontsize{7.5}{7.5}\selectfont\color{popmuted}CE COURS SIGNÉ CONTIENT}\par\vspace{7pt}
  \noindent\begin{minipage}[t]{0.235\linewidth}\poptile{popblue}{Cours}{complet, illustré, pas à pas}\end{minipage}\hfill
  \begin{minipage}[t]{0.235\linewidth}\poptile{poporange}{Méthodes}{et exercices résolus}\end{minipage}\hfill
  \begin{minipage}[t]{0.235\linewidth}\poptile{poppink}{Exercices}{type examen + corrigés}\end{minipage}\hfill
  \begin{minipage}[t]{0.235\linewidth}\poptile{popgreen}{Fiche bilan}{l'essentiel à retenir}\end{minipage}\par}

% Sommaire
\newtcolorbox{sommaire}{enhanced,colback=white,colframe=popink,boxrule=2.2pt,arc=10pt,
  drop shadow=popink,left=18pt,right=18pt,top=13pt,bottom=13pt,before skip=6pt,after skip=20pt,
  before upper={\poppill{Sommaire}{poppurple}{white}\par\vspace{9pt}\popsemi\fontsize{10.6}{14.5}\selectfont\color{popink}}}

% Signature de fin de cours — poussée tout en bas de la dernière page
\newcommand{\findecours}{%
  \par\vfill\nopagebreak
  \begin{center}\popsquiggle{poporange}\end{center}\vspace{4pt}
  \signatureonbuch
}
\AtBeginDocument{\def\llbracket{\mathopen{[\mkern-3.5mu[}}\def\rrbracket{\mathclose{]\mkern-3.5mu]}}} % intervalles d'entiers (glyphe ⟦ absent de la police)
% Glyphes absents de Fira Math : redéfinis à partir de glyphes disponibles
\AtBeginDocument{%
  \def\setminus{\mathbin{\backslash}}\let\smallsetminus\setminus
  \def\vdots{\mathord{\vbox{\baselineskip3.2pt\lineskiplimit0pt\kern4pt\hbox{.}\hbox{.}\hbox{.}}}}%
  \def\ddots{\mathinner{\mkern1mu\raise6.5pt\hbox{.}\mkern2mu\raise3.6pt\hbox{.}\mkern2mu\raise0.7pt\hbox{.}\mkern1mu}}%
  \def\triangle{\mathord{\scalebox{0.9}{$\symup{\Delta}$}}}%
  \def\nabla{\mathord{\raisebox{\depth}{\scalebox{1}[-1]{$\symup{\Delta}$}}}}%
  \def\checkmark{\popcheck{popdark}}%
  \def\star{\mathbin{\ast}}\def\bigstar{\mathord{\ast}}%
  \def\diamond{\mathbin{\scalebox{0.6}{$\Diamond$}}}%
  \def\bigcup{\mathop{\mathchoice{\raisebox{-0.3em}{\scalebox{1.7}{$\cup$}}}{\scalebox{1.25}{$\cup$}}{\cup}{\cup}}\displaylimits}%
  \def\bigcap{\mathop{\mathchoice{\raisebox{-0.3em}{\scalebox{1.7}{$\cap$}}}{\scalebox{1.25}{$\cap$}}{\cap}{\cap}}\displaylimits}%
  \def\lesseqqgtr{\mathrel{\lessgtr}}\def\bigoplus{\mathop{\oplus}\displaylimits}%
}
\providecommand{\ssi}{si et seulement si} % abréviation fréquente
\AtBeginDocument{\providecommand{\wideparen}{\overparen}} % arc de cercle

\begin{document}

\pagegarde{La gestion de l'énergie}{Allemand}{1ère A}{Chapitre 3 : Environnement, santé et bien-être}{ALA11}{2 h}{Cette leçon propose la lecture et l'exploitation d'un texte allemand sur la gestion de l'énergie et les économies d'électricité. Elle permet d'acquérir le vocabulaire thématique de l'énergie et de maîtriser les verbes de modalité ainsi que l'expression de la finalité.
\vspace{4pt}
\begin{multicols}{2}\small
\begin{itemize}
\item À la fin de cette leçon, je sais comprendre globalement puis en détail un texte allemand sur la gestion de l'énergie.
\item Je sais employer correctement le vocabulaire du thème : die Energie, der Strom, die Sonne, das Wasserkraftwerk, sparen, die Steckdose, avec articles et pluriels.
\item Je sais reconnaître et conjuguer les verbes de modalité (können, müssen, dürfen, sollen, wollen, mögen) au présent et au prétérit.
\item Je sais construire une phrase exprimant la finalité avec um ... zu + infinitif et damit + proposition.
\item Je sais répondre par écrit et à l'oral à des questions de compréhension sur un texte.
\item Je sais résumer un texte en quelques phrases en réutilisant son lexique.
\end{itemize}\end{multicols}}

\begin{sommaire}
\begin{multicols}{2}
\begin{enumerate}
\item Texte support : Energie sparen im Alltag
\item Compréhension du texte
\item Lexique thématique : die Energie
\item Grammaire 1 : les verbes de modalité
\item Grammaire 2 : la finalité (um ... zu / damit)
\item Méthode du commentaire et production guidée
\item Activité d'intégration
\item Méthodes et exercices résolus
\item Exercices
\item Corrigés des exercices
\item Fiche bilan
\end{enumerate}
\end{multicols}
\end{sommaire}

\begin{objectifs}
\begin{itemize}
\item À la fin de cette leçon, je sais comprendre globalement puis en détail un texte allemand sur la gestion de l'énergie.
\item Je sais employer correctement le vocabulaire du thème : die Energie, der Strom, die Sonne, das Wasserkraftwerk, sparen, die Steckdose, avec articles et pluriels.
\item Je sais reconnaître et conjuguer les verbes de modalité (können, müssen, dürfen, sollen, wollen, mögen) au présent et au prétérit.
\item Je sais construire une phrase exprimant la finalité avec um ... zu + infinitif et damit + proposition.
\item Je sais répondre par écrit et à l'oral à des questions de compréhension sur un texte.
\item Je sais résumer un texte en quelques phrases en réutilisant son lexique.
\item Je sais produire un court paragraphe guidé sur les économies d'énergie au Cameroun.
\item Je sais éviter les erreurs fréquentes des francophones (place de l'infinitif, oubli de zu, confusion damit/um ... zu).
\end{itemize}
\end{objectifs}

\begin{prerequis}
\begin{itemize}
\item La conjugaison des verbes au présent (Seconde)
\item L'ordre des mots en proposition principale et subordonnée (Seconde)
\item Les articles définis et indéfinis et le pluriel des noms (Seconde)
\item Le vocabulaire de la maison et du quotidien (Seconde)
\item L'infinitif et la particule zu (début de Première)
\end{itemize}
\end{prerequis}

\newpage

\coursec{Texte support : Energie sparen im Alltag}

\courssub{Situation d'accroche et problématique}

À Yaoundé comme à Douala, les délestages rappellent chaque jour que l'énergie est une ressource précieuse : quand le courant se coupe au milieu du repas ou pendant les révisions du soir, chacun comprend que l'électricité n'est pas inépuisable. En Allemagne, le débat sur la sortie du nucléaire (\all{der Atomausstieg}) et les énergies renouvelables est au cœur de l'actualité : panneaux solaires sur les toits, éoliennes dans le nord du pays, factures d'électricité qui augmentent. Comment les Allemands gèrent-ils leur consommation d'électricité ? Que peut-on faire, chacun chez soi, pour économiser l'énergie ?

Ce texte nous plonge dans le quotidien d'une famille allemande, la famille Berger à Munich, soucieuse de réduire sa consommation. À travers elle, nous allons apprendre le vocabulaire de l'énergie et préparer deux notions de grammaire : les verbes de modalité et l'expression de la finalité.

\textbf{Problématique :} \emph{Comment une famille allemande peut-elle réduire sa consommation d'électricité au quotidien ?}

\courssub{Lecture globale du texte}

\begin{definition}[La lecture globale]
La \cle{lecture globale} est la première lecture d'un texte : elle vise à saisir le \textbf{sens général} sans s'arrêter sur les mots inconnus. On répond à quatre questions simples : \emph{qui ? où ? quel problème ? quelles solutions ?} Les mots inconnus seront étudiés lors de la lecture détaillée.
\end{definition}

\begin{methode}[Comment lire un texte allemand]
\begin{enumerate}
\item Lire le \textbf{titre} et anticiper : de quoi va parler le texte ?
\item Lire une première fois \textbf{sans dictionnaire}, en silence.
\item Repérer les \textbf{mots transparents} (proches du français) : \all{die Energie}, \all{die Sonne}, \all{das Problem}, \all{die Familie}.
\item Relire avec le \textbf{glossaire} pour comprendre les détails.
\item \textbf{Résumer} le texte en une phrase, en français ou en allemand.
\end{enumerate}
\end{methode}

\begin{textebox}[Energie sparen im Alltag]
Die Familie Berger aus München will weniger Strom verbrauchen. Sie muss die Rechnung jeden Monat bezahlen, und sie wird immer teurer. Deshalb schaltet der Vater das Licht ab, wenn er ein Zimmer verlässt. Die Tochter zieht den Stecker aus der Steckdose, um Energie zu sparen. Die Mutter kauft eine Solaranlage, damit die Familie die Sonne nutzen kann. In der Schule lernt der Sohn, dass Wasserkraftwerke Strom aus Wasser machen. „Wir können alle etwas tun“, sagt er, „wir sollen die Erde schützen!“
\end{textebox}

\textbf{Glossaire (mots nouveaux) :}

\begin{center}
\begin{tabularx}{\linewidth}{|X|X|}\hline
\rowcolor{popblueL} \textbf{Deutsch} & \textbf{Français} \\ \hline
der Strom (nur Singular) & l'électricité, le courant \\ \hline
verbrauchen & consommer \\ \hline
die Rechnung, -en & la facture \\ \hline
abschalten (schaltete ab, hat abgeschaltet) & éteindre, couper \\ \hline
der Stecker, - & la fiche (électrique) \\ \hline
die Steckdose, -n & la prise de courant \\ \hline
sparen & économiser, épargner \\ \hline
die Solaranlage, -n & l'installation solaire \\ \hline
nutzen & utiliser \\ \hline
das Wasserkraftwerk, -e & la centrale hydroélectrique \\ \hline
schützen & protéger \\ \hline
die Erde (nur Singular) & la Terre \\ \hline
\end{tabularx}
\end{center}

\begin{attention}[Faux ami : der Strom]
\all{der Strom} ne signifie pas « la tempête » ! \all{der Strom} = l'électricité, le courant (aussi : le grand fleuve). La tempête se dit \all{der Sturm}. De même, ne traduisez jamais mot à mot : \all{sie wird immer teurer} = « elle devient de plus en plus chère » (et non « elle sera toujours chère »). Ici, \all{wird} est le présent du verbe \all{werden} (devenir), pas le futur.
\end{attention}

\textbf{Repérage des informations générales (lecture globale) :}

\begin{center}
\begin{tabularx}{\linewidth}{|X|X|}\hline
\rowcolor{popblueL} \textbf{Question} & \textbf{Réponse} \\ \hline
Qui ? (Wer ?) & la famille Berger (le père, la mère, la fille, le fils) \\ \hline
Où ? (Wo ?) & à Munich, en Allemagne (et à l'école du fils) \\ \hline
Quel problème ? & la facture d'électricité devient de plus en plus chère \\ \hline
Quelles solutions ? & éteindre la lumière, débrancher les appareils, installer des panneaux solaires \\ \hline
\end{tabularx}
\end{center}

\courssub{Lecture détaillée phrase par phrase}

\textbf{Phrase 1 :} \all{Die Familie Berger aus München will weniger Strom verbrauchen.}

« La famille Berger, de Munich, veut consommer moins d'électricité. » Mots-clés à souligner : \cle{will} (verbe de modalité : veut), \cle{weniger} (moins), \cle{Strom verbrauchen} (consommer de l'électricité). Cette phrase annonce le sujet et l'intention de la famille. Observez la construction : le verbe modal \all{will} est conjugué en deuxième position, l'infinitif \all{verbrauchen} va à la fin de la phrase.

\textbf{Phrase 2 :} \all{Sie muss die Rechnung jeden Monat bezahlen, und sie wird immer teurer.}

« Elle doit payer la facture chaque mois, et elle devient de plus en plus chère. » Mots-clés : \cle{muss} (doit — verbe de modalité), \cle{die Rechnung} (la facture), \cle{immer teurer} (de plus en plus chère). C'est ici qu'apparaît le \emph{problème}. Notez : \all{jeden Monat} = « chaque mois » (accusatif de temps).

\textbf{Phrase 3 :} \all{Deshalb schaltet der Vater das Licht ab, wenn er ein Zimmer verlässt.}

« C'est pourquoi le père éteint la lumière quand il quitte une pièce. » Mots-clés : \cle{deshalb} (c'est pourquoi), \cle{abschalten} (éteindre — verbe à particule séparable : \all{ab} va à la fin de la proposition principale). Première solution concrète. Après \all{wenn} (quand), le verbe conjugué \all{verlässt} va à la fin de la subordonnée.

\textbf{Phrase 4 :} \all{Die Tochter zieht den Stecker aus der Steckdose, um Energie zu sparen.}

« La fille retire la fiche de la prise, pour économiser de l'énergie. » Mots-clés : \cle{der Stecker} (la fiche), \cle{die Steckdose} (la prise), \cle{um ... zu sparen} (pour économiser — construction de \textbf{finalité}, à retenir pour la section 5). Attention : \all{den Stecker} est à l'accusatif (complément d'objet direct) ; \all{aus der Steckdose} est au datif après la préposition \all{aus}.

\textbf{Phrase 5 :} \all{Die Mutter kauft eine Solaranlage, damit die Familie die Sonne nutzen kann.}

« La mère achète une installation solaire, afin que la famille puisse utiliser le soleil. » Mots-clés : \cle{die Solaranlage}, \cle{die Sonne nutzen}, \cle{damit ... kann} (afin que ... puisse — autre construction de \textbf{finalité}). Après \all{damit}, le verbe conjugué \all{kann} va à la fin, et l'infinitif \all{nutzen} se place juste avant lui.

\textbf{Phrases 6 et 7 :} \all{In der Schule lernt der Sohn, dass Wasserkraftwerke Strom aus Wasser machen. „Wir können alle etwas tun“, sagt er, „wir sollen die Erde schützen!“}

« À l'école, le fils apprend que les centrales hydroélectriques produisent de l'électricité à partir de l'eau. "Nous pouvons tous faire quelque chose, dit-il, nous devons protéger la Terre !" » Mots-clés : \cle{das Wasserkraftwerk} (la centrale hydroélectrique), \cle{können} (pouvoir), \cle{sollen} (devoir), \cle{die Erde schützen} (protéger la Terre). Après \all{dass} (que), le verbe conjugué \all{machen} va à la fin. Le texte se termine par un message d'espoir et de responsabilité collective.

\begin{exoresolu}[Compréhension globale]
Énoncé : Répondez en allemand, par une phrase complète : \emph{Warum will die Familie Berger weniger Strom verbrauchen?}
\tcblower
Solution : On cherche la cause dans la phrase 2 du texte : la facture devient de plus en plus chère. Réponse correcte : \all{Sie will weniger Strom verbrauchen, weil die Rechnung immer teurer wird.} « Elle veut consommer moins d'électricité parce que la facture devient de plus en plus chère. » Attention à la place du verbe : après \all{weil}, le verbe conjugué (\all{wird}) va à la fin de la phrase.
\end{exoresolu}

\courssub{Écoute et relecture à voix haute}

\begin{experience}[Écouter puis lire à voix haute]
\begin{enumerate}
\item \textbf{Écoute} : l'enseignant lit le texte à voix haute, les livres fermés. Les élèves notent uniquement les mots qu'ils reconnaissent.
\item \textbf{Deuxième écoute}, livres ouverts : chaque élève suit le texte du doigt et souligne les mots-clés.
\item \textbf{Relecture individuelle} : chaque élève lit une phrase à voix haute. L'enseignant corrige la prononciation.
\end{enumerate}
Points de prononciation à surveiller : \all{Strom} « chtrom » (le \all{St-} initial se prononce « cht »), \all{Steckdose} « chtek-do-zé », \all{Wasserkraftwerk} « va-seur-kraft-verk » (le \all{W} se prononce « v »), \all{München} « mu-n-heur », \all{verbrauchen} « fer-braou-rheur » (\all{au} se prononce « aou »).
\end{experience}

\courssub{Carte mentale du vocabulaire}

\begin{popfigure}
\begin{center}
\begin{tikzpicture}[mindmap, grow cyclic, every node/.style={concept, text=white, align=center, font=\small}, root concept/.append style={concept color=popblue, font=\bfseries\small}, level 1/.append style={level distance=3.2cm, sibling angle=72}]
\node[root concept]{Energie sparen}
  child[concept color=poporange]{ node[align=center]{der Strom\\die Steckdose\\der Stecker} }
  child[concept color=popgold]{ node[align=center]{die Sonne\\die Solaranlage} }
  child[concept color=popgreen]{ node[align=center]{das Wasser\\das Wasserkraftwerk} }
  child[concept color=poppurple]{ node[align=center]{zu Hause\\das Licht abschalten} }
  child[concept color=poppink]{ node[align=center]{in der Schule\\die Erde schützen} };
\end{tikzpicture}
\end{center}
\legende{Carte mentale : le vocabulaire de l'énergie autour du thème \all{Energie sparen}}
\end{popfigure}

\courssub{Repérage grammatical : préparation des sections 4 et 5}

Le texte contient des \textbf{verbes de modalité} (ils expriment la volonté, l'obligation, la possibilité) et des \textbf{constructions de finalité} (elles expriment le but). Repérez-les avant de les étudier en détail :

\begin{center}
\begin{tabularx}{\linewidth}{|X|X|X|}\hline
\rowcolor{popblueL} \textbf{Phrase du texte} & \textbf{Forme} & \textbf{Valeur} \\ \hline
\all{will weniger Strom verbrauchen} & wollen + infinitif & volonté \\ \hline
\all{muss die Rechnung bezahlen} & müssen + infinitif & obligation \\ \hline
\all{die Sonne nutzen kann} & können + infinitif & possibilité \\ \hline
\all{Wir können alle etwas tun} & können + infinitif & possibilité \\ \hline
\all{wir sollen die Erde schützen} & sollen + infinitif & devoir \\ \hline
\all{um Energie zu sparen} & um ... zu + infinitif & but (même sujet) \\ \hline
\all{damit die Familie ... kann} & damit + phrase & but (sujets différents) \\ \hline
\end{tabularx}
\end{center}

\begin{aretenir}[Deux constructions pour le but]
Pour exprimer le but (« pour », « afin de »), l'allemand utilise :
\begin{itemize}
\item \all{um ... zu} + infinitif quand le sujet est le même : \all{Sie zieht den Stecker aus der Steckdose, um Energie zu sparen.} « Elle retire la fiche de la prise pour économiser de l'énergie. »
\item \all{damit} + proposition (verbe conjugué à la fin) quand les sujets sont différents : \all{Die Mutter kauft eine Solaranlage, damit die Familie die Sonne nutzen kann.} « La mère achète une installation solaire afin que la famille puisse utiliser le soleil. »
\end{itemize}
\end{aretenir}

\begin{savaistu}[L'Allemagne et l'énergie]
L'Allemagne a fermé ses dernières centrales nucléaires en 2023 : c'est \all{der Atomausstieg} (la sortie du nucléaire). Le pays mise désormais sur les \all{erneuerbaren Energien} (les énergies renouvelables) : le soleil (\all{die Sonnenenergie}), le vent (\all{die Windenergie}) et l'eau (\all{die Wasserkraft}). Au Cameroun, l'hydroélectricité est aussi essentielle : les barrages sur le fleuve Sanaga fournissent une grande partie de l'électricité du pays.
\end{savaistu}

\begin{exercice}[À vous]
\begin{enumerate}
\item Résumez le texte en une phrase en français.
\item Soulignez dans le texte les cinq formes de verbes de modalité et donnez leur infinitif.
\item Traduisez : \all{Die Tochter zieht den Stecker aus der Steckdose, um Energie zu sparen.}
\item Répondez en allemand : \emph{Was macht der Vater, wenn er ein Zimmer verlässt?}
\end{enumerate}
\end{exercice}

\begin{corrige}[Exercice « À vous »]
\begin{enumerate}
\item La famille Berger de Munich veut consommer moins d'électricité car sa facture augmente ; chacun agit à sa manière (éteindre, débrancher, solaire).
\item \all{will} (wollen), \all{muss} (müssen), \all{kann} (können), \all{können} (können), \all{sollen} (sollen).
\item « La fille retire la fiche de la prise pour économiser de l'énergie. »
\item \all{Er schaltet das Licht ab, wenn er ein Zimmer verlässt.} « Il éteint la lumière quand il quitte une pièce. »
\end{enumerate}
\end{corrige}

\coursec{Compréhension du texte}

Dans la section précédente, nous avons découvert le texte support \all{Energie sparen im Alltag}, l'histoire de la famille Berger qui décide de réduire sa consommation d'électricité. Nous allons maintenant exploiter ce texte en profondeur : comprendre d'abord le sens général, puis les détails, apprendre à répondre correctement aux questions, à repérer les informations vraies ou fausses, et enfin à résumer le texte en quelques phrases. C'est exactement la démarche attendue à l'examen de fin d'année et, plus tard, au baccalauréat.

\courssub{La méthode de compréhension d'un texte}

Lire un texte allemand ne veut pas dire traduire chaque mot. Un bon lecteur procède par étapes : il regarde d'abord le titre, lit le texte une première fois pour l'idée générale, puis une deuxième fois pour les détails, et répond ensuite aux questions en s'appuyant sur le texte.

\begin{popfigure}
\begin{center}
\begin{tikzpicture}[
  node distance=0.55cm,
  etape/.style={rectangle, rounded corners=3pt, draw=popblue, fill=popblueL, thick, align=center, text width=4.6cm, minimum height=0.9cm, font=\small},
  fleche/.style={-{Stealth[length=2.5mm]}, thick, popdark}
]
\node[etape, align=center] (titre){\textbf{1. Titre}\\ \all{Energie sparen im Alltag}\\ Que va raconter le texte ?};
\node[etape, below=of titre, align=center] (globale){\textbf{2. Lecture globale}\\ De quoi parle le texte ? Qui ? Où ?};
\node[etape, below=of globale, align=center] (detail){\textbf{3. Lecture détaillée}\\ Souligner les mots-clés, les chiffres, les actions};
\node[etape, below=of detail, align=center] (questions){\textbf{4. Questions}\\ Repérer la ligne, reformuler, phrase complète};
\node[etape, below=of questions, align=center] (resume){\textbf{5. Résumé}\\ 3 à 5 phrases au présent, sans détails};
\draw[fleche] (titre) -- (globale);
\draw[fleche] (globale) -- (detail);
\draw[fleche] (detail) -- (questions);
\draw[fleche] (questions) -- (resume);
\end{tikzpicture}
\end{center}
\legende{Organigramme de la méthode de compréhension d'un texte}
\end{popfigure}

\begin{methode}[Comment répondre à une question de compréhension]
\begin{enumerate}
\item \textbf{Repérer la ligne du texte} qui contient la réponse : souligne-la au crayon.
\item \textbf{Reformuler avec tes propres mots} : ne recopie pas toute la phrase du texte, prends seulement l'information demandée.
\item \textbf{Répondre par une phrase complète} : sujet + verbe conjugué en 2\up{e} position. Jamais un seul mot.
\item \textbf{Vérifier} : la phrase répond-elle vraiment à la question posée (\all{Wer} ? \all{Was} ? \all{Warum} ? \all{Wann} ?) ?
\end{enumerate}
\end{methode}

\begin{attention}[Erreurs fréquentes des francophones]
À l'écrit, réponds toujours en allemand par une \textbf{phrase complète} : sujet + verbe conjugué en \textbf{2\up{e} position}. Ne recopie jamais la question telle quelle.\\
\emph{Faux :} \all{Warum will die Familie weniger Strom verbrauchen? — Weniger Strom verbrauchen.}\\
\emph{Correct :} \all{Die Familie will weniger Strom verbrauchen, weil die Rechnung immer teurer wird.} « La famille veut consommer moins d'électricité parce que la facture devient de plus en plus chère. »\\
Attention aussi à l'ordre des mots : après \all{weil} (parce que), le verbe conjugué va \textbf{à la fin} de la proposition. En revanche, après \all{denn} (car), l'ordre reste celui d'une phrase principale : \all{..., denn die Rechnung wird immer teurer.}
\end{attention}

\courssub{Compréhension globale}

Relisons le texte support dans sa version intégrale avant de répondre aux questions.

\begin{textebox}[Energie sparen im Alltag]
Die Familie Berger wohnt in einem kleinen Haus in Bamberg. Am Ende des Monats kommt die Stromrechnung — und sie wird immer teurer! „Wir müssen weniger Strom verbrauchen“, sagt der Vater. „Aber wie?“, fragt die Tochter Lena.\\
Der Vater hat eine Idee: „Wir können das Licht ausschalten, wenn wir ein Zimmer verlassen.“ Lena zieht jetzt immer den Stecker heraus, denn der Fernseher und der Computer brauchen auch Strom, wenn sie aus sind. Die Mutter kauft neue Lampen: LED-Lampen brauchen weniger Energie als alte Lampen. Der Sohn Tim lernt in der Schule viel über die Umwelt. Er sagt: „Wir sollen auch die Sonne nutzen! Ein Wasserkraftwerk und Sonnenenergie sind gut für die Natur.“\\
Am Abend sitzt die Familie zusammen und macht einen Plan. Jeder darf eine Idee bringen. Lena will das Handy nicht die ganze Nacht an der Steckdose lassen. Tim will kürzer duschen, denn warmes Wasser braucht viel Energie. Die Mutter will die Waschmaschine nur noch voll machen. Der Vater will im Winter die Heizung nicht so hoch stellen.\\
Nach drei Monaten kommt die neue Rechnung. Die Familie hat zwanzig Prozent weniger Strom verbraucht! Alle sind stolz. „Energie sparen ist nicht schwer“, sagt Lena, „man muss nur daran denken.“
\end{textebox}

\textbf{Glossaire :}
\begin{itemize}
\item \all{die Stromrechnung, -en} : la facture d'électricité
\item \all{verbrauchen (verbraucht, verbrauchte, hat verbraucht)} : consommer
\item \all{ausschalten (schaltet aus, schaltete aus, hat ausgeschaltet)} : éteindre
\item \all{den Stecker herausziehen} : débrancher la prise ; \all{der Stecker, -} : la fiche, la prise (mâle)
\item \all{die Steckdose, -n} : la prise de courant (murale)
\item \all{nutzen} : utiliser, exploiter
\item \all{das Wasserkraftwerk, -e} : la centrale hydroélectrique
\item \all{die Heizung, -en} : le chauffage
\item \all{stolz (auf + Akkusativ)} : fier (de)
\end{itemize}

\textbf{Questions de compréhension globale :}
\begin{enumerate}
\item \all{Wovon handelt der Text?} (De quoi parle le texte ?)
\item \all{Welches Problem hat die Familie Berger?} (Quel est le problème de la famille Berger ?)
\item \all{Was macht die Familie am Ende?} (Que fait la famille à la fin ?)
\end{enumerate}

\begin{exemplebox}[Réponses modèles traduites]
\all{Der Text handelt von einer Familie, die Strom sparen will.} « Le texte parle d'une famille qui veut économiser l'électricité. »\\
\all{Das Problem der Familie ist die teure Stromrechnung.} « Le problème de la famille, c'est la facture d'électricité chère. »\\
\all{Am Ende hat die Familie zwanzig Prozent weniger Strom verbraucht.} « À la fin, la famille a consommé vingt pour cent d'électricité en moins. »
\end{exemplebox}

\begin{attention}[La question globale \all{Wovon handelt der Text?}]
Pour répondre à \all{Wovon handelt der Text?}, utilise la formule toute faite \all{Der Text handelt von + Datif} : \all{Der Text handelt von einer Familie...} Le verbe \all{handeln} se construit avec \all{von} + datif, jamais avec l'accusatif. Ne traduis pas mot à mot « le texte parle de » par \all{der Text spricht über} dans ce contexte d'examen : la formule \all{handeln von} est la plus sûre.
\end{attention}

\courssub{Compréhension détaillée}

Pour la lecture détaillée, on cherche des informations précises : une personne, une action, une raison. Souligne dans le texte la phrase qui contient la réponse avant d'écrire.

\textbf{Questions de compréhension détaillée :}
\begin{enumerate}
\item \all{Wer hat die Idee, das Licht auszuschalten?} (Qui a l'idée d'éteindre la lumière ?)
\item \all{Warum zieht die Tochter den Stecker heraus?} (Pourquoi la fille débranche-t-elle la prise ?)
\item \all{Was kauft die Mutter?} (Qu'achète la mère ?)
\item \all{Was lernt der Sohn in der Schule?} (Qu'apprend le fils à l'école ?)
\end{enumerate}

\begin{exoresolu}[Question détaillée résolue]
Énoncé : réponds par une phrase complète à la question \all{Warum zieht die Tochter den Stecker heraus?}
\tcblower
\textbf{Étape 1 — Repérage :} je souligne dans le texte : \all{Lena zieht jetzt immer den Stecker heraus, denn der Fernseher und der Computer brauchen auch Strom, wenn sie aus sind.}\\
\textbf{Étape 2 — Reformulation :} l'information clé est « le téléviseur et l'ordinateur consomment du courant même éteints ».\\
\textbf{Étape 3 — Phrase complète :} \all{Die Tochter zieht den Stecker heraus, weil der Fernseher und der Computer auch Strom brauchen, wenn sie aus sind.} « La fille débranche la prise parce que le téléviseur et l'ordinateur consomment aussi du courant quand ils sont éteints. »\\
\textbf{Vérification :} sujet (\all{die Tochter}) + verbe en 2\up{e} position (\all{zieht}), particule séparable \all{heraus} à la fin de la principale, verbe de la subordonnée en \all{weil} à la fin (\all{brauchen}). Correct.
\end{exoresolu}

\begin{exemplebox}[Autres réponses modèles]
\all{Der Vater hat die Idee, das Licht auszuschalten.} « C'est le père qui a l'idée d'éteindre la lumière. »\\
\all{Die Mutter kauft neue LED-Lampen.} « La mère achète de nouvelles lampes LED. »\\
\all{Der Sohn lernt in der Schule viel über die Umwelt.} « Le fils apprend beaucoup de choses sur l'environnement à l'école. »\\
\all{Warum will die Familie weniger Strom verbrauchen? — Weil die Rechnung immer teurer wird.} « Pourquoi la famille veut-elle consommer moins d'électricité ? — Parce que la facture devient de plus en plus chère. »
\end{exemplebox}

\begin{attention}[Le mot interrogatif commande la réponse]
\all{Wer} (qui) demande une \textbf{personne} ; \all{was} (quoi) demande une \textbf{chose ou une action} ; \all{warum} (pourquoi) demande une \textbf{raison}, qu'on donne avec \all{weil} + verbe à la fin ; \all{wann} (quand) demande un \textbf{moment} ; \all{wo} (où) demande un \textbf{lieu}. Répondre « où » à une question « pourquoi » est une erreur classique !
\end{attention}

\courssub{Vrai ou faux : justifier par le texte}

L'exercice \all{richtig oder falsch} est très fréquent à l'examen. La règle d'or : \textbf{chaque réponse doit être justifiée par une citation du texte} entre guillemets allemands „…“.

\begin{exercice}[Richtig oder falsch?]
Coche \all{richtig} (vrai) ou \all{falsch} (faux) et justifie par une citation du texte.
\begin{enumerate}
\item \all{Die Familie Berger wohnt in Berlin.}
\item \all{Lena lässt das Handy die ganze Nacht an der Steckdose.}
\item \all{LED-Lampen brauchen weniger Energie als alte Lampen.}
\item \all{Die Familie hat fünfzig Prozent weniger Strom verbraucht.}
\item \all{Energie sparen ist nicht schwer.}
\end{enumerate}
\end{exercice}

\begin{corrige}[Exercice « Richtig oder falsch? »]
\begin{enumerate}
\item \all{Falsch.} Citation : \all{„Die Familie Berger wohnt in einem kleinen Haus in Bamberg.“} (Bamberg, pas Berlin.)
\item \all{Falsch.} Citation : \all{„Lena will das Handy nicht die ganze Nacht an der Steckdose lassen.“}
\item \all{Richtig.} Citation : \all{„LED-Lampen brauchen weniger Energie als alte Lampen.“}
\item \all{Falsch.} Citation : \all{„Die Familie hat zwanzig Prozent weniger Strom verbraucht.“} (vingt, pas cinquante.)
\item \all{Richtig.} Citation : \all{„Energie sparen ist nicht schwer.“}
\end{enumerate}
\end{corrige}

\begin{attention}[Attention aux pièges du vrai/faux]
Les phrases fausses changent souvent \textbf{un seul détail} : un chiffre (\all{zwanzig} → \all{fünfzig}), une ville (\all{Bamberg} → \all{Berlin}), une négation. Lis donc chaque phrase mot par mot et compare-la au texte. Ne réponds jamais de mémoire !
\end{attention}

\courssub{Questions à réponse personnelle}

Après la compréhension du texte, l'examinateur te demande souvent ton avis ou ton expérience. Ici, la réponse n'est pas dans le texte : tu dois produire tes propres phrases, avec le vocabulaire du thème.

\textbf{Question modèle :} \all{Was machst du, um Energie zu sparen?} (Que fais-tu pour économiser l'énergie ?)

\begin{exemplebox}[Réponses personnelles possibles]
\all{Ich schalte das Licht aus, wenn ich das Zimmer verlasse.} « J'éteins la lumière quand je quitte la pièce. »\\
\all{Ich lasse das Handy nicht die ganze Nacht an der Steckdose.} « Je ne laisse pas le portable toute la nuit sur la prise. »\\
\all{Ich dusche kürzer, um warmes Wasser zu sparen.} « Je me douche plus vite pour économiser l'eau chaude. »\\
\all{In Yaoundé nutzen wir die Sonne, denn sie scheint oft.} « À Yaoundé, nous utilisons le soleil, car il brille souvent. »
\end{exemplebox}

\begin{center}
\begin{tabularx}{\linewidth}{|X|X|X|}
\hline
\rowcolor{popblueL} Français & Deutsch & Remarque \\
\hline
Que fais-tu pour économiser l'énergie ? & \all{Was machst du, um Energie zu sparen?} & \all{um ... zu} = « pour » (finalité) \\
\hline
J'éteins la lumière. & \all{Ich schalte das Licht aus.} & verbe à particule : \all{aus-} va à la fin \\
\hline
Je débranche la prise. & \all{Ich ziehe den Stecker heraus.} & particule \all{heraus} à la fin \\
\hline
Parce que la facture est chère. & \all{Weil die Rechnung teuer ist.} & verbe \all{ist} à la fin après \all{weil} \\
\hline
Nous devons économiser l'énergie. & \all{Wir müssen Energie sparen.} & modal \all{müssen} + infinitif à la fin \\
\hline
\end{tabularx}
\end{center}

\courssub{La technique du résumé}

Résumer, c'est dire l'essentiel en peu de mots. Au niveau Première, on attend un résumé de \textbf{3 à 5 phrases}, au \textbf{présent}, sans détails ni exemples.

\begin{methode}[Comment rédiger un résumé]
\begin{enumerate}
\item Relis le texte et relève \textbf{une idée principale par paragraphe} (ici : 5 idées).
\item Rédige \textbf{une phrase simple par idée}, au présent, sujet + verbe en 2\up{e} position.
\item Supprime les détails : chiffres précis, dialogues, exemples, noms propres secondaires.
\item Relie les phrases avec des connecteurs simples : \all{dann} (puis), \all{auch} (aussi), \all{am Ende} (à la fin).
\item Compte : 3 à 5 phrases maximum.
\end{enumerate}
\end{methode}

\textbf{Les 5 idées principales du texte :}
\begin{enumerate}
\item La famille Berger a un problème : la facture d'électricité est de plus en plus chère.
\item Chaque membre de la famille trouve une idée pour économiser le courant (lumière, prises, lampes LED, chauffage).
\item Le fils parle des énergies naturelles (soleil, eau).
\item La famille fait un plan commun.
\item Après trois mois, la famille a économisé vingt pour cent d'électricité.
\end{enumerate}

\begin{exemplebox}[Résumé modèle (5 phrases)]
\all{Die Familie Berger will weniger Strom verbrauchen, weil die Rechnung teuer wird. Jeder in der Familie hat eine Idee: das Licht ausschalten, den Stecker herausziehen, neue Lampen kaufen. Der Sohn lernt in der Schule viel über Sonnenenergie und Wasserkraft. Die Familie macht zusammen einen Plan. Am Ende spart die Familie zwanzig Prozent Strom.}\\
« La famille Berger veut consommer moins d'électricité parce que la facture devient chère. Chacun dans la famille a une idée : éteindre la lumière, débrancher les prises, acheter de nouvelles lampes. Le fils apprend beaucoup de choses à l'école sur l'énergie solaire et l'hydroélectricité. La famille fait un plan ensemble. À la fin, la famille économise vingt pour cent d'électricité. »
\end{exemplebox}

\begin{attention}[Ce qu'un résumé n'est pas]
Un résumé n'est \textbf{pas} une traduction du texte, ni une copie des phrases du texte, ni ton opinion personnelle. On n'écrit pas \all{Ich finde...} dans un résumé. Pas de guillemets, pas de dialogue, pas de détails comme \all{in Bamberg} ou \all{am Abend}.
\end{attention}

\courssub{Entraînement oral : questions-réponses en binôme}

\begin{experience}[Interview en binôme : la famille Berger]
Travaillez à deux. L'élève A est journaliste, l'élève B est un membre de la famille Berger (père, mère, Lena ou Tim). Le journaliste pose les questions, l'autre répond par des phrases complètes, puis on échange les rôles.
\begin{enumerate}
\item \all{Warum wollen Sie weniger Strom verbrauchen?}
\item \all{Was machen Sie, um Energie zu sparen?}
\item \all{Was macht Ihre Tochter / Ihr Sohn?}
\item \all{Was haben Sie am Ende gespart?}
\item \all{Was machst du zu Hause, um Energie zu sparen?} (question personnelle)
\end{enumerate}
Critères de réussite : répondre par une phrase complète, verbe en 2\up{e} position, prononciation claire (\all{Strom} « chtrom », \all{sparen} « chpa-ren », \all{Steckdose} « chtek-do-ze »).
\end{experience}

\begin{aretenir}[L'essentiel de la section]
\begin{itemize}
\item Comprendre un texte se fait en étapes : titre → lecture globale → lecture détaillée → questions → résumé.
\item On répond à une question par une \textbf{phrase complète} : sujet + verbe conjugué en 2\up{e} position ; après \all{weil}, le verbe va à la fin ; après \all{denn}, l'ordre ne change pas.
\item Au vrai/faux, on \textbf{justifie toujours par une citation} du texte entre guillemets „…“.
\item Un résumé = 3 à 5 phrases au présent, une idée par phrase, sans détails ni opinion.
\end{itemize}
\end{aretenir}

\coursec{Lexique thématique : die Energie}

Après avoir travaillé la compréhension globale et détaillée du texte \emph{Energie sparen im Alltag}, nous allons maintenant fixer et structurer le vocabulaire essentiel du thème. Ce lexique vous servira de base pour les exercices de grammaire (verbes de modalité, expression de la finalité) et pour vos productions écrites et orales. Nous l'organisons en trois champs sémantiques logiques : les sources d'énergie, la consommation domestique et les actions.

\courssub{Champ 1 : Les sources d'énergie (die Energiequellen)}

Observons d'abord comment le texte et le langage courant nomment les différentes origines de l'énergie. Notez systématiquement l'article, le pluriel et, pour les mots composés, la décomposition.

\begin{center}
\begin{tabularx}{\linewidth}{|X|X|X|X|}
\hline
\rowcolor{popblueL}
\textbf{Allemand (Singulier / Pluriel)} & \textbf{Décomposition} & \textbf{Français} & \textbf{Remarque} \\
\hline
die Energie, die Energien & — & l'énergie & Nom féminin, pluriel en \textbf{-n}. \\
\hline
der Strom & — & l'électricité (le courant) & Masculin, pas de pluriel au sens général. \\
\hline
die Sonne, die Sonnen & — & le soleil & Féminin. Composés : \emph{die Sonnenenergie, das Sonnenlicht}. \\
\hline
das Wasser & — & l'eau & Neutre, pas de pluriel usuel. Composés : \emph{die Wasserkraft, der Wasserstoff}. \\
\hline
der Wind, die Winde & — & le vent & Masculin, pluriel \textbf{-e}. \\
\hline
das Wasserkraftwerk, die Wasserkraftwerke & Wasser + Kraft + Werk & la centrale hydroélectrique & Neutre (dernier élément : \emph{das Werk}). \\
\hline
das Solarpanel, die Solarpanels & Solar + Panel & le panneau solaire & Neutre, pluriel \textbf{-s} (mot d'origine anglaise). \\
\hline
die Solaranlage, die Solaranlagen & Solar + Anlage & l'installation solaire & Féminin (dernier élément : \emph{die Anlage}). \\
\hline
das Kohlekraftwerk, die Kohlekraftwerke & Kohle + Kraft + Werk & la centrale à charbon & Neutre. \emph{die Kohle} = le charbon. \\
\hline
die erneuerbare Energie, die erneuerbaren Energien & erneuerbar + Energie & l'énergie renouvelable & Adjectif \emph{erneuerbar} (renouvelable) décliné. \\
\hline
\end{tabularx}
\end{center}

\begin{definition}[Énergie renouvelable — \emph{die erneuerbare Energie}]
On appelle \cle{énergies renouvelables} (\all{die erneuerbaren Energien}) les sources d'énergie qui se régénèrent naturellement à l'échelle humaine : le rayonnement solaire (\all{die Sonnenenergie}), le vent (\all{die Windenergie}), la force de l'eau (\all{die Wasserkraft}), la biomasse (\all{die Biomasse}) et la géothermie (\all{die Geothermie}). Contrairement aux énergies fossiles (\all{die fossilen Energieträger} : le charbon \all{die Kohle}, le pétrole \all{das Erdöl}, le gaz \all{das Gas}), elles n'épuisent pas les réserves de la planète.
\end{definition}

\begin{savaistu}[La mécanique des mots composés : \emph{das Wasserkraftwerk}]
L'allemand adore les mots composés (\cle{die Komposita}). La règle d'or pour le genre : \textbf{le dernier élément détermine l'article et le pluriel}.
\begin{itemize}
    \item \all{das Werk} (l'usine, l'ouvrage) $\rightarrow$ neutre, pluriel \emph{-e} : \emph{die Werke}.
    \item \all{das Wasserkraftwerk} = \emph{Wasser} (eau) + \emph{Kraft} (force, puissance) + \emph{Werk} (usine) = littéralement « usine de force hydraulique ».
    \item \all{das Kohlekraftwerk} = \emph{Kohle} + \emph{Kraft} + \emph{Werk}.
    \item \all{die Solaranlage} = \emph{Solar} + \emph{Anlage} (installation) $\rightarrow$ féminin (\emph{die Anlage}), pluriel \emph{-n}.
\end{itemize}
À l'examen, on vous demandera souvent de décomposer un mot composé ou d'en trouver le genre. Regardez toujours la fin du mot !
\end{savaistu}

\begin{exemplebox}[Familles de mots : \emph{Energie} et \emph{Strom}]
\begin{tabular}{lll}
\all{die Energie} & $\rightarrow$ & \all{energisch} (adj., énergique) ; \all{der Energiesparmodus} (mode économie d'énergie). \\
\all{der Strom} & $\rightarrow$ & \all{stromlos} (adj., sans courant) ; \all{der Stromverbrauch} (consommation d'électricité) ; \all{der Strompreis} (prix de l'électricité).
\end{tabular}
\end{exemplebox}

\courssub{Champ 2 : La consommation à la maison (der Verbrauch im Haushalt)}

Le texte mentionne plusieurs appareils et éléments concrets de la vie quotidienne. Voici le lexique indispensable pour décrire votre propre consommation.

\begin{center}
\begin{tabularx}{\linewidth}{|X|X|X|}
\hline
\rowcolor{poporangeL}
\textbf{Allemand (Singulier / Pluriel)} & \textbf{Français} & \textbf{Contexte / Exemple} \\
\hline
die Steckdose, die Steckdosen & la prise (murale) & \all{Die Steckdose ist hinter dem Schrank.} \\
\hline
der Stecker, die Stecker & la fiche (au bout du câble) & \all{Zieh den Stecker heraus!} \\
\hline
das Licht, die Lichter & la lumière, l'éclairage & \all{Das Licht brennt noch.} \\
\hline
der Kühlschrank, die Kühlschränke & le réfrigérateur & Umlaut au pluriel : \emph{ä}. \all{Der Kühlschrank läuft Tag und Nacht.} \\
\hline
die Heizung, die Heizungen & le chauffage (installation) & \all{Die Heizung ist zu warm eingestellt.} \\
\hline
das Gerät, die Geräte & l'appareil (électroménager) & Umlaut au pluriel. \all{Alte Geräte verbrauchen viel Strom.} \\
\hline
die Rechnung, die Rechnungen & la facture & \all{Die Stromrechnung ist hoch.} \\
\hline
der Verbrauch (kein Plural) & la consommation & \all{Der Verbrauch ist gesunken.} \\
\hline
\end{tabularx}
\end{center}

\begin{popfigure}
\begin{tikzpicture}[
    scale=0.85, transform shape,
    haus/.style={draw=popdark, thick, fill=popcream, rounded corners=3pt},
    labelstyle/.style={font=\small\bfseries, text=popdark, fill=popcream, inner sep=2pt, rounded corners=1pt},
    arrowstyle/.style={-Stealth, poporange, thick}
]
% Maison
\draw[haus] (0,0) rectangle (8,5);
\draw[haus, fill=poporange!20] (0,5) -- (4,8) -- (8,5) -- cycle;

% Solarpanel sur le toit
\draw[fill=popblue!30, draw=popblue, thick] (1.5,5.5) rectangle (6.5,6.2);
\node[labelstyle] at (4,5.85) {Solarpanel};

% Fenêtres / Porte
\draw[fill=popblueL, draw=popdark] (1,1) rectangle (2.5,2.5);
\draw[fill=popblueL, draw=popdark] (5.5,1) rectangle (7,2.5);
\draw[fill=popgold!50, draw=popdark] (3.2,0) rectangle (4.8,2.2);

% Étiquettes intérieures
\node[labelstyle, anchor=west] at (0.2, 4.2) {Heizung};
\node[labelstyle, anchor=west] at (0.2, 3.0) {Kühlschrank};
\node[labelstyle, anchor=east] at (7.8, 4.2) {Steckdose};
\node[labelstyle, anchor=east] at (7.8, 2.5) {Licht};

% Flèches de liaison
\draw[arrowstyle] (4, 6.2) -- (4, 5.0) node[midway, right, labelstyle] {Strom};
\draw[arrowstyle, dashed] (4, 5.0) -- (1.5, 4.2);
\draw[arrowstyle, dashed] (4, 5.0) -- (6.5, 4.2);
\draw[arrowstyle, dashed] (4, 5.0) -- (1.5, 3.0);
\draw[arrowstyle, dashed] (4, 5.0) -- (6.5, 2.5);
\end{tikzpicture}
\legende{Schéma d'une maison équipée : le vocabulaire clé du Champ 2 en images.}
\end{popfigure}

\begin{attention}[Piège n°1 : \emph{der Stecker} ou \emph{die Steckdose} ?]
C'est l'erreur n°1 des francophones ! En français, « la prise » désigne les deux objets. En allemand, la distinction est stricte :
\begin{itemize}
    \item \cle{die Steckdose} (féminin) = la prise \textbf{fixe au mur} (celle qui reçoit).
    \item \cle{der Stecker} (masculin) = la fiche \textbf{mobile} au bout du câble de l'appareil (celle qu'on enfiche).
\end{itemize}
On tire la fiche hors de la prise : \all{den Stecker aus der Steckdose ziehen}. \\
Ne dites jamais « \emph{den Stecker in den Stecker stecken} » !
\end{attention}

\courssub{Champ 3 : Les actions (die Handlungen) — les verbes clés}

La plupart de ces verbes sont à \cle{particule séparable} (\all{trennbare Verben}). Au présent de l'indicatif, la particule se détache et va en \textbf{fin de phrase}. Nous les classons pour faciliter la mémorisation.

\begin{center}
\begin{tabularx}{\linewidth}{|X|X|X|X|}
\hline
\rowcolor{popgreenL}
\textbf{Verbe (Infinitif)} & \textbf{Particule + Radical} & \textbf{Présent (ich / er)} & \textbf{Français} \\
\hline
\all{abschalten} & ab + schalten & ich schalte ab / er schaltet ab & éteindre, couper (appareil) \\
\hline
\all{einschalten} & ein + schalten & ich schalte ein / er schaltet ein & allumer (appareil, lumière) \\
\hline
\all{ausschalten} & aus + schalten & ich schalte aus / er schaltet aus & éteindre (appareil, lumière) \\
\hline
\all{herausziehen} & heraus + ziehen & ich ziehe heraus / er zieht heraus & tirer dehors, débrancher (la fiche) \\
\hline
\all{sparen} & (inséparable) & ich spare / er spart & économiser \\
\hline
\all{verbrauchen} & (inséparable) & ich verbrauche / er verbraucht & consommer \\
\hline
\all{nutzen} & (inséparable) & ich nutze / er nutzt & utiliser, exploiter \\
\hline
\all{produzieren} & (inséparable) & ich produziere / er produziert & produire \\
\hline
\all{schützen} & (inséparable) & ich schütze / er schützt & protéger \\
\hline
\all{bezahlen} & (inséparable) & ich bezahle / er bezahlt & payer \\
\hline
\end{tabularx}
\end{center}

\begin{propriete}[Verbes à particule séparable : règle de position]
Pour les verbes \textbf{abschalten, einschalten, ausschalten, herausziehen} (et tous les verbes à particule accentuée : \emph{ab-, an-, auf-, aus-, ein-, heraus-, mit-, vor-, weg-, zu-, zurück-}) :
\begin{enumerate}
    \item \textbf{Phrase principale (présent, prétérit, impératif) :} le radical conjugué est en 2\textsuperscript{e} position (1\textsuperscript{re} à l'impératif) et la particule va à la \textbf{fin de la phrase}.
    \item \textbf{Phrase subordonnée (avec \emph{dass, weil, wenn...}) :} la particule \textbf{reste collée} au radical et le verbe conjugué complet va à la fin.
    \item \textbf{Infinitif avec \emph{zu} :} \emph{zu} s'insère entre la particule et le radical (\emph{abzuschalten}, \emph{herauszuziehen}).
    \item \textbf{Après un verbe de modalité :} l'infinitif complet (particule collée) va à la fin : \all{Wir müssen das Licht ausschalten.}
\end{enumerate}
\end{propriete}

\begin{exemplebox}[Exemples avec verbes séparables]
\begin{enumerate}
    \item \all{Wir schalten das Licht aus, wenn wir den Raum verlassen.} \\
    (Nous éteignons la lumière quand nous quittons la pièce.) — \emph{Phrase principale : aus à la fin.}
    \item \all{Ich muss den Stecker herausziehen, bevor ich das Gerät repariere.} \\
    (Je dois débrancher la fiche avant de réparer l'appareil.) — \emph{Après le modal müssen : infinitif complet herausziehen à la fin.}
    \item \all{Weil alte Geräte viel Strom verbrauchen, schalten wir sie ab.} \\
    (Parce que les vieux appareils consomment beaucoup d'électricité, nous les éteignons.) — \emph{Subordonnée : verbe à la fin ; principale : ab à la fin.}
    \item \all{Schalten Sie bitte den Fernseher aus!} \\
    (Éteignez la télévision, s'il vous plaît !) — \emph{Impératif : aus à la fin.}
\end{enumerate}
\end{exemplebox}

\begin{attention}[Piège n°2 : \emph{ausschalten}, \emph{abschalten} ou \emph{ausmachen} ?]
\begin{itemize}
    \item \all{ausschalten} / \all{abschalten} : registre standard et neutre, pour les appareils, la lumière, les moteurs.
    \item \all{ausmachen} (séparable) : familier, très oral, surtout pour la lumière (\all{das Licht ausmachen}). À éviter dans une copie d'examen.
    \item Pour « allumer », les correspondants sont : \all{einschalten}, \all{anschalten} (standard) et \all{anmachen} (familier).
\end{itemize}
Pour l'examen : privilégiez \textbf{\emph{einschalten / ausschalten}} ou \textbf{\emph{abschalten}}.
\end{attention}

\courssub{Familles de mots (die Wortfamilien) et expressions figées}

Travailler par familles permet de multiplier votre vocabulaire sans apprendre des listes isolées.

\begin{center}
\begin{tabularx}{\linewidth}{|X|X|X|}
\hline
\rowcolor{poppurpleL}
\textbf{Mot de base} & \textbf{Dérivés (Nom / Adjectif)} & \textbf{Sens / Traduction} \\
\hline
\all{sparen} (v.) & \all{die Ersparnis, -se} (f.), \all{sparsam} (adj.), \all{der Sparer, -} (m.) & économie ; économe ; celui qui économise \\
\hline
\all{verbrauchen} (v.) & \all{der Verbrauch} (m.), \all{der Verbraucher, -} (m.), \all{verbrauchsarm} (adj.) & consommation ; consommateur ; peu consommateur \\
\hline
\all{die Energie} (f.) & \all{energisch} (adj.), \all{die Energieeffizienz} (f.), \all{energiesparend} (adj.) & énergique ; efficacité énergétique ; économe en énergie \\
\hline
\all{produzieren} (v.) & \all{die Produktion} (f.), \all{der Produzent, -en} (m.), \all{produktiv} (adj.) & production ; producteur ; productif \\
\hline
\all{schützen} (v.) & \all{der Schutz} (m.), \all{der Umweltschutz} (m.), \all{schützenswert} (adj.) & protection ; protection de l'environnement ; digne d'être protégé \\
\hline
\end{tabularx}
\end{center}

\begin{aretenir}[Expressions utiles pour l'oral et l'écrit (Redemittel)]
Apprenez ces blocs par cœur : ils sont la clé de la fluidité.
\begin{itemize}
    \item \all{Strom sparen} — économiser de l'électricité (collocation forte).
    \item \all{Energie verbrauchen} — consommer de l'énergie.
    \item \all{das Licht ausschalten / abschalten} — éteindre la lumière.
    \item \all{den Stecker (aus der Steckdose) ziehen / herausziehen} — débrancher la fiche.
    \item \all{die Umwelt schützen} — protéger l'environnement.
    \item \all{erneuerbare Energien nutzen} — utiliser les énergies renouvelables.
    \item \all{die Stromrechnung bezahlen} — payer la facture d'électricité.
    \item \all{energiesparende Geräte} — des appareils économes en énergie.
\end{itemize}
\end{aretenir}

\courssub{Entraînement sur texte : \emph{Mein Beitrag zum Klimaschutz}}

Le texte suivant réemploie tout le vocabulaire étudié ci-dessus. Lisez-le, soulignez les mots connus, déduisez les mots inconnus du contexte, puis répondez aux questions.

\begin{textebox}[Texte original : Mein Beitrag zum Klimaschutz]
„Klimaschutz beginnt zu Hause“, sagt Amina, eine Schülerin aus Yaoundé. „In Kamerun haben wir viel Sonne, aber der Strom ist oft teuer, und manchmal fällt er ganz aus. Deshalb haben meine Eltern vor zwei Jahren eine \textbf{Solaranlage} auf dem Dach installiert. Die \textbf{Solarpanels} produzieren jetzt einen großen Teil unseres \textbf{Stroms}. Das \textbf{Wasserkraftwerk} in der Nähe liefert den Rest.

Früher haben wir nicht auf den \textbf{Verbrauch} geachtet. Der \textbf{Kühlschrank} war alt, und die \textbf{Heizung} lief im Winter auf höchster Stufe. Heute ist alles anders. Wir \textbf{schalten} das \textbf{Licht} konsequent \textbf{aus}, wenn wir einen Raum verlassen. Wir \textbf{ziehen} die \textbf{Stecker} vom Fernseher und vom Computer \textbf{heraus}, weil die Geräte im \textbf{Standby-Modus} auch \textbf{Strom verbrauchen}. Meine Mutter \textbf{kauft} nur noch \textbf{energiesparende Geräte} der Klasse A+++.

Am Ende des Monats ist die \textbf{Stromrechnung} viel niedriger. Die \textbf{Ersparnis} geben wir für Schulbücher aus. So \textbf{schützen} wir nicht nur die \textbf{Umwelt}, sondern auch unseren Geldbeutel. Jeder kann etwas tun: \textbf{Strom sparen} ist der erste Schritt.“
\end{textebox}

\begin{definition}[Glossaire du texte (mots difficiles)]
\begin{itemize}
    \item \all{das Dach, die Dächer} — le toit (Umlaut au pluriel).
    \item \all{installieren} — installer (verbe inséparable, Partizip II : \emph{installiert}).
    \item \all{liefern} — fournir, livrer.
    \item \all{früher} — autrefois, avant.
    \item \all{auf etwas achten} — faire attention à quelque chose (\emph{achten auf} + accusatif).
    \item \all{konsequent} — systématiquement, de manière conséquente.
    \item \all{der Standby-Modus, die Standby-Modi} — le mode veille.
    \item \all{die Klasse A+++} — la classe énergétique A+++ (la meilleure).
    \item \all{niedrig} — bas, faible (prix, niveau).
    \item \all{der Geldbeutel, -} — le porte-monnaie (ici : le budget).
    \item \all{der Schritt, -e} — le pas, l'étape.
\end{itemize}
\end{definition}

\begin{exoresolu}[Exercice résolu : compléter et traduire]
\textbf{Consigne :} Complétez les phrases avec la forme correcte du verbe entre parenthèses (Présent) et le vocabulaire adéquat. Traduisez ensuite en français.

1. Wir \_\_\_\_\_\_\_\_\_ (ausschalten) das Licht, wenn wir das Haus verlassen. \\
2. Der alte Kühlschrank \_\_\_\_\_\_\_\_\_ (verbrauchen) zu viel \_\_\_\_\_\_\_\_\_ (Strom). \\
3. Bitte \_\_\_\_\_\_\_\_\_ (herausziehen) Sie den \_\_\_\_\_\_\_\_\_ (Stecker) aus der \_\_\_\_\_\_\_\_\_ (Steckdose)! \\
4. Die neue Solaranlage \_\_\_\_\_\_\_\_\_ (produzieren) saubere \_\_\_\_\_\_\_\_\_ (Energie). \\
5. Wir wollen die Umwelt \_\_\_\_\_\_\_\_\_ (schützen) und Geld \_\_\_\_\_\_\_\_\_ (sparen).

\tcblower
\textbf{Solution détaillée :}
\begin{enumerate}
    \item \all{Wir schalten das Licht aus, wenn wir das Haus verlassen.} \\
    Verbe séparable : radical \emph{schalten} en 2\textsuperscript{e} position, particule \emph{aus} en fin de phrase principale. \\
    \emph{Traduction :} Nous éteignons la lumière quand nous quittons la maison.
    \item \all{Der alte Kühlschrank verbraucht zu viel Strom.} \\
    Verbe inséparable, 3\textsuperscript{e} personne du singulier : \emph{verbraucht}. \emph{Strom} est masculin, sans article après \emph{viel}. \\
    \emph{Traduction :} Le vieux réfrigérateur consomme trop d'électricité.
    \item \all{Bitte ziehen Sie den Stecker aus der Steckdose heraus!} \\
    Impératif formel (\emph{Sie}) : verbe en 1\textsuperscript{re} position, particule \emph{heraus} à la fin. \emph{den Stecker} = accusatif masculin ; \emph{aus der Steckdose} = datif féminin après la préposition \emph{aus}. \\
    \emph{Traduction :} Veuillez débrancher la fiche de la prise !
    \item \all{Die neue Solaranlage produziert saubere Energie.} \\
    Verbe inséparable, 3\textsuperscript{e} personne du singulier : \emph{produziert}. \emph{Energie} = accusatif féminin singulier, adjectif \emph{saubere} décliné. \\
    \emph{Traduction :} La nouvelle installation solaire produit de l'énergie propre.
    \item \all{Wir wollen die Umwelt schützen und Geld sparen.} \\
    Après le verbe de modalité \emph{wollen}, les deux infinitifs \emph{schützen} et \emph{sparen} vont à la fin. \\
    \emph{Traduction :} Nous voulons protéger l'environnement et économiser de l'argent.
\end{enumerate}
\end{exoresolu}

\begin{exercice}[Entraînement personnel : thème (français $\rightarrow$ allemand)]
Traduisez les phrases suivantes en allemand en utilisant \textbf{impérativement} le vocabulaire de la leçon (Champs 1, 2, 3) et en respectant la place des verbes séparables.

\begin{enumerate}
    \item Le panneau solaire sur le toit produit de l'électricité.
    \item Éteins la lumière ! (tu / \emph{du})
    \item Nous devons débrancher la fiche de la prise.
    \item L'ancienne centrale à charbon consomme beaucoup de charbon.
    \item Ma mère paie la facture d'électricité chaque mois.
    \item Les appareils économes en énergie protègent l'environnement.
    \item Pourquoi le chauffage consomme-t-il tant d'énergie ?
    \item Utilisons les énergies renouvelables ! (impératif à la 1\textsuperscript{re} personne du pluriel)
\end{enumerate}
\end{exercice}

\begin{corrige}[Exercice : thème]
\begin{enumerate}
    \item \all{Das Solarpanel auf dem Dach produziert Strom.} \\
    Sujet neutre singulier, verbe en 2\textsuperscript{e} position ; \emph{Strom} sans article (sens général).
    \item \all{Schalte das Licht aus!} \\
    Impératif \emph{du} : radical \emph{schalte} en 1\textsuperscript{re} position, particule \emph{aus} à la fin. (Forme familière possible : \all{Mach das Licht aus!})
    \item \all{Wir müssen den Stecker aus der Steckdose ziehen.} \\
    Verbe de modalité \emph{müssen} en 2\textsuperscript{e} position, infinitif \emph{ziehen} à la fin ; \emph{den Stecker} = accusatif ; \emph{aus der Steckdose} = datif féminin.
    \item \all{Das alte Kohlekraftwerk verbraucht viel Kohle.} \\
    Sujet neutre (\emph{das Kraftwerk}), adjectif \emph{alte} décliné, verbe inséparable.
    \item \all{Meine Mutter bezahlt die Stromrechnung jeden Monat.} \\
    Possessif \emph{meine}, complément \emph{die Stromrechnung} à l'accusatif, complément de temps \emph{jeden Monat} (accusatif masculin).
    \item \all{Energiesparende Geräte schützen die Umwelt.} \\
    Adjectif \emph{energiesparende} décliné au nominatif pluriel sans article (terminaison \emph{-e}).
    \item \all{Warum verbraucht die Heizung so viel Energie?} \\
    Question en \emph{Warum} : verbe en 2\textsuperscript{e} position, sujet après le verbe.
    \item \all{Nutzen wir erneuerbare Energien!} \\
    Impératif à la 1\textsuperscript{re} personne du pluriel : verbe conjugué en 1\textsuperscript{re} position suivi de \emph{wir} (comme en français « Utilisons... ! »).
\end{enumerate}
\end{corrige}

\begin{methode}[Comment réviser ce lexique pour l'examen ?]
\begin{enumerate}
    \item \textbf{Carte mentale :} dessinez une maison (comme le schéma ci-dessus) et placez les mots \emph{Steckdose, Stecker, Kühlschrank, Heizung, Licht, Solarpanel} aux bons endroits, avec leur article.
    \item \textbf{Tableau des verbes séparables :} faites quatre colonnes — Infinitif / Particule / Exemple en phrase principale / Exemple en subordonnée — et remplissez-le pour \emph{abschalten, einschalten, ausschalten, herausziehen}.
    \item \textbf{Familles de mots :} pour chaque verbe du Champ 3, cherchez le nom et l'adjectif correspondants (ex. \emph{produzieren} $\rightarrow$ \emph{die Produktion}, \emph{produktiv}, \emph{der Produzent}).
    \item \textbf{Collocations :} apprenez \emph{Strom sparen}, \emph{die Umwelt schützen}, \emph{die Rechnung bezahlen}, \emph{den Stecker ziehen} comme des blocs indissociables, et réemployez-les dans vos productions écrites.
\end{enumerate}
\end{methode}

\coursec{Grammaire 1 : les verbes de modalité}

Dans la section précédente, nous avons découvert le vocabulaire du thème \all{die Energie} : \all{der Strom}, \all{die Sonne}, \all{das Wasserkraftwerk}, \all{sparen}, \all{die Steckdose}. Ces mots nous ont permis de comprendre le texte \emph{Energie sparen im Alltag}. Mais en lisant ce texte, tu as peut-être remarqué des petits verbes qui reviennent sans cesse : \all{muss}, \all{können}, \all{sollen}, \all{will}. Ce sont les \cle{verbes de modalité} (en allemand : \all{die Modalverben}). Ils sont absolument essentiels : sans eux, impossible de dire ce que l'on \emph{peut}, \emph{doit}, \emph{veut} ou \emph{est autorisé} à faire. Cette section va te les faire maîtriser pas à pas.

\courssub{Observation : que remarque-t-on dans le texte ?}

Reprenons trois phrases tirées de notre texte support :

\begin{exemplebox}[Les modaux dans le texte]
\all{Sie muss die Rechnung bezahlen.} « Elle doit payer la facture. »\par
\all{Wir können alle etwas tun.} « Nous pouvons tous faire quelque chose. »\par
\all{Wir sollen die Erde schützen.} « Nous devrions protéger la Terre. »
\end{exemplebox}

Observons attentivement ces trois phrases :

\begin{itemize}
\item Dans chaque phrase, il y a \textbf{deux verbes} : un verbe conjugué (\all{muss}, \all{können}, \all{sollen}) et un verbe à l'infinitif (\all{bezahlen}, \all{tun}, \all{schützen}).
\item Le verbe conjugué occupe la \textbf{deuxième position} de la phrase, exactement comme n'importe quel verbe conjugué en allemand.
\item L'infinitif est renvoyé \textbf{à la toute fin} de la phrase.
\item Le verbe conjugué n'exprime pas une action : il exprime une \emph{attitude} face à l'action (obligation, possibilité, devoir moral). C'est l'infinitif final qui porte le sens de l'action.
\end{itemize}

\begin{definition}[Le verbe de modalité]
Un \cle{verbe de modalité} (\all{das Modalverb, die Modalverben}) est un verbe qui accompagne un infinitif et qui précise le \emph{rapport} du sujet à l'action : la capacité, la possibilité, l'obligation, la permission, la volonté ou le désir. En français, ce rôle est joué par les semi-auxiliaires « pouvoir, devoir, vouloir ».
\end{definition}

\courssub{Les six verbes de modalité et leur sens}

L'allemand compte \textbf{six} verbes de modalité. Il faut connaître le sens de chacun par cœur, car ils ne se recouvrent pas exactement avec le français.

\begin{center}
\begin{tabularx}{\linewidth}{|l|X|X|}
\hline
\rowcolor{popblueL} \textbf{Modalverb} & \textbf{Sens principal} & \textbf{Exemple} \\
\hline
\all{können} & capacité, possibilité (« pouvoir, savoir ») & \all{Wir können die Sonne nutzen.} « Nous pouvons utiliser le soleil. » \\
\hline
\all{müssen} & obligation forte, nécessité (« devoir ») & \all{Du musst das Licht ausschalten.} « Tu dois éteindre la lumière. » \\
\hline
\all{dürfen} & permission, autorisation (« avoir le droit de ») & \all{Man darf hier nicht parken.} « Il est interdit de se garer ici. » \\
\hline
\all{sollen} & devoir moral, conseil, consigne reçue (« devoir, être censé ») & \all{Wir sollen die Umwelt schützen.} « Nous devrions protéger l'environnement. » \\
\hline
\all{wollen} & volonté, intention ferme (« vouloir ») & \all{Er will ein Solarpanel kaufen.} « Il veut acheter un panneau solaire. » \\
\hline
\all{mögen} (et \all{möchte}) & goût, désir ; \all{möchte} = souhait poli (« aimer, vouloir ») & \all{Ich möchte Energie sparen.} « Je voudrais économiser de l'énergie. » \\
\hline
\end{tabularx}
\end{center}

\begin{attention}[\all{müssen} ou \all{sollen} ?]
La distinction \all{müssen} / \all{sollen} est un piège classique. \all{müssen} exprime une obligation qui vient de la situation ou de soi-même : \all{Ich muss die Rechnung bezahlen.} « Je dois payer la facture » (sinon, coupure de courant !). \all{sollen} exprime un devoir moral, un conseil ou une consigne venant d'une autre personne : \all{Der Lehrer sagt, wir sollen Strom sparen.} « Le professeur dit que nous devons économiser le courant. » En français, les deux se traduisent souvent par « devoir », mais l'allemand distingue nettement.
\end{attention}

\begin{attention}[La négation : un faux ami dangereux]
Avec la négation, le sens change de façon inattendue pour un francophone :
\begin{itemize}
\item \all{Du musst das nicht tun.} = « Tu n'es pas obligé de faire cela » (ce n'est pas nécessaire).
\item \all{Du darfst das nicht tun.} = « Tu ne dois pas faire cela, c'est interdit ».
\end{itemize}
Pour dire « il ne faut pas », l'allemand utilise donc \all{nicht dürfen} (ou \all{nicht sollen}), et non \all{nicht müssen}. Compare : \all{Man darf hier nicht parken.} « Il est interdit de se garer ici » et \all{Man muss hier nicht parken.} « On n'est pas obligé de se garer ici ».
\end{attention}

\courssub{Conjugaison des six modaux au présent}

Voici le tableau complet à apprendre par cœur. Observe bien les deux grandes particularités : la \textbf{voyelle change au singulier} et la \textbf{troisième personne du singulier ne prend pas de -t... ni de -s} : elle est identique à la première personne (\all{ich kann} / \all{er kann}).

\begin{center}
\begin{tabular}{|l|l|l|l|l|l|l|}
\hline
\rowcolor{popgreenL} \textbf{Personne} & \textbf{können} & \textbf{müssen} & \textbf{dürfen} & \textbf{sollen} & \textbf{wollen} & \textbf{mögen} \\
\hline
ich & kann & muss & darf & soll & will & mag \\
\hline
du & kannst & musst & darfst & sollst & willst & magst \\
\hline
er/sie/es & kann & muss & darf & soll & will & mag \\
\hline
wir & können & müssen & dürfen & sollen & wollen & mögen \\
\hline
ihr & könnt & müsst & dürft & sollt & wollt & mögt \\
\hline
sie/Sie & können & müssen & dürfen & sollen & wollen & mögen \\
\hline
\end{tabular}
\end{center}

\begin{propriete}[Deux particularités de la conjugaison des modaux]
\begin{enumerate}
\item \textbf{Changement de voyelle au singulier} (ich, du, er/sie/es) : \all{können} $\rightarrow$ \all{kann} ; \all{müssen} $\rightarrow$ \all{muss} ; \all{dürfen} $\rightarrow$ \all{darf} ; \all{sollen} $\rightarrow$ \all{soll} ; \all{wollen} $\rightarrow$ \all{will} ; \all{mögen} $\rightarrow$ \all{mag}. Au pluriel (wir, ihr, sie/Sie), on retrouve la voyelle de l'infinitif.
\item \textbf{Pas de terminaison à la 3e personne du singulier} : on dit \all{er kann}, \all{er muss}, \all{er will}, jamais « er kannt » ni « er kannst ». La forme de \all{er} est identique à celle de \all{ich}.
\end{enumerate}
\end{propriete}

\begin{attention}[Erreurs fréquentes des francophones]
\begin{itemize}
\item Ne dis jamais « er kannst » ou « er musst » : la terminaison \all{-st} est réservée à \all{du}. Correct : \all{er kann}, \all{er muss}.
\item Attention à \all{wollen} : on dit \all{ich will} et \all{er will}, jamais « ich wolle ». La voyelle \all{o} devient \all{i} au singulier.
\item N'oublie pas le \all{ß} de \all{müssen} : \all{du musst} (avec \all{ss} car la voyelle est courte), mais \all{wir müssen} (avec \all{ß}).
\end{itemize}
\end{attention}

\courssub{La parenthèse verbale : le modal en 2e position, l'infinitif à la fin}

C'est LA règle capitale de cette leçon, et c'est elle qui déroute le plus les francophones.

\begin{propriete}[Règle de la parenthèse verbale (\all{Satzklammer})]
Dans une phrase déclarative avec un verbe de modalité :
\begin{enumerate}
\item le \textbf{modal conjugué} occupe la \textbf{2e position} (juste après le sujet ou le premier élément de la phrase) ;
\item l'\textbf{infinitif} se place \textbf{à la toute fin} de la phrase, quelle que soit sa longueur ;
\item tout le reste (compléments, adverbes) se loge \emph{entre} les deux verbes, comme dans une parenthèse.
\end{enumerate}
Exemple : \all{Die Tochter zieht den Stecker heraus, weil sie Energie sparen will.} « La fille débranche la prise parce qu'elle veut économiser de l'énergie. » Dans la subordonnée introduite par \all{weil}, le modal conjugué \all{will} passe lui aussi à la fin, juste après l'infinitif.
\end{propriete}

\begin{popfigure}
\begin{tikzpicture}[
box/.style={draw=popblue, fill=popblueL, rounded corners, align=center, font=\small, inner sep=4pt},
ibox/.style={draw=poporange, fill=poporangeL, rounded corners, align=center, font=\small, inner sep=4pt},
mid/.style={draw=popmuted, fill=popcream, rounded corners, align=center, font=\small, inner sep=4pt}]
\node[box, align=center] (s) at (0,0){\textbf{1} Sujet\\ \all{Wir}};
\node[box, fill=popgreenL, draw=popgreen, align=center] (m) at (2.6,0){\textbf{2} MODAL\\ \all{können}};
\node[mid, align=center] (c) at (6.2,0){\textbf{3} le reste\\ \all{heute im Haushalt}};
\node[ibox, align=center] (i) at (9.8,0){\textbf{dernière place} INFINITIF\\ \all{sparen.}};
\draw[-{Stealth[length=3mm]}, very thick, popgreen] (m.south) .. controls (5,-1.6) and (8,-1.6) .. node[below, font=\small, align=center, text=popdark]{la parenthèse verbale (\all{Satzklammer})} (i.south);
\end{tikzpicture}
\legende{La parenthèse verbale : \all{Wir können heute im Haushalt sparen.} « Nous pouvons économiser aujourd'hui dans la maison. » Le modal ouvre la parenthèse, l'infinitif la ferme.}
\end{popfigure}

Comparons maintenant les deux langues : c'est ici que tout se joue.

\begin{center}
\begin{tabularx}{\linewidth}{|X|X|X|}
\hline
\rowcolor{popblueL} \textbf{Français} & \textbf{Deutsch} & \textbf{Remarque} \\
\hline
Je dois payer la facture. & \all{Ich muss die Rechnung bezahlen.} & Le français conjugue « dois » et enchaîne « payer » ; l'allemand renvoie \all{bezahlen} à la fin. \\
\hline
Nous pouvons économiser de l'énergie. & \all{Wir können Energie sparen.} & « économiser de l'énergie » = un seul verbe en allemand : \all{Energie sparen}. \\
\hline
Elle veut acheter un panneau solaire. & \all{Sie will ein Solarpanel kaufen.} & Le complément \all{ein Solarpanel} se place avant l'infinitif. \\
\hline
Tu dois éteindre la lumière. & \all{Du musst das Licht ausschalten.} & \all{ausschalten} est un verbe à particule : il reste entier à l'infinitif final. \\
\hline
\end{tabularx}
\end{center}

\begin{attention}[Le calque du français : l'erreur numéro un]
En français, l'infinitif suit immédiatement le verbe conjugué : « je peux économiser ». Beaucoup d'élèves écrivent donc « Ich kann sparen Energie ». C'est \textbf{faux} ! En allemand, l'infinitif va \textbf{en dernier} : \all{Ich kann Energie sparen.} Astuce : prononce d'abord le modal, mets tout le reste au milieu, et garde l'infinitif pour la fin, comme la cerise sur le gâteau. Autre exemple : \all{Er will morgen ein Solarpanel kaufen.} « Il veut acheter un panneau solaire demain. »
\end{attention}

\begin{propriete}[Le modal dans la question]
Dans une question sans mot interrogatif, le modal conjugué passe en \textbf{première position} (comme tout verbe conjugué allemand), et l'infinitif reste à la fin : \all{Kannst du das Licht ausschalten?} « Peux-tu éteindre la lumière ? » Avec un mot interrogatif, celui-ci occupe la première position : \all{Was müssen wir tun?} « Que devons-nous faire ? »
\end{propriete}

\courssub{\all{möchte} : le souhait poli}

Le verbe \all{mögen} possède une forme spéciale, extrêmement fréquente : \all{möchte}. C'est le \textbf{Konjunktiv II} de \all{mögen}, et il sert à exprimer un souhait de manière polie, comme le conditionnel français « je voudrais ».

\begin{exemplebox}[\all{möchte} en action]
\all{Ich möchte Energie sparen.} « Je voudrais économiser de l'énergie. »\par
\all{Möchtest du eine Solaranlage?} « Voudrais-tu une installation solaire ? »\par
\all{Wir möchten in Yaoundé weniger Strom verbrauchen.} « Nous voudrions consommer moins de courant à Yaoundé. »
\end{exemplebox}

Conjugaison de \all{möchte} : \all{ich möchte, du möchtest, er/sie/es möchte, wir möchten, ihr möchtet, sie/Sie möchten}. La règle de la parenthèse verbale s'applique exactement de la même façon : \all{Ich möchte die Rechnung sofort bezahlen.} « Je voudrais payer la facture tout de suite. »

\begin{aretenir}[\all{mögen} ou \all{möchte} ?]
\all{mögen} (au présent) signifie « aimer » et s'emploie surtout avec un nom : \all{Ich mag Suppe.} « J'aime la soupe. » Devant un infinitif, pour exprimer un souhait poli, utilise toujours \all{möchte} : \all{Ich möchte Strom sparen.} « Je voudrais économiser le courant. »
\end{aretenir}

\courssub{Pour aller plus loin : le prétérit des modaux (utile pour le Bac)}

Pour raconter au passé, les modaux s'emploient très souvent au \textbf{prétérit}. Cette notion sera approfondie en classe de Terminale, mais il est bon de la reconnaître dès maintenant dans les textes.

\begin{center}
\begin{tabular}{|l|l|l|}
\hline
\rowcolor{popgoldL} \textbf{Présent} & \textbf{Präteritum} & \textbf{Exemple au passé} \\
\hline
kann & konnte & \all{Ich konnte nicht kommen.} « Je n'ai pas pu venir. » \\
\hline
muss & musste & \all{Wir mussten Strom sparen.} « Nous devions économiser le courant. » \\
\hline
darf & durfte & \all{Sie durfte nicht fernsehen.} « Elle n'avait pas le droit de regarder la télé. » \\
\hline
soll & sollte & \all{Er sollte die Rechnung bezahlen.} « Il était censé payer la facture. » \\
\hline
will & wollte & \all{Ich wollte ein Solarpanel kaufen.} « Je voulais acheter un panneau solaire. » \\
\hline
mag & mochte & \all{Sie mochte das neue Haus.} « Elle aimait la nouvelle maison. » \\
\hline
\end{tabular}
\end{center}

Bonne nouvelle : au prétérit, les modaux se conjuguent \textbf{régulièrement}, sans changement de voyelle : \all{ich konnte, du konntest, er konnte, wir konnten, ihr konntet, sie konnten}. La parenthèse verbale reste identique : \all{Früher musste man viel für Strom bezahlen.} « Autrefois, on devait payer beaucoup pour le courant. »

\courssub{Texte d'application : les modaux en contexte}

\begin{textebox}[Familie Ndi spart Energie]
Familie Ndi wohnt in Yaoundé. Die Stromrechnung ist jeden Monat sehr hoch. „Wir müssen etwas ändern“, sagt der Vater. „Wir können nicht so viel Geld für Strom ausgeben.“ Die Tochter Amina will helfen. „Ich kann das Licht ausschalten, wenn ich aus dem Zimmer gehe“, sagt sie. „Und ich möchte den Computer nicht die ganze Nacht laufen lassen.“ Der Sohn Junior soll den Stecker vom Fernseher herausziehen, aber er will das nicht: „Ich will doch Fußball sehen!“, ruft er. Die Mutter lacht: „Du darfst Fußball sehen, aber du musst danach den Fernseher ausschalten. Wir sollen alle zusammenarbeiten.“ Am Ende des Monats ist die Familie glücklich: die Rechnung ist viel kleiner. „Seht ihr“, sagt der Vater, „wir können alle etwas tun. Und nächstes Jahr möchte ich ein Solarpanel kaufen. Dann können wir die Sonne nutzen!“
\end{textebox}

\textbf{Glossaire} : \all{die Stromrechnung, -en} la facture d'électricité ; \all{ändern} changer ; \all{ausgeben (gibt aus)} dépenser ; \all{ausschalten} éteindre ; \all{laufen lassen} laisser marcher ; \all{der Stecker, -} la fiche (prise mâle) ; \all{herausziehen} débrancher, retirer ; \all{der Fernseher, -} le téléviseur ; \all{zusammenarbeiten} travailler ensemble ; \all{nutzen} utiliser.

\textbf{Questions de compréhension} :
\begin{enumerate}
\item \all{Warum muss Familie Ndi etwas ändern?} (Pourquoi la famille Ndi doit-elle changer quelque chose ?)
\item \all{Was will Amina tun?} (Que veut faire Amina ?)
\item \all{Was darf Junior? Was muss er danach tun?} (Qu'a le droit de faire Junior ? Que doit-il faire ensuite ?)
\item \all{Was möchte der Vater nächstes Jahr kaufen?} (Que veut acheter le père l'année prochaine ?)
\end{enumerate}

\courssub{Exercice résolu et entraînement}

\begin{exoresolu}[Conjugue et place l'infinitif correctement]
Énoncé : mets le modal à la bonne forme et construis la phrase complète.\par
1. \all{ich / können / Energie sparen} \par
2. \all{er / müssen / die Rechnung bezahlen} \par
3. \all{wir / sollen / die Umwelt schützen} \par
4. \all{du / wollen / ein Solarpanel kaufen} \par
5. Traduis : « Tu dois éteindre la lumière. »
\tcblower
Solution détaillée :\par
1. \all{Ich kann Energie sparen.} — \all{ich} $\rightarrow$ \all{kann} (changement de voyelle ö $\rightarrow$ a) ; l'infinitif \all{sparen} à la fin.\par
2. \all{Er muss die Rechnung bezahlen.} — \all{er} $\rightarrow$ \all{muss}, sans \all{-st} ni \all{-t} ; \all{bezahlen} en dernière position.\par
3. \all{Wir sollen die Umwelt schützen.} — au pluriel, on retrouve l'infinitif \all{sollen}.\par
4. \all{Du willst ein Solarpanel kaufen.} — \all{du} $\rightarrow$ \all{willst} (o $\rightarrow$ i + \all{-st}).\par
5. \all{Du musst das Licht ausschalten.} — attention au calque français : pas « Du musst ausschalten das Licht » ; le verbe à particule \all{ausschalten} reste entier à la fin.
\end{exoresolu}

\begin{exercice}[À toi de jouer]
1. Conjugue : \all{dürfen} (du), \all{wollen} (er), \all{mögen} (ihr), \all{können} (sie, pluriel).\par
2. Remets les mots dans l'ordre : \all{sparen / wir / können / Strom / im Haushalt}.\par
3. Traduis en allemand : a) « Nous devons protéger la Terre. » b) « Elle veut économiser de l'énergie. » c) « Je voudrais utiliser le soleil. »\par
4. Trouve l'erreur et corrige : \all{Er kannst das Licht ausschalten.} / \all{Ich muss bezahlen die Rechnung.}\par
5. Traduis en allemand : « Il est interdit de jeter des déchets ici. » (utilise \all{dürfen} + négation ; « jeter des déchets » = \all{Müll wegwerfen}).
\end{exercice}

\begin{corrige}[Exercice « À toi de jouer »]
1. \all{du darfst} ; \all{er will} ; \all{ihr mögt} ; \all{sie können}.\par
2. \all{Wir können Strom im Haushalt sparen.} « Nous pouvons économiser le courant dans la maison. »\par
3. a) \all{Wir müssen die Erde schützen.} (ou \all{Wir sollen die Erde schützen.} pour un devoir moral) ; b) \all{Sie will Energie sparen.} ; c) \all{Ich möchte die Sonne nutzen.}\par
4. \all{Er kann das Licht ausschalten.} (pas de \all{-st} à la 3e personne) ; \all{Ich muss die Rechnung bezahlen.} (l'infinitif à la fin).\par
5. \all{Man darf hier keinen Müll wegwerfen.} « Il est interdit de jeter des déchets ici » : « il ne faut pas » se rend par \all{nicht dürfen}, jamais par \all{nicht müssen}.
\end{corrige}

\begin{aretenir}[L'essentiel de la section]
\begin{itemize}
\item Six modaux : \all{können} (pouvoir), \all{müssen} (devoir, obligation), \all{dürfen} (avoir le droit), \all{sollen} (devoir moral, conseil), \all{wollen} (vouloir), \all{mögen}/\all{möchte} (aimer / vouloir poliment).
\item Au singulier, la voyelle change (\all{kann, muss, darf, soll, will, mag}) et la 3e personne ne prend aucune terminaison : \all{er kann}.
\item La parenthèse verbale : modal conjugué en 2e position, infinitif à la toute fin : \all{Wir können alle etwas tun.}
\item « Il ne faut pas » = \all{nicht dürfen} ; \all{nicht müssen} = « ne pas être obligé ».
\item \all{möchte} = souhait poli : \all{Ich möchte Energie sparen.}
\item Au passé, retiens le prétérit : \all{konnte, musste, durfte, sollte, wollte, mochte}.
\end{itemize}
\end{aretenir}

Maintenant que tu sais exprimer ce que l'on \emph{peut}, \emph{doit} et \emph{veut} faire pour l'énergie, la section suivante te permettra d'expliquer \emph{pourquoi} on le fait : nous étudierons l'expression de la finalité avec \all{um ... zu} et \all{damit}.

\coursec{Grammaire 2 : la finalité (um ... zu / damit)}

Dans la section précédente, nous avons étudié les verbes de modalité (\all{können, müssen, sollen, dürfen, wollen}), qui permettent d'exprimer la capacité, l'obligation ou le souhait. Mais quand on parle de la gestion de l'énergie, on ne dit pas seulement ce que l'on \emph{peut} ou \emph{doit} faire : on explique aussi \emph{pourquoi} on le fait, dans quel \emph{but}. « J'éteins la lumière \textbf{pour} économiser de l'énergie. » Cette expression du but, c'est la \cle{finalité}. En allemand, elle se construit de deux façons différentes selon une question très simple : le sujet des deux actions est-il le même ? C'est toute la logique de \all{um ... zu} et \all{damit}, que nous allons maintenant découvrir pas à pas.

\courssub{Observation : partons du texte}

Reprenons deux phrases tirées de notre texte support \emph{Energie sparen im Alltag} et observons-les attentivement.

\begin{exemplebox}[Deux phrases à observer]
\all{Die Tochter zieht den Stecker heraus, um Energie zu sparen.} « La fille débranche la prise pour économiser de l'énergie. »

\all{Die Mutter kauft eine Solaranlage, damit die Familie die Sonne nutzen kann.} « La mère achète une installation solaire pour que la famille puisse utiliser le soleil. »
\end{exemplebox}

Posons-nous trois questions :

\begin{enumerate}
\item Dans la première phrase, qui débranche la prise ? \emph{Die Tochter}. Et qui économise l'énergie ? \emph{Die Tochter} aussi. Le sujet est \textbf{identique} pour les deux actions.
\item Dans la deuxième phrase, qui achète l'installation solaire ? \emph{Die Mutter}. Et qui utilise le soleil ? \emph{Die Familie}. Les sujets sont \textbf{différents}.
\item Où se trouve le verbe dans chaque cas ? Dans la première phrase, l'infinitif \all{zu sparen} est à la toute fin. Dans la deuxième, le verbe conjugué \all{kann} est lui aussi à la fin, mais il est précédé d'un sujet, \all{die Familie}.
\end{enumerate}

Nous tenons la règle : le choix entre \all{um ... zu} et \all{damit} dépend uniquement de l'identité ou de la différence des sujets.

\courssub{Règle 1 : um ... zu + infinitif (même sujet)}

\begin{propriete}[La finalité avec um ... zu]
Quand le sujet des deux actions est \textbf{le même}, on exprime le but avec la construction \all{um ... zu + Infinitiv} :

\begin{itemize}
\item \all{um} ouvre le groupe de finalité, généralement après une virgule ;
\item \all{zu} se place \textbf{juste avant l'infinitif} ;
\item l'infinitif précédé de \all{zu} occupe la \textbf{dernière place} de la phrase ;
\item on ne répète \textbf{jamais} le sujet dans le groupe \all{um ... zu}.
\end{itemize}

Formule : \all{Sujet + verbe conjugué + ..., um + compléments + zu + Infinitiv.}
\end{propriete}

\begin{exemplebox}[Exemples traduits]
\all{Ich schalte das Licht ab, um Strom zu sparen.} « J'éteins la lumière pour économiser l'électricité. »

\all{Sie kauft eine Solaranlage, um die Sonne zu nutzen.} « Elle achète une installation solaire pour utiliser le soleil. »

\all{Wir gehen zu Fuß zur Schule, um Energie zu sparen.} « Nous allons à l'école à pied pour économiser de l'énergie. »

\all{Er öffnet das Fenster, um frische Luft hereinzulassen.} « Il ouvre la fenêtre pour laisser entrer l'air frais. »
\end{exemplebox}

Remarquez que le français utilise simplement « pour + infinitif » : la structure est donc très proche, ce qui est une bonne nouvelle pour les francophones. La seule différence, mais elle est capitale, c'est la présence obligatoire de \all{zu} devant l'infinitif allemand.

\courssub{Règle 2 : damit + subordonnée (sujets différents)}

\begin{propriete}[La finalité avec damit]
Quand les sujets des deux actions sont \textbf{différents}, on ne peut pas utiliser \all{um ... zu}. On construit alors une \textbf{proposition subordonnée} introduite par la conjonction \all{damit} (« pour que, afin que ») :

\begin{itemize}
\item \all{damit} est une conjonction de subordination : le \textbf{verbe conjugué passe en position finale} ;
\item la subordonnée possède son propre sujet, placé juste après \all{damit} ;
\item la subordonnée est séparée de la principale par une virgule ;
\item le verbe de la subordonnée se conjugue à l'\textbf{indicatif} (l'allemand n'a pas de subjonctif après \all{damit}, contrairement au français).
\end{itemize}

Formule : \all{Sujet 1 + verbe + ..., damit + Sujet 2 + compléments + verbe conjugué.}
\end{propriete}

\begin{exemplebox}[Exemples traduits]
\all{Der Vater arbeitet viel, damit die Familie die Rechnung bezahlen kann.} « Le père travaille beaucoup pour que la famille puisse payer la facture. »

\all{Die Lehrerin erklärt die Regel, damit die Schüler sie verstehen.} « La professeure explique la règle pour que les élèves la comprennent. »

\all{Ich gebe dir eine Lampe, damit du im Dunkeln lesen kannst.} « Je te donne une lampe pour que tu puisses lire dans le noir. »

\all{Die Regierung baut Wasserkraftwerke, damit die Dörfer Strom bekommen.} « Le gouvernement construit des centrales hydroélectriques pour que les villages reçoivent de l'électricité. »
\end{exemplebox}

\begin{attention}[Ne pas confondre les deux damit]
Le mot \all{damit} existe aussi comme \textbf{pronom adverbial} signifiant « avec cela, avec ça » : \all{Ich schreibe damit.} « J'écris avec ça. » Dans ce cas, le verbe reste à la deuxième place et il n'y a pas de subordonnée. Le \all{damit} de finalité, lui, introduit toujours une proposition complète avec un sujet et un verbe conjugué renvoyé à la fin : \all{..., damit du Strom sparst.} Le contexte permet de ne jamais les confondre.
\end{attention}

\courssub{Cas particulier 1 : les verbes à particule séparable}

Que se passe-t-il quand l'infinitif est un verbe séparable comme \all{ausschalten} (éteindre), \all{herausziehen} (débrancher) ou \all{aufstehen} (se lever) ? La particule et le verbe se réunissent, et \all{zu} vient \textbf{s'insérer entre les deux}, le tout s'écrivant \textbf{en un seul mot}.

\begin{propriete}[zu et les verbes séparables]
Avec un verbe à particule séparable, \all{zu} se place entre le préfixe et le radical, et l'ensemble forme un seul mot :

\all{ausschalten} $\rightarrow$ \all{aus\textbf{zu}schalten} ; \all{herausziehen} $\rightarrow$ \all{heraus\textbf{zu}ziehen} ; \all{einkaufen} $\rightarrow$ \all{ein\textbf{zu}kaufen}.
\end{propriete}

\begin{exemplebox}[Exemples traduits]
\all{Sie geht ins Zimmer, um das Licht auszuschalten.} « Elle entre dans la chambre pour éteindre la lumière. »

\all{Die Tochter zieht den Stecker heraus, um das Gerät auszuschalten.} « La fille débranche la prise pour éteindre l'appareil. »

\all{Ich stehe früh auf, um das Wasser für den Tee zu kochen.} « Je me lève tôt pour faire bouillir l'eau du thé. »
\end{exemplebox}

\begin{attention}[Un seul mot !]
Erreur fréquente des francophones : écrire \all{um das Licht aus zu schalten} en trois morceaux. C'est faux. À l'infinitif avec \all{zu}, le verbe séparable se soude : \all{auszuschalten}, \all{herauszuziehen}, \all{einzukaufen}. Un seul mot, sans espace.
\end{attention}

\courssub{Cas particulier 2 : um ... zu avec un verbe de modalité}

On peut tout à fait exprimer un but qui contient un modal (« pour \emph{pouvoir} économiser »). Dans ce cas, le modal se met lui aussi à l'infinitif et se place \textbf{après} l'infinitif principal, \all{zu} restant devant le modal.

\begin{exemplebox}[um ... zu + modal]
\all{Wir bauen ein Wasserkraftwerk, um Strom sparen zu können.} « Nous construisons une centrale hydroélectrique pour pouvoir économiser l'électricité. »

\all{Er lernt Deutsch, um in Deutschland studieren zu können.} « Il apprend l'allemand pour pouvoir étudier en Allemagne. »
\end{exemplebox}

Ordre à retenir : \all{um + compléments + infinitif principal + zu + modal}. Ce cas de figure est un \textbf{complément de niveau Bac} : il n'est pas exigé en Première, mais le reconnaître à la lecture est un vrai atout.

\courssub{Comparaison avec le français}

Le français et l'allemand fonctionnent de manière très semblable pour exprimer le but, ce qui facilite l'apprentissage. Voici le tableau de correspondance :

\begin{center}
\begin{tabularx}{\linewidth}{|X|X|X|}
\hline
\rowcolor{popblueL} Français & Deutsch & Remarque \\
\hline
pour économiser & um zu sparen & même sujet ; \all{zu} obligatoire \\
\hline
pour éteindre la lumière & um das Licht auszuschalten & verbe séparable : un seul mot \\
\hline
pour que tu économises & damit du sparst & sujets différents ; verbe conjugué à la fin \\
\hline
pour que la famille puisse payer & damit die Familie bezahlen kann & le modal \all{kann} se place après l'infinitif \\
\hline
afin d'utiliser le soleil & um die Sonne zu nutzen & « afin de » = registre soutenu de « pour » \\
\hline
\end{tabularx}
\end{center}

\begin{attention}[Les deux erreurs classiques des francophones]
Première erreur : oublier \all{zu}. La phrase \all{Ich schalte das Licht ab, um Energie sparen.} est \textbf{fausse} : il faut \all{um Energie \textbf{zu} sparen}. En français, « pour » suffit ; en allemand, \all{um} ne suffit jamais seul devant un infinitif.

Deuxième erreur : mettre un sujet après \all{um ... zu}. On ne dit \textbf{jamais} \all{um ich zu sparen}. Dès qu'il y a un deuxième sujet, on doit passer à \all{damit} : \all{..., damit ich Strom spare.}
\end{attention}

\courssub{Schéma de décision : um ... zu ou damit ?}

La figure suivante résume la démarche à suivre chaque fois que vous voulez exprimer un but en allemand.

\begin{popfigure}
\begin{center}
\begin{tikzpicture}[
node distance=0.9cm and 1.6cm,
question/.style={diamond, draw=popblue, fill=popblueL, align=center, inner sep=1.5pt, aspect=2.2, font=\small},
result/.style={rectangle, rounded corners=4pt, draw=popgreen, fill=popgreenL, align=center, font=\small, text width=4.6cm, inner sep=5pt},
lbl/.style={font=\small\bfseries, text=popdark}
]
\node[question, align=center] (q){Même sujet\\pour les deux actions ?};
\node[result, below left=1.1cm and 0.8cm of q, align=center] (oui){\all{um ... zu + Infinitiv}\\[2pt] \all{Ich schalte das Licht ab, um Strom zu sparen.}};
\node[result, below right=1.1cm and 0.8cm of q, align=center] (non){\all{damit + verbe conjugué final}\\[2pt] \all{Der Vater arbeitet, damit die Familie zahlen kann.}};
\draw[-{Stealth}, thick, popblue] (q) -- node[lbl, left] {oui} (oui);
\draw[-{Stealth}, thick, poporange] (q) -- node[lbl, right] {non} (non);
\end{tikzpicture}
\end{center}
\legende{Organigramme de choix : même sujet $\rightarrow$ \all{um ... zu} ; sujets différents $\rightarrow$ \all{damit}}
\end{popfigure}

\courssub{Texte d'application : la finalité en contexte}

Lisons maintenant un court texte qui réutilise \all{um ... zu} et \all{damit} dans un contexte camerounais. Repérez chaque construction de finalité et justifiez le choix fait par l'auteur.

\begin{textebox}[Eine Solaranlage für die Schule]
In einem Dorf bei Yaoundé gibt es eine neue Solaranlage. Die Schüler und Schülerinnen kommen jetzt früh in die Schule, um die neuen Computer zu benutzen. Der Direktor hat die Anlage gekauft, damit die Klassenzimmer am Abend Licht haben. Viele Eltern helfen mit, um die Rechnung zu bezahlen. Eine Lehrerin erklärt den Kindern jeden Tag, wie man Energie spart: „Schaltet die Lampen aus, um Strom zu sparen!“, sagt sie. Die Familien im Dorf kaufen jetzt auch kleine Solarpanels, damit ihre Handys immer geladen sind. Ein Junge namens Samuel lernt besonders viel, um später Ingenieur zu werden. Sein Vater arbeitet auf dem Markt, damit Samuel die Schule besuchen kann. Alle sind stolz: das Dorf nutzt die Sonne, um eine bessere Zukunft zu bauen. Die Lehrerin sagt am Ende der Stunde: „Wir sparen Energie, damit unsere Kinder morgen besser leben.“ Und alle Schüler antworten: „Wir schalten das Licht aus, um die Erde zu schützen!“
\end{textebox}

\begin{definition}[Glossaire]
\all{die Solaranlage, -n} : l'installation solaire ; \all{das Dorf, die Dörfer} : le village ; \all{der Direktor, -en} : le directeur ; \all{die Rechnung, -en} : la facture ; \all{mithelfen (half mit, hat mitgeholfen)} : aider, participer ; \all{das Solarpanel, -s} : le panneau solaire ; \all{geladen} : chargé ; \all{der Ingenieur, -e} : l'ingénieur ; \all{stolz} : fier ; \all{die Zukunft} : l'avenir ; \all{die Erde, -n} : la Terre ; \all{schützen} : protéger.
\end{definition}

Questions de repérage : (1) Combien de constructions \all{um ... zu} le texte contient-il ? (2) Combien de subordonnées avec \all{damit} ? (3) Pour chaque \all{damit}, identifiez les deux sujets différents. (4) Trouvez le verbe séparable du texte et dites sous quelle forme il apparaît.

\begin{corrige}[Questions de repérage]
(1) Six constructions \all{um ... zu} : \all{um die neuen Computer zu benutzen}, \all{um die Rechnung zu bezahlen}, \all{um Strom zu sparen}, \all{um später Ingenieur zu werden}, \all{um eine bessere Zukunft zu bauen}, \all{um die Erde zu schützen}. (2) Quatre subordonnées \all{damit} : \all{damit die Klassenzimmer am Abend Licht haben}, \all{damit ihre Handys immer geladen sind}, \all{damit Samuel die Schule besuchen kann}, \all{damit unsere Kinder morgen besser leben}. (3) Exemples : \emph{der Direktor} achète, mais ce sont \emph{die Klassenzimmer} qui ont de la lumière ; \emph{der Vater} travaille, mais c'est \emph{Samuel} qui peut aller à l'école. (4) Le verbe séparable : \all{ausschalten}. Il apparaît deux fois à l'impératif (\all{Schaltet die Lampen aus}, \all{Wir schalten das Licht aus}), donc sans \all{zu} ; à l'infinitif avec \all{zu}, on écrirait \all{auszuschalten} en un seul mot.
\end{corrige}

\courssub{Exercice résolu}

\begin{exoresolu}[um ... zu ou damit ?]
Complétez avec \all{um ... zu} ou \all{damit} et mettez le verbe à la bonne forme et à la bonne place.

1. Ich kaufe eine Lampe, \ldots\ Strom \ldots\ (sparen). (sujet : \emph{ich})
2. Der Lehrer spricht langsam, \ldots\ die Schüler alles \ldots\ (verstehen).
3. Wir fahren mit dem Fahrrad, \ldots\ die Umwelt \ldots\ (schützen).
4. Die Mutter gibt dem Sohn Geld, \ldots\ er eine Taschenlampe \ldots\ (kaufen).
5. Sie geht in die Küche, \ldots\ das Radio \ldots\ (ausschalten).

\tcblower

\textbf{Solution détaillée.}

1. Même sujet (\emph{ich} dans les deux cas) $\rightarrow$ \all{um ... zu} : \all{Ich kaufe eine Lampe, \textbf{um} Strom \textbf{zu sparen}.} « J'achète une lampe pour économiser l'électricité. »

2. Sujets différents (\emph{der Lehrer} / \emph{die Schüler}) $\rightarrow$ \all{damit} + verbe conjugué final : \all{Der Lehrer spricht langsam, \textbf{damit} die Schüler alles \textbf{verstehen}.} « Le professeur parle lentement pour que les élèves comprennent tout. »

3. Même sujet (\emph{wir}) $\rightarrow$ \all{um ... zu} : \all{Wir fahren mit dem Fahrrad, \textbf{um} die Umwelt \textbf{zu schützen}.} « Nous roulons à vélo pour protéger l'environnement. »

4. Sujets différents (\emph{die Mutter} / \emph{er}) $\rightarrow$ \all{damit} : \all{Die Mutter gibt dem Sohn Geld, \textbf{damit} er eine Taschenlampe \textbf{kauft}.} « La mère donne de l'argent au fils pour qu'il achète une lampe de poche. » Attention : \all{kaufen} est conjugué à la troisième personne du singulier et renvoyé à la fin.

5. Même sujet (\emph{sie}) et verbe séparable $\rightarrow$ \all{zu} s'insère dans le verbe : \all{Sie geht in die Küche, \textbf{um} das Radio \textbf{auszuschalten}.} « Elle va dans la cuisine pour éteindre la radio. » Un seul mot : \all{auszuschalten}.
\end{exoresolu}

\begin{aretenir}[L'essentiel de la finalité]
\begin{itemize}
\item But avec \textbf{même sujet} : \all{um ... zu + Infinitiv}, \all{zu} juste avant l'infinitif, en fin de phrase : \all{..., um Energie zu sparen.}
\item But avec \textbf{sujets différents} : \all{damit} + subordonnée, verbe conjugué \textbf{en position finale} : \all{..., damit die Familie spart.}
\item Verbe séparable : \all{zu} s'insère et tout se soude en un mot : \all{auszuschalten}.
\item Avec un modal (niveau Bac) : \all{um Strom sparen zu können}.
\item Correspondances : « pour + infinitif » = \all{um ... zu} ; « pour que + subjonctif » = \all{damit} + indicatif.
\item Jamais de sujet après \all{um ... zu} ; jamais \all{um} sans \all{zu} devant un infinitif.
\end{itemize}
\end{aretenir}

Vous savez désormais exprimer le but de vos actions en allemand, une compétence indispensable pour parler de la gestion de l'énergie : on n'éteint pas la lumière par hasard, mais \all{um Strom zu sparen} ! Dans la dernière partie de la leçon, nous mettrons ce point de grammaire en pratique dans la méthode du commentaire de texte et dans une production guidée.

\coursec{Méthode du commentaire et production guidée}

Dans la section précédente, tu as appris à exprimer la finalité avec \all{um ... zu} et \all{damit}. Ces constructions, comme les verbes de modalité, ne sont pas seulement des règles à connaître : ce sont des \cle{outils} pour parler et écrire. Cette dernière section te montre comment les utiliser dans les deux grandes tâches de l'examen : le \cle{commentaire de texte} et la \cle{production écrite guidée}. Tu vas d'abord apprendre une méthode en quatre étapes, lire un commentaire modèle, puis rédiger et présenter ton propre texte sur l'économie d'énergie.

\courssub{La méthode du commentaire de texte en quatre étapes}

Au lycée, commenter un texte allemand ne signifie pas le traduire. Commenter, c'est \emph{présenter}, \emph{résumer}, \emph{analyser} et \emph{donner son avis}. Observe d'abord la frise suivante : elle résume le parcours complet.

\begin{popfigure}
\begin{center}
\begin{tikzpicture}[node distance=0.6cm, every node/.style={align=center}]
\node[draw=popblue, fill=popblueL, rounded corners, minimum width=3cm, minimum height=1.4cm, text width=2.8cm] (e1) {\textbf{1. Introduction}\\ \all{Der Text handelt von...}};
\node[draw=popgreen, fill=popgreenL, rounded corners, minimum width=3cm, minimum height=1.4cm, text width=2.8cm, right=of e1] (e2) {\textbf{2. Résumé}\\ \all{Die Familie will...}};
\node[draw=poporange, fill=poporangeL, rounded corners, minimum width=3cm, minimum height=1.4cm, text width=2.8cm, right=of e2] (e3) {\textbf{3. Analyse}\\ idées + vocabulaire + grammaire};
\node[draw=poppurple, fill=poppurpleL, rounded corners, minimum width=3cm, minimum height=1.4cm, text width=2.8cm, right=of e3] (e4) {\textbf{4. Conclusion}\\ \all{Meiner Meinung nach...}};
\draw[-{Stealth}, very thick, popdark] (e1) -- (e2);
\draw[-{Stealth}, very thick, popdark] (e2) -- (e3);
\draw[-{Stealth}, very thick, popdark] (e3) -- (e4);
\end{tikzpicture}
\end{center}
\legende{Les quatre étapes du commentaire de texte : de la présentation à l'avis personnel.}
\end{popfigure}

\begin{methode}[Commenter un texte en quatre étapes]
\begin{enumerate}
\item \textbf{Présenter le texte en deux phrases} : indique le thème, la source (un journal, un magazine, un dialogue) et les personnages. Utilise : \all{Der Text handelt von...} (le texte parle de...), \all{Der Text kommt aus einer Zeitung.} (le texte vient d'un journal).
\item \textbf{Résumer en trois phrases} : raconte l'essentiel avec tes propres mots, au présent. Qui fait quoi, et pourquoi ? Utilise : \all{Die Familie will...}, \all{Der Vater schaltet das Licht ab.}
\item \textbf{Analyser} : relève deux ou trois idées principales et illustre-les avec le vocabulaire du texte (\all{der Strom, die Steckdose, sparen}) et les structures grammaticales remarquables (modaux, \all{um ... zu}). Utilise : \all{Der Autor benutzt die Wörter...}, \all{Man findet hier das Modalverb „müssen“.}
\item \textbf{Conclure avec ton avis personnel} : dis ce que tu penses du texte et fais un lien avec ta vie, par exemple au Cameroun. Utilise : \all{Meiner Meinung nach...}, \all{Ich finde es wichtig, dass...}, \all{Ich finde den Text interessant, weil...}
\end{enumerate}
\end{methode}

\courssub{Phrases toutes faites pour le commentaire}

Ces phrases-modèles (\emph{Redemittel}) sont tes amies : apprends-les par cœur et réutilise-les avec n'importe quel texte.

\begin{center}
\begin{tabularx}{\linewidth}{|X|X|}
\hline
\rowcolor{popblueL} \textbf{Deutsch} & \textbf{Français} \\
\hline
\all{Der Text handelt von einer Familie.} & Le texte parle d'une famille. \\
\hline
\all{Die Familie will weniger Strom verbrauchen.} & La famille veut consommer moins d'électricité. \\
\hline
\all{Der Autor erklärt, wie man Energie sparen kann.} & L'auteur explique comment on peut économiser l'énergie. \\
\hline
\all{Man findet hier viele Modalverben.} & On trouve ici beaucoup de verbes de modalité. \\
\hline
\all{Ich finde den Text interessant.} & Je trouve le texte intéressant. \\
\hline
\all{Meiner Meinung nach müssen wir alle die Umwelt schützen.} & À mon avis, nous devons tous protéger l'environnement. \\
\hline
\all{Ich finde es wichtig, dass wir Strom sparen.} & Je trouve important que nous économisions l'électricité. \\
\hline
\end{tabularx}
\end{center}

\begin{exemplebox}[Deux phrases modèles traduites]
\all{Der Text handelt von einer Familie, die Strom sparen will.} — « Le texte parle d'une famille qui veut économiser l'électricité. » Observe la place du verbe \all{will} à la fin de la subordonnée relative introduite par \all{die}.

\all{Meiner Meinung nach müssen wir alle die Umwelt schützen.} — « À mon avis, nous devons tous protéger l'environnement. » Observe : après l'expression d'opinion en tête de phrase, le verbe conjugué \all{müssen} vient en deuxième position, puis le sujet \all{wir}.
\end{exemplebox}

\begin{attention}[Le piège de \all{ich finde, dass...}]
Après \all{dass}, le verbe conjugué passe \textbf{en fin de proposition}. C'est l'erreur numéro un des francophones, car en français l'ordre ne change pas.

\all{Ich finde, dass wir Energie sparen müssen.} — « Je trouve que nous devons économiser l'énergie. »

Ne dis jamais : \emph{Ich finde, dass wir müssen Energie sparen.} Le modal conjugué (\all{müssen}) doit aller tout à la fin, et l'infinitif (\all{sparen}) se place juste avant lui. Même règle après \all{weil} : \all{Ich finde den Text interessant, weil wir in Kamerun auch Energie sparen müssen.}
\end{attention}

\courssub{Commentaire modèle}

Lis ce commentaire du texte \all{Energie sparen im Alltag}. Il suit exactement les quatre étapes de la méthode. Repère-les pendant la lecture.

\begin{textebox}[Commentaire modèle : „Energie sparen im Alltag“]
\all{Der Text „Energie sparen im Alltag“ handelt von der Familie Berger aus München. Die Familie will weniger Strom verbrauchen, weil die Rechnung immer teurer wird. Der Vater schaltet das Licht ab, die Tochter zieht den Stecker heraus, und die Mutter kauft eine Solaranlage. Der Autor benutzt oft Modalverben wie „müssen“ und „sollen“, und er erklärt, wie man im Alltag Energie sparen kann. Ich finde den Text interessant, weil wir in Kamerun auch Energie sparen müssen. Meiner Meinung nach sollen alle Menschen die Erde schützen.}
\end{textebox}

\textbf{Analyse guidée du modèle.} Le commentaire compte \textbf{six phrases}. Phrase 1 : introduction (\all{Der Text handelt von...}). Phrases 2 et 3 : résumé au présent, avec une subordonnée en \all{weil}. Phrase 4 : analyse (vocabulaire et grammaire du texte). Phrases 5 et 6 : conclusion avec avis personnel (\all{Ich finde...}, \all{Meiner Meinung nach...}) et lien avec le Cameroun. Remarque la place des verbes : \all{wird}, \all{sparen kann}, \all{sparen müssen}, \all{schützen sollen} en fin de subordonnée.

\begin{savaistu}[Comment la production écrite est-elle notée au Bac ?]
Au baccalauréat, la production écrite d'allemand est notée sur trois critères : la \textbf{richesse du vocabulaire} (utilise les mots du thème : \all{die Energie, der Strom, sparen}), la \textbf{correction grammaticale} (place du verbe, modaux conjugués, majuscules aux noms) et le \textbf{respect de la consigne} (sujet traité, nombre de mots demandé). Dernier conseil de correcteur : garde toujours cinq minutes pour \cle{relire ta copie}. Une relecture attentive fait gagner des points !
\end{savaistu}

\courssub{Production écrite guidée}

\textbf{Sujet :} \all{Was kannst du tun, um Energie zu sparen?} (Que peux-tu faire pour économiser l'énergie ?)

\textbf{Consigne :} rédige cinq à huit phrases avec \textbf{au moins trois verbes de modalité} et \textbf{au moins une construction \all{um ... zu}}.

\begin{methode}[Planifier sa production en trois temps]
\begin{enumerate}
\item \textbf{Avant d'écrire} : liste trois ou quatre actions concrètes (éteindre la lumière, débrancher les appareils, prendre le vélo, fermer le robinet) et choisis tes modaux : \all{können} (pouvoir), \all{müssen} (devoir), \all{sollen} (devoir, recommandation), \all{wollen} (vouloir).
\item \textbf{En écrivant} : une idée par phrase. Modal conjugué en deuxième position, infinitif en fin de phrase : \all{Ich muss das Licht ausschalten.}
\item \textbf{Après avoir écrit} : vérifie avec la grille d'auto-évaluation ci-dessous.
\end{enumerate}
\end{methode}

\begin{exoresolu}[Un élève de Première rédige sa production]
Énoncé : écris une production de cinq à huit phrases sur le sujet \all{Was kannst du tun, um Energie zu sparen?} avec au moins trois modaux et une construction \all{um ... zu}.
\tcblower
\textbf{Solution rédigée (modèle) :}

\all{Ich kann viel tun, um Energie zu sparen. Zuerst muss ich das Licht ausschalten, wenn ich das Zimmer verlasse. Dann soll ich den Stecker herausziehen, weil die Geräte auch im Stand-by Strom verbrauchen. Ich will auch zu Fuß zur Schule gehen, um kein Benzin zu verbrauchen. Meiner Meinung nach müssen alle Schüler in Yaoundé die Umwelt schützen. Ich finde es wichtig, dass wir zusammen die Erde retten.}

« Je peux faire beaucoup pour économiser l'énergie. D'abord, je dois éteindre la lumière quand je quitte la pièce. Ensuite, je devrais débrancher la prise, parce que les appareils consomment du courant même en veille. Je veux aussi aller à l'école à pied, pour ne pas consommer d'essence. À mon avis, tous les élèves de Yaoundé doivent protéger l'environnement. Je trouve important que nous sauvions la Terre ensemble. »

\textbf{Vérification :} cinq modaux au total (\all{kann, muss, soll, will} en proposition principale ; \all{müssen} en subordonnée), deux constructions \all{um ... zu}, vocabulaire du thème (\all{die Energie, das Licht, der Stecker, der Strom, die Umwelt}), infinitifs en fin de phrase, verbe conjugué en fin de subordonnée (\all{verlasse, verbrauchen, müssen, retten}). Consigne respectée !
\end{exoresolu}

\courssub{Grille d'auto-évaluation}

Avant de rendre ta copie ou de passer à l'oral, coche chaque critère.

\begin{center}
\begin{tabularx}{\linewidth}{|X|l|}
\hline
\rowcolor{popblueL} \textbf{Critère} & \textbf{Oui / Non} \\
\hline
J'ai utilisé au moins trois mots du thème (\all{die Energie, der Strom, die Steckdose, sparen...}). & \\
\hline
J'ai employé au moins trois verbes de modalité correctement conjugués (\all{ich kann, ich muss, ich soll...}). & \\
\hline
L'infinitif est bien en fin de phrase (\all{Ich muss das Licht ausschalten.}). & \\
\hline
J'ai écrit au moins une construction \all{um ... zu} avec l'infinitif à la fin. & \\
\hline
Tous les noms commencent par une majuscule (\all{das Licht, die Schule, der Strom}). & \\
\hline
Après \all{weil} et \all{dass}, le verbe conjugué est à la fin. & \\
\hline
J'ai respecté le nombre de phrases demandé (cinq à huit). & \\
\hline
\end{tabularx}
\end{center}

\begin{attention}[Les trois fautes qui coûtent le plus de points]
\begin{itemize}
\item \textbf{Le modal mal placé} : \emph{Ich muss ausschalten das Licht} est faux. Dis : \all{Ich muss das Licht ausschalten.} (l'infinitif à la fin, la particule séparable recollée à l'infinitif : \all{ausschalten}).
\item \textbf{La minuscule aux noms} : en allemand, \all{das Licht, die Energie, der Strom} prennent toujours une majuscule, même au milieu de la phrase.
\item \textbf{Oublier le sujet} : en allemand, on ne peut pas sauter le sujet comme parfois en français parlé. Chaque phrase a besoin de son sujet : \all{ich, wir, man}.
\end{itemize}
\end{attention}

\courssub{Production orale : ta présentation en une minute}

\begin{experience}[Présente ton texte à la classe]
\begin{enumerate}
\item \textbf{Préparation (5 minutes)} : relis ta production écrite et souligne tes trois modaux et ta construction \all{um ... zu}. Apprends ta première et ta dernière phrase par cœur.
\item \textbf{Présentation (1 minute)} : devant la classe ou en binôme, présente ton texte en regardant ton public. Commence par \all{Ich kann viel tun, um Energie zu sparen.} et termine par ton avis : \all{Meiner Meinung nach...}
\item \textbf{Écoute active} : pendant la présentation d'un camarade, note un mot du vocabulaire du thème qu'il a bien utilisé et un modal. Tu lui donneras ce retour positif à la fin.
\item \textbf{Variante en binôme} : l'un pose la question \all{Was kannst du tun, um Energie zu sparen?}, l'autre répond en trois phrases, puis on échange les rôles.
\end{enumerate}
\end{experience}

\begin{aretenir}[L'essentiel de la section]
\begin{itemize}
\item Un commentaire de texte suit quatre étapes : \textbf{introduction} (\all{Der Text handelt von...}), \textbf{résumé} au présent, \textbf{analyse} (idées + vocabulaire + grammaire), \textbf{conclusion} avec avis personnel (\all{Meiner Meinung nach...}).
\item Après \all{dass} et \all{weil}, le verbe conjugué va en fin de proposition : \all{Ich finde, dass wir Energie sparen müssen.}
\item Une bonne production écrite contient le vocabulaire du thème, des modaux bien conjugués, l'infinitif en fin de phrase, des majuscules à tous les noms — et elle est toujours relue.
\item À l'oral comme à l'écrit, les constructions \all{um ... zu} et les modaux montrent que tu maîtrises la leçon : utilise-les !
\end{itemize}
\end{aretenir}

\begin{exercice}[À toi de jouer]
\begin{enumerate}
\item Remets dans l'ordre les quatre étapes du commentaire : Analyse — Introduction — Conclusion — Résumé.
\item Complète avec le modal qui convient : \all{Ich ... das Licht ausschalten, wenn ich gehe. (devoir)} ; \all{Wir ... die Umwelt schützen. (pouvoir)}.
\item Traduis en allemand : « Je trouve que nous devons économiser l'énergie. »
\item Rédige ta propre production : \all{Was kannst du tun, um Energie zu sparen?} (cinq à huit phrases, trois modaux, une construction \all{um ... zu}), puis auto-évalue-toi avec la grille.
\end{enumerate}
\end{exercice}

\begin{corrige}[Exercice : À toi de jouer]
\begin{enumerate}
\item Introduction — Résumé — Analyse — Conclusion.
\item \all{Ich muss das Licht ausschalten, wenn ich gehe.} ; \all{Wir können die Umwelt schützen.}
\item \all{Ich finde, dass wir Energie sparen müssen.} (verbe conjugué \all{müssen} à la fin, infinitif \all{sparen} juste avant).
\item Production individuelle : compare ta copie avec le modèle de l'exercice résolu et coche chaque ligne de la grille d'auto-évaluation.
\end{enumerate}
\end{corrige}

\newpage

\coursec{Activité d'intégration}

\coursec{Leçon ALA11 : La gestion de l'énergie — Activité d'intégration et production finale}

\courssub{Objectifs de l'activité}

Cette activité d'intégration clôt la leçon sur la gestion de l'énergie. Elle permet aux élèves de réinvestir, dans une tâche authentique et collective, l'ensemble des acquis de la leçon :

\begin{itemize}
\item \textbf{Lexique} : employer correctement au moins cinq mots du champ lexical de l'énergie (\all{die Energie}, \all{der Strom}, \all{die Steckdose}, \all{das Solarpanel}, \all{sparen}, \all{ausschalten}...) avec le bon article.
\item \textbf{Grammaire 1} : utiliser les verbes de modalité (\all{müssen}, \all{sollen}, \all{können}, \all{dürfen}, \all{wollen}) à la bonne place dans la phrase, avec l'infinitif en fin de proposition.
\item \textbf{Grammaire 2} : exprimer la finalité avec \all{um ... zu} et \all{damit}.
\item \textbf{Expression écrite} : rédiger des consignes courtes, claires et percutantes, adaptées au genre de l'affiche.
\item \textbf{Expression orale} : présenter un document collectif en deux minutes, avec une prononciation soignée.
\end{itemize}

\courssub{Situation}

Tu es élève en classe de Première A au Lycée de Nkol-Eton à Yaoundé. La direction du lycée constate que la facture d'électricité de l'établissement augmente chaque mois : les lumières restent allumées dans les salles vides, les ventilateurs tournent pendant les récréations, et les téléphones portables restent branchés toute la journée dans les salles de classe. Le proviseur lance donc un concours entre les classes de germanistes : chaque classe doit créer la meilleure affiche en allemand pour sensibiliser les élèves à l'économie d'énergie. L'affiche gagnante sera affichée dans toutes les salles du lycée.

Pour t'inspirer, voici le message que le proviseur a fait passer dans les classes :

\begin{textebox}[Ankündigung des Direktors — L'annonce du proviseur]
\all{Liebe Schülerinnen und Schüler,}

\all{unsere Schule verbraucht zu viel Strom. Das kostet viel Geld, und das ist schlecht für die Umwelt. Deshalb organisieren wir einen Wettbewerb: Macht ein Plakat auf Deutsch! Das Plakat heißt „Energie sparen in unserer Schule“. Zeigt, wie wir Strom sparen können: Licht ausschalten, Handys nicht den ganzen Tag laden, Ventilatoren ausmachen. Das beste Plakat hängt in allen Klassenzimmern. Viel Erfolg!}

\emph{Traduction : Chers élèves, notre école consomme trop d'électricité. Cela coûte beaucoup d'argent et c'est mauvais pour l'environnement. C'est pourquoi nous organisons un concours : faites une affiche en allemand ! L'affiche s'intitule « Économiser l'énergie dans notre école ». Montrez comment nous pouvons économiser l'électricité : éteindre la lumière, ne pas charger les téléphones toute la journée, éteindre les ventilateurs. La meilleure affiche sera accrochée dans toutes les salles de classe. Bonne chance !}
\end{textebox}

\textbf{Glossaire :} \all{der Strom} (l'électricité), \all{verbrauchen} (consommer), \all{der Wettbewerb} (le concours), \all{das Plakat} (l'affiche), \all{ausschalten} (éteindre), \all{laden} (charger), \all{das Klassenzimmer} (la salle de classe), \all{ausmachen} (éteindre, familier).

\courssub{Consignes}

Travaillez par groupes de quatre. Répartissez les rôles : un secrétaire (il écrit), un correcteur de grammaire, un dessinateur, un présentateur. Suivez les étapes dans l'ordre.

\begin{enumerate}
\item \textbf{Brainstorming (5 min).} En groupe, listez en allemand cinq problèmes d'énergie que vous observez dans votre lycée (lumière, ventilateur, téléphone, ordinateur, eau...). Utilisez le dictionnaire si nécessaire.

\item \textbf{Le slogan (5 min).} Trouvez un slogan court et fort, à l'impératif ou avec un modal. Exemples : \all{Schaltet das Licht ab!} ; \all{Wir müssen Energie sparen!} ; \all{Strom sparen — Zukunft sichern!}

\item \textbf{Les cinq conseils (15 min).} Rédigez cinq conseils. Chaque conseil doit utiliser un verbe de modalité différent parmi : \all{müssen}, \all{sollen}, \all{können}, \all{dürfen}, \all{wollen}. Attention : le modal est conjugué en deuxième position, l'infinitif va à la fin de la phrase. Exemple : \all{Wir müssen das Licht ausschalten.}

\item \textbf{La finalité (10 min).} Ajoutez au moins deux constructions de finalité : une avec \all{um ... zu} (même sujet) et une avec \all{damit} (sujets différents). Exemple : \all{Wir schalten das Licht ab, um Strom zu sparen.} / \all{Wir sparen Energie, damit die Schule weniger Geld bezahlt.}

\item \textbf{Le lexique illustré (10 min).} Choisissez cinq mots du vocabulaire de la leçon (par exemple \all{die Steckdose}, \all{das Solarpanel}, \all{das Licht}, \all{der Strom}, \all{die Sonne}) et illustrez-les par des dessins simples avec une étiquette en allemand (avec l'article !).

\item \textbf{La présentation orale (2 min par groupe).} Le présentateur (aidé du groupe) présente l'affiche : il lit le slogan, explique deux conseils et une phrase de finalité. Parlez lentement et clairement.

\item \textbf{Le vote (5 min).} La classe vote pour la meilleure affiche selon trois critères : la correction de la langue, la clarté du message, la beauté du dessin. L'affiche gagnante sera affichée dans la salle.
\end{enumerate}

\courssub{Ressources à mobiliser}

\begin{propriete}[Conjugaison des verbes de modalité au Présent (Präsens)]
Les verbes de modalité sont irréguliers au singulier (modification de voyelle pour \all{können}, \all{müssen}, \all{dürfen}, \all{sollen}, \all{wollen}) et réguliers au pluriel. L'infinitif principal se place \cle{toujours à la fin} de la phrase.
\begin{center}
\begin{tabular}{|l|c|c|c|c|c|}\hline
\rowcolor{popblueL} \textbf{Personne} & \textbf{müssen} & \textbf{können} & \textbf{dürfen} & \textbf{sollen} & \textbf{wollen} \\ \hline
ich & muss & kann & darf & soll & will \\ \hline
du & musst & kannst & darfst & sollst & willst \\ \hline
er/sie/es & muss & kann & darf & soll & will \\ \hline
wir & müssen & können & dürfen & sollen & wollen \\ \hline
ihr & müsst & könnt & dürft & sollt & wollt \\ \hline
sie/Sie & müssen & können & dürfen & sollen & wollen \\ \hline
\end{tabular}
\end{center}
\end{propriete}

\begin{aretenir}[Les verbes de modalité : la structure en « cadre »]
Le modal conjugué occupe la \cle{deuxième position}, l'infinitif se place \cle{à la fin} de la proposition : \all{Wir \textbf{müssen} das Licht \textbf{ausschalten}.} 
Attention aux verbes à particule séparable : à l'infinitif (avec modal), la particule reste soudée : \all{... das Licht ausschalten}. Sans modal (verbe conjugué principal), la particule se sépare et va à la fin : \all{Wir schalten das Licht \textbf{ab}.}
\end{aretenir}

\begin{aretenir}[La finalité : \all{um ... zu} vs \all{damit}]
\begin{itemize}
\item \textbf{\all{um ... zu} + infinitif} : même sujet dans les deux actions. La construction \all{um ... zu} encadre l'infinitif (avec \all{zu} devant l'infinitif à la fin). Exemple : \all{Wir sparen Strom, um die Umwelt zu schützen.} (Nous économisons de l'électricité pour protéger l'environnement.)
\item \textbf{\all{damit} + subordonnée} : sujets différents. \all{damit} introduit une subordonnée : le verbe conjugué se place \cle{à la fin}. Exemple : \all{Wir machen das Licht aus, damit die Schule Geld spart.} (Nous éteignons la lumière pour que l'école économise de l'argent.)
\end{itemize}
\end{aretenir}

\begin{center}
\begin{tabularx}{\linewidth}{|l|X|l|}\hline
\rowcolor{popblueL} \textbf{Mot} & \textbf{Traduction} & \textbf{Pluriel} \\ \hline
\all{die Energie} & l'énergie & \all{die Energien} \\ \hline
\all{der Strom} & l'électricité & --- \\ \hline
\all{die Steckdose} & la prise électrique & \all{die Steckdosen} \\ \hline
\all{das Solarpanel} & le panneau solaire & \all{die Solarpanels} \\ \hline
\all{die Sonne} & le soleil & \all{die Sonnen} \\ \hline
\all{das Licht} & la lumière / le feu & \all{die Lichter} \\ \hline
\all{der Ventilator} & le ventilateur & \all{die Ventilatoren} \\ \hline
\all{sparen} & économiser & --- \\ \hline
\all{ausschalten / ausmachen} & éteindre & --- \\ \hline
\all{verbrauchen} & consommer & --- \\ \hline
\all{laden} & charger (batterie) & --- \\ \hline
\all{verlassen} & quitter (un lieu) & --- \\ \hline
\all{nutzen} & utiliser / exploiter & --- \\ \hline
\end{tabularx}
\end{center}

\courssub{Proposition de production et corrigé commenté}

\begin{exoresolu}[Une affiche modèle]
Rédige l'affiche complète « Energie sparen in unserer Schule » avec : un slogan, cinq conseils (cinq modaux différents), deux constructions de finalité et cinq mots du lexique illustrés.
\tcblower
\textbf{Production modèle :}

\all{\textbf{ENERGIE SPAREN IN UNSERER SCHULE!}}

\all{\textbf{Schaltet das Licht ab — die Sonne ist unser Freund!}}

\all{1. Wir \textbf{müssen} das Licht ausschalten, wenn wir das Klassenzimmer verlassen.}

\all{2. Ihr \textbf{sollt} die Handys nicht den ganzen Tag an der Steckdose laden.}

\all{3. Ihr \textbf{könnt} die Ventilatoren ausmachen, um Strom zu sparen.}

\all{4. Wir \textbf{dürfen} nicht zu viel Energie verbrauchen, damit die Schule weniger Geld bezahlt.}

\all{5. Wir \textbf{wollen} Solarpanels auf das Dach bauen, um die Energie der Sonne zu nutzen.}

\all{[Dessins : eine Steckdose, ein Solarpanel, ein Licht, die Sonne, ein Ventilator]}

\textbf{Commentaire des choix de langue :}
\begin{itemize}
\item \textbf{Le slogan} utilise l'impératif pluriel (\all{Schaltet ... ab!}), forme adaptée à une affiche qui s'adresse à tous les élèves (\all{ihr}-form) ; le verbe à particule \all{abschalten} est correctement séparé (\all{Schaltet ... ab}). \textbf{Attention} : la forme \all{Schalt} est l'impératif singulier (\all{du}), incorrect ici.
\item \textbf{Les cinq modaux} sont tous différents et correctement conjugués : \all{müssen} (1re pers. plur.), \all{sollt} (2e pers. plur. de \all{sollen}), \all{könnt} (2e pers. plur. de \all{können}), \all{dürfen} (1re pers. plur.), \all{wollen} (1re pers. plur.). Chaque infinitif est bien placé en fin de phrase : \all{ausschalten}, \all{laden}, \all{ausmachen}, \all{verbrauchen}, \all{bauen}.
\item \textbf{La finalité} apparaît deux fois : 
      \begin{itemize}
      \item \all{um Strom zu sparen} (même sujet : \all{ihr} $\rightarrow$ infinitif \all{sparen} avec \all{zu} en fin de groupe).
      \item \all{damit die Schule weniger Geld bezahlt} (sujets différents : \all{wir} / \all{die Schule}, verbe conjugué \all{bezahlt} à la fin de la subordonnée).
      \end{itemize}
\item \textbf{Le lexique thématique} est bien présent avec les bons articles : \all{das Licht}, \all{die Steckdose}, \all{der Strom}, \all{die Energie}, \all{die Sonne}, \all{das Solarpanel}, \all{das Klassenzimmer}.
\item \textbf{Pièges évités} : 
      \begin{itemize}
      \item Pas de \all{*wir müssen ausschalten das Licht} (l'infinitif doit fermer la phrase).
      \item Pas de \all{*um zu Strom sparen} (l'infinitif avec \all{zu} se place en fin : \all{um Strom zu sparen}).
      \item Pas de \all{*damit die Schule weniger Geld bezahlen} (verbe conjugué \all{bezahlt} 3e pers. sg. à la fin).
      \item Respect de la casse : majuscule à tous les noms (\all{Licht}, \all{Strom}, \all{Schule}...).
      \end{itemize}
\end{itemize}
\end{exoresolu}

\courssub{Bilan de l'activité}

À l'issue de cette activité, l'élève a mobilisé simultanément les quatre savoir-faire de la leçon : comprendre une consigne authentique (l'annonce du proviseur), réutiliser le lexique de l'énergie avec les bons articles, appliquer les deux points de grammaire du programme (modaux et finalité) et produire un texte court mais réel, destiné à un vrai public. La dimension collective (répartition des rôles, vote final) développe en outre la coopération et la prise de parole en public, compétences transversales essentielles au lycée. Le professeur peut prolonger l'activité en faisant réellement afficher la production gagnante et en demandant à chaque élève de recopier dans son cahier trois conseils de l'affiche de son choix, en soulignant le modal et la construction de finalité.

\newpage

\coursec{Méthodes et exercices résolus}

\courssub{Savoir-faire 1 : Comprendre un texte sur la gestion de l'énergie}

\begin{methode}[Lire et comprendre un texte documentaire en deux lectures]
\begin{enumerate}
\item \textbf{Première lecture (globale).} Lis le titre, la source et le premier paragraphe : ils annoncent presque toujours le thème. Demande-toi : \emph{Worum geht es ?} (De quoi s'agit-il ?) Ici : \all{Energie sparen im Alltag} = économiser l'énergie au quotidien.
\item \textbf{Repérage des mots-clés du champ lexical.} Souligne les mots de la famille de \all{die Energie} : \all{der Strom}, \all{die Steckdose}, \all{sparen}, \all{das Wasserkraftwerk}, \all{die Sonne}. Ils forment un réseau qui t'aide à deviner les mots inconnus.
\item \textbf{Deuxième lecture (détaillée).} Lis phrase par phrase en repérant la structure : sujet -- verbe conjugué en 2\textsuperscript{e} position -- compléments. Ne traduis pas mot à mot : traduis par blocs de sens.
\item \textbf{Exploitation du contexte.} Un mot inconnu ? Devine grâce au contexte : \all{Man muss das Licht ausschalten, wenn man das Zimmer verlässt.} : \all{ausschalten} se comprend grâce à \all{das Licht} (éteindre).
\item \textbf{Réponse aux questions.} Réponds toujours par une phrase complète en allemand, en reprenant les mots de la question, jamais par un simple « oui » ou « non ».
\end{enumerate}
\end{methode}

\begin{exoresolu}[Compréhension de l'écrit : questions sur le texte]
\textbf{Énoncé.} Lis le texte \emph{Energie sparen im Alltag} et réponds en allemand par des phrases complètes :
\begin{enumerate}
\item Warum müssen wir Energie sparen?
\item Was kann man im Haushalt machen, um Strom zu sparen?
\item Welche Energiequellen nennt der Text?
\end{enumerate}
\tcblower
\textbf{Raisonnement.} On repère d'abord les passages du texte qui contiennent la réponse : pour la question 1, la phrase avec \all{weil} ou \all{denn} ; pour la question 2, les exemples concrets (lumière, appareils, prises) ; pour la question 3, la liste des sources d'énergie.

\textbf{Réponses rédigées.}
\begin{enumerate}
\item \all{Wir müssen Energie sparen, weil die Energie teuer ist und weil die Produktion von Strom die Umwelt belastet.} « Nous devons économiser l'énergie parce qu'elle est chère et parce que la production d'électricité pollue l'environnement. » (Justification : subordonnée de cause avec \all{weil}, verbe conjugué à la fin.)
\item \all{Man kann das Licht ausschalten, wenn man das Zimmer verlässt, und man soll die Geräte nicht auf Standby lassen.} « On peut éteindre la lumière en quittant la pièce et on ne doit pas laisser les appareils en veille. »
\item \all{Der Text nennt die Sonne, das Wasser (das Wasserkraftwerk) und den Wind.} « Le texte mentionne le soleil, l'eau (la centrale hydroélectrique) et le vent. »
\end{enumerate}
\textbf{Conclusion.} Chaque réponse reprend les mots de la question, contient un verbe conjugué à la bonne place et reste fidèle au texte : c'est ce que le correcteur attend.
\end{exoresolu}

\courssub{Savoir-faire 2 : Employer le vocabulaire du thème « die Energie »}

\begin{methode}[Mémoriser et réutiliser un champ lexical]
\begin{enumerate}
\item \textbf{Apprends chaque nom avec son article et son pluriel} : \all{die Energie, -n} ; \all{der Strom} (sans pluriel) ; \all{die Steckdose, -n} ; \all{das Wasserkraftwerk, -e} ; \all{die Sonne, -n}.
\item \textbf{Regroupe par familles de mots} : \all{sparen} (économiser) $\rightarrow$ \all{sparsam} (économe), \all{die Ersparnis, -se} (l'économie, l'épargne) ; \all{die Energie} $\rightarrow$ \all{die Energiequelle, -n} (source d'énergie), \all{die Energiewende} (transition énergétique), \all{die Solarenergie} (énergie solaire).
\item \textbf{Apprends les verbes avec leur construction} : \all{Strom sparen} (économiser l'électricité), \all{das Licht ausschalten} (éteindre la lumière, particule séparable !), \all{ein Gerät an die Steckdose anschließen} (brancher un appareil), \all{den Stecker aus der Steckdose ziehen} (débrancher la fiche de la prise).
\item \textbf{Réutilise dans tes propres phrases} : pour chaque mot, écris une phrase personnelle sur ta vie au Cameroun (coupures d'électricité, lampes solaires, etc.).
\item \textbf{Vérifie les majuscules} : tous les noms prennent une majuscule en allemand, même au milieu de la phrase.
\end{enumerate}
\end{methode}

\begin{exoresolu}[Lexique : compléter un texte à trous]
\textbf{Énoncé.} Complète avec : \all{die Steckdose} -- \all{sparen} -- \all{der Strom} -- \all{die Sonne} -- \all{das Wasserkraftwerk}.

\all{In Kamerun produziert ein (1) \dots\ am Fluss Sanaga viel (2) \dots\ . Viele Familien wollen Energie (3) \dots\ , denn sie ist teuer. Auf dem Land benutzt man oft (4) \dots\ für die Lampen. Zu Hause zieht man abends den Stecker aus (5) \dots\ .}
\tcblower
\textbf{Raisonnement.} On analyse chaque trou : article, sens, cas.
\begin{enumerate}
\item \all{ein Wasserkraftwerk} : article indéfini neutre \all{ein} $\rightarrow$ nom neutre \all{das Wasserkraftwerk}. Sens : une centrale au bord d'un fleuve produit de l'électricité.
\item \all{viel Strom} : après \all{viel} + nom masculin sans article au singulier ; \all{Strom produzieren} = produire de l'électricité (collocation à connaître).
\item \all{Energie sparen} : infinitif après \all{wollen} (verbe modal), placé à la fin ; sens « économiser ».
\item \all{die Sonne} : accusatif féminin après \all{benutzen} ; contexte : lampes solaires à la campagne.
\item \all{der Steckdose} : attention au cas ! \all{aus} + datif $\rightarrow$ \all{aus der Steckdose} (on débranche la fiche \emph{de} la prise).
\end{enumerate}
\textbf{Texte complété.} \all{In Kamerun produziert ein Wasserkraftwerk am Fluss Sanaga viel Strom. Viele Familien wollen Energie sparen, denn sie ist teuer. Auf dem Land benutzt man oft die Sonne für die Lampen. Zu Hause zieht man abends den Stecker aus der Steckdose.}
\textbf{Conclusion.} Le choix du mot dépend du sens, mais la \emph{forme} dépend du cas : toujours vérifier la préposition et le verbe qui gouvernent le nom.
\end{exoresolu}

\courssub{Savoir-faire 3 : Employer les verbes de modalité}

\begin{methode}[Construire une phrase avec un verbe modal]
\begin{enumerate}
\item \textbf{Connais les six modaux et leur sens} : \all{können} (pouvoir, capacité), \all{müssen} (devoir, obligation), \all{dürfen} (avoir le droit), \all{sollen} (devoir, recommandation), \all{wollen} (vouloir), \all{mögen/möchte} (aimer, vouloir poliment).
\item \textbf{Conjugue le modal à la 2\textsuperscript{e} position}, l'infinitif de l'action va \textbf{à la fin} de la phrase : \all{Wir müssen Energie sparen.}
\item \textbf{Attention aux singuliers irréguliers} : \all{ich kann, du kannst, er kann} ; \all{ich muss, er muss} ; \all{ich darf, er darf} ; \all{ich soll, er soll} ; \all{ich will, er will} ; \all{ich mag, er mag}. La 1\textsuperscript{re} et la 3\textsuperscript{e} personnes du singulier sont identiques, sans -t.
\item \textbf{Dans une subordonnée}, le modal conjugué va à la toute fin, après l'infinitif : \all{..., weil wir Energie sparen müssen.}
\item \textbf{Choisis le bon modal selon le sens} : obligation $\rightarrow$ \all{müssen} ; conseil $\rightarrow$ \all{sollen} ; possibilité $\rightarrow$ \all{können}.
\end{enumerate}
\end{methode}

\begin{center}
\begin{tabular}{|l|l|l|l|}
\hline
\rowcolor{popblueL} Infinitif & ich / er / sie / es & du & wir / sie / Sie \\
\hline
\all{können} (pouvoir) & \all{kann} & \all{kannst} & \all{können} \\
\hline
\all{müssen} (devoir) & \all{muss} & \all{musst} & \all{müssen} \\
\hline
\all{dürfen} (avoir le droit) & \all{darf} & \all{darfst} & \all{dürfen} \\
\hline
\all{sollen} (devoir, conseil) & \all{soll} & \all{sollst} & \all{sollen} \\
\hline
\all{wollen} (vouloir) & \all{will} & \all{willst} & \all{wollen} \\
\hline
\all{mögen} (aimer) & \all{mag} & \all{magst} & \all{mögen} \\
\hline
\end{tabular}
\end{center}

\begin{exoresolu}[Thème : traduction avec verbes de modalité]
\textbf{Énoncé.} Traduis en allemand :
\begin{enumerate}
\item Nous devons économiser l'énergie.
\item On peut utiliser le soleil pour produire de l'électricité.
\item Tu ne dois pas laisser la lumière allumée. (conseil)
\item Je veux débrancher l'appareil parce que je veux économiser de l'argent.
\end{enumerate}
\tcblower
\textbf{Raisonnement et solutions.}
\begin{enumerate}
\item Obligation $\rightarrow$ \all{müssen}. \all{Wir müssen Energie sparen.} (modal conjugué en 2\textsuperscript{e} position, infinitif \all{sparen} à la fin).
\item Possibilité $\rightarrow$ \all{können}. \all{Man kann die Sonne benutzen, um Strom zu produzieren.} (« produire de l'électricité » = \all{Strom produzieren} ; but exprimé par \all{um ... zu}).
\item Conseil/interdiction $\rightarrow$ \all{sollen} + négation. \all{Du sollst das Licht nicht anlassen.} Attention : \all{anlassen} est un verbe à particule séparable ; avec un modal, le verbe reste entier à l'infinitif final. Ne pas calquer « allumée » par un participe : l'allemand dit simplement \all{das Licht anlassen}.
\item Volonté $\rightarrow$ \all{wollen}. \all{Ich will das Gerät ausstecken, weil ich Geld sparen will.} Dans la subordonnée avec \all{weil}, l'ordre est : sujet -- complément -- infinitif -- modal conjugué à la fin : \all{..., weil ich Geld sparen will.}
\end{enumerate}
\textbf{Conclusion.} Deux réflexes sauvent la phrase : modal conjugué en position 2 (ou en fin de subordonnée) et infinitif renvoyé à la fin de la phrase principale.
\end{exoresolu}

\courssub{Savoir-faire 4 : Exprimer la finalité avec \all{um ... zu} et \all{damit}}

\begin{methode}[Choisir entre \all{um ... zu} et \all{damit}]
\begin{enumerate}
\item \textbf{Même sujet dans les deux parties} $\rightarrow$ \all{um ... zu} + infinitif à la fin : \all{Ich schalte das Licht aus, um Strom zu sparen.} (c'est moi qui éteins et moi qui économise).
\item \textbf{Sujets différents} $\rightarrow$ \all{damit} + subordonnée (verbe conjugué à la fin) : \all{Der Lehrer erklärt die Regel, damit die Schüler sie verstehen.} (le professeur / les élèves).
\item \textbf{Place de l'infinitif} : avec \all{um ... zu}, le \all{zu} se place devant l'infinitif ; pour un verbe à particule séparable, \all{zu} s'insère au milieu : \all{um das Licht auszuschalten}.
\item \textbf{Traduction française} : « pour + infinitif ».
\item \textbf{Vérifie toujours le sujet} avant de choisir : c'est le seul critère.
\end{enumerate}
\end{methode}

\begin{exoresolu}[Transformation : relier deux phrases par la finalité]
\textbf{Énoncé.} Relie les phrases avec \all{um ... zu} ou \all{damit}.
\begin{enumerate}
\item Wir benutzen Solarpanels. Wir produzieren Strom.
\item Die Regierung baut ein Wasserkraftwerk. Die Dörfer bekommen Elektrizität.
\item Er zieht den Stecker aus der Steckdose. Er will Energie sparen.
\end{enumerate}
\tcblower
\textbf{Raisonnement et solutions.}
\begin{enumerate}
\item Même sujet (\all{wir}) $\rightarrow$ \all{um ... zu} : \all{Wir benutzen Solarpanels, um Strom zu produzieren.} « Nous utilisons des panneaux solaires pour produire de l'électricité. »
\item Sujets différents (\all{die Regierung} / \all{die Dörfer}) $\rightarrow$ \all{damit} : \all{Die Regierung baut ein Wasserkraftwerk, damit die Dörfer Elektrizität bekommen.} Verbe conjugué \all{bekommen} à la fin de la subordonnée.
\item Même sujet (\all{er}) $\rightarrow$ \all{um ... zu} : \all{Er zieht den Stecker aus der Steckdose, um Energie zu sparen.} Noter le datif \all{aus der Steckdose}.
\end{enumerate}
\textbf{Conclusion.} Le critère unique est l'identité des sujets : identiques $\rightarrow$ \all{um ... zu} ; différents $\rightarrow$ \all{damit} avec verbe conjugué final.
\end{exoresolu}

\courssub{Savoir-faire 5 : Résumer et produire un court texte guidé}

\begin{methode}[Rédiger un résumé ou un court paragraphe en allemand]
\begin{enumerate}
\item \textbf{Repère les 3 ou 4 idées principales} du texte (une par paragraphe en général) : ici problème (énergie chère, pollution) -- solutions au quotidien -- énergies renouvelables.
\item \textbf{Reformule avec tes propres mots}, sans copier les phrases du texte : utilise des connecteurs simples : \all{zuerst} (d'abord), \all{dann} (ensuite), \all{außerdem} (de plus), \all{deshalb} (c'est pourquoi), \all{zum Schluss} (pour finir).
\item \textbf{Utilise les structures de la leçon} : au moins un verbe modal et une finalité (\all{um ... zu}) : cela montre ta maîtrise de la grammaire.
\item \textbf{Reste bref} : 5 à 8 phrases simples et correctes valent mieux qu'une longue phrase fautive.
\item \textbf{Relis} : verbe en 2\textsuperscript{e} position, majuscules aux noms, infinitif final après un modal.
\end{enumerate}
\end{methode}

\begin{exoresolu}[Production écrite guidée : court paragraphe]
\textbf{Énoncé.} Rédige un court texte (5 à 7 phrases) : « Was kannst du zu Hause und in der Schule machen, um Energie zu sparen? » Utilise au moins deux verbes de modalité et une construction avec \all{um ... zu}.
\tcblower
\textbf{Plan.} 1) Introduction (le problème) ; 2) actions à la maison ; 3) actions à l'école ; 4) conclusion personnelle.

\textbf{Production modèle.}

\all{Energie ist teuer, deshalb müssen wir alle Energie sparen. Zu Hause kann ich das Licht ausschalten, wenn ich das Zimmer verlasse. Ich soll auch die Geräte nicht auf Standby lassen und den Stecker aus der Steckdose ziehen. In der Schule wollen wir die Fenster öffnen, um kein elektrisches Licht zu benutzen. Außerdem können wir Solarenergie benutzen, denn in Kamerun scheint oft die Sonne. Zum Schluss: Jeder muss helfen, um unsere Umwelt zu schützen.}

« L'énergie est chère, c'est pourquoi nous devons tous économiser l'énergie. À la maison, je peux éteindre la lumière quand je quitte la pièce. Je ne dois pas non plus laisser les appareils en veille et je dois débrancher la fiche de la prise. À l'école, nous voulons ouvrir les fenêtres pour ne pas utiliser la lumière électrique. De plus, nous pouvons utiliser l'énergie solaire, car au Cameroun le soleil brille souvent. Pour finir : chacun doit aider pour protéger notre environnement. »

\textbf{Vérification grammaticale.} Modaux utilisés : \all{müssen}, \all{kann}, \all{soll}, \all{wollen}, \all{können}, \all{muss} (infinitifs à la fin : \all{sparen}, \all{ausschalten}, \all{lassen}, \all{ziehen}, \all{öffnen}, \all{benutzen}, \all{helfen}). Finalité : \all{um kein elektrisches Licht zu benutzen}, \all{um unsere Umwelt zu schützen}. Subordonnées correctes : \all{wenn ich das Zimmer verlasse} (verbe conjugué à la fin) ; \all{denn in Kamerun scheint oft die Sonne} (\all{denn} = conjonction de coordination, verbe en position 2). Majuscules à tous les noms : \all{Energie}, \all{Licht}, \all{Zimmer}, \all{Geräte}, \all{Stecker}, \all{Steckdose}, \all{Schule}, \all{Fenster}, \all{Solarenergie}, \all{Sonne}, \all{Schluss}, \all{Umwelt}.
\textbf{Conclusion.} Un texte simple, structuré par des connecteurs et sans faute de place du verbe obtient la totalité des points.
\end{exoresolu}

\begin{attention}[Les 5 erreurs qui coûtent des points au Bac]
\begin{enumerate}
\item \textbf{Oublier la majuscule aux noms.} \all{der strom}, \all{die steckdose} sont fautifs : écris \all{der Strom}, \all{die Steckdose}. En allemand, \emph{tous} les noms prennent une majuscule.
\item \textbf{Mal placer l'infinitif avec un modal.} Ne dis jamais \all{Wir müssen sparen Energie} : l'infinitif va à la fin : \all{Wir müssen Energie sparen.}
\item \textbf{Confondre \all{um ... zu} et \all{damit}.} Avec deux sujets différents, \all{um ... zu} est impossible : \all{Der Staat baut ein Kraftwerk, damit die Menschen Strom haben} (et non \all{um zu haben}).
\item \textbf{Calquer le français.} « Économiser de l'énergie » ne se traduit pas par \all{von Energie sparen} : \all{sparen} prend l'accusatif direct : \all{Energie sparen}. De même, « éteindre la lumière » = \all{das Licht ausschalten} (et non \all{löschen}, qui concerne un feu).
\item \textbf{Faux amis et cas après préposition.} \all{aus} exige le datif : \all{aus der Steckdose} (et non \all{aus die Steckdose}). Attention aussi à \all{die Energie} qui est féminin, contrairement à ce que la terminaison pourrait faire croire.
\end{enumerate}
\end{attention}

\newpage

\coursec{Exercices}

\courssub{Je vérifie mes connaissances}

\begin{aretenir}[Rappel express avant les exercices]
\begin{itemize}
\item Les six verbes de modalité au présent : \all{ich kann / muss / darf / soll / will / mag}. Le modal conjugué est en \cle{2\up{e} position}, l'infinitif à la \cle{fin} de la phrase : \all{Ich muss das Licht ausschalten.} « Je dois éteindre la lumière. »
\item La finalité : \all{um ... zu} + infinitif quand le sujet est le même (\all{Ich lerne Deutsch, um in Deutschland zu arbeiten.}) ; \all{damit} + verbe conjugué à la fin quand les sujets sont différents (\all{Die Schule kauft Lampen, damit die Kinder besser lernen.}).
\end{itemize}
\end{aretenir}

\begin{exercice}[QCM : la gestion de l'énergie]
Entoure la bonne réponse.
\begin{enumerate}
\item \all{die Energie} signifie : a) l'argent \quad b) l'énergie \quad c) l'eau
\item \all{die Steckdose} signifie : a) la prise de courant \quad b) l'ampoule \quad c) le compteur
\item \all{das Wasserkraftwerk} est : a) une éolienne \quad b) un panneau solaire \quad c) une centrale hydroélectrique
\item \all{Strom sparen} signifie : a) gaspiller l'électricité \quad b) économiser l'électricité \quad c) produire du vent
\item \all{die Sonne} signifie : a) la lune \quad b) le soleil \quad c) la pluie
\item Quel verbe modal exprime l'obligation ? a) \all{wollen} \quad b) \all{müssen} \quad c) \all{mögen}
\end{enumerate}
\end{exercice}

\begin{exercice}[Vrai ou faux ?]
Réponds par \all{richtig} ou \all{falsch} et justifie ta réponse par une courte phrase en français.
\begin{enumerate}
\item Le verbe modal conjugué se place à la deuxième position de la phrase et l'infinitif à la fin. \hfill \all{richtig / falsch}
\item On utilise \all{um ... zu} quand les deux propositions ont des sujets différents. \hfill \all{richtig / falsch}
\item \all{Ich will} signifie « je veux ». \hfill \all{richtig / falsch}
\item \all{Er kann das Licht ausschalten} signifie « Il doit éteindre la lumière ». \hfill \all{richtig / falsch}
\item Après \all{damit}, le verbe conjugué se place à la fin de la subordonnée. \hfill \all{richtig / falsch}
\end{enumerate}
\end{exercice}

\begin{exercice}[Vocabulaire : articles et pluriels]
Complète le tableau avec l'article et le pluriel de chaque nom, puis traduis-le en français.

\begin{center}
\begin{tabular}{|l|l|l|}\hline
\rowcolor{popblueL} Nom (Singular) & Plural & Français \\ \hline
\ldots{} Energie & \ldots & \ldots \\ \hline
\ldots{} Strom & \ldots & \ldots \\ \hline
\ldots{} Sonne & \ldots & \ldots \\ \hline
\ldots{} Wasserkraftwerk & \ldots & \ldots \\ \hline
\ldots{} Steckdose & \ldots & \ldots \\ \hline
\ldots{} Lampe & \ldots & \ldots \\ \hline
\end{tabular}
\end{center}
\end{exercice}

\begin{exercice}[Conjugaison : les modaux au présent]
Complète chaque phrase avec le verbe modal entre parenthèses, correctement conjugué au présent.
\begin{enumerate}
\item Ich \ldots{} das Licht ausschalten. (müssen)
\item \ldots{} du mir bitte helfen? (können)
\item Wir \ldots{} Energie sparen. (wollen)
\item Ihr \ldots{} hier nicht rauchen. (dürfen)
\item Er \ldots{} mehr Sport machen, sagt der Arzt. (sollen)
\item Sie (pluriel) \ldots{} gern Tee. (mögen)
\end{enumerate}
\end{exercice}

\courssub{Je m'entraîne}

\begin{exercice}[Le bon modal]
Complète les dix phrases avec le modal qui convient au présent : \all{können, müssen, dürfen, sollen, wollen, mögen} (certains modaux peuvent servir deux fois).
\begin{enumerate}
\item Man \ldots{} die Geräte ganz ausschalten, das spart Strom. (conseil)
\item Kinder \ldots{} nicht an der Steckdose spielen. (interdiction)
\item Wir \ldots{} ein Wasserkraftwerk besuchen. (désir)
\item \ldots{} du schwimmen? (capacité)
\item Ich \ldots{} jetzt leider gehen, der Bus kommt. (nécessité)
\item Mein Bruder \ldots{} keine Suppe. (goût)
\item Ihr \ldots{} mehr Wasser trinken, sagt der Lehrer. (recommandation)
\item \ldots{} ich das Fenster öffnen? (permission)
\item Sie \ldots{} Deutsch lernen, um in Deutschland zu studieren. (volonté)
\item Man \ldots{} hier nicht parken. (interdiction)
\end{enumerate}
\end{exercice}

\begin{exercice}[Réécriture avec un modal]
Transforme chaque phrase en y intégrant le modal entre parenthèses. Attention à la place de l'infinitif !

Exemple : \all{Er schaltet das Licht ab. (müssen)} $\rightarrow$ \all{Er muss das Licht abschalten.}
\begin{enumerate}
\item Wir sparen Energie. (sollen)
\item Sie kauft eine neue Lampe. (wollen)
\item Ich schalte den Computer aus. (können)
\item Die Schüler lernen die Wörter. (müssen)
\item Ihr benutzt die Sonnenenergie. (können)
\item Er repariert die Steckdose. (dürfen)
\end{enumerate}
\end{exercice}

\begin{exercice}[Texte à trous]
Complète le texte avec les mots de la liste : \all{Strom -- sparen -- Sonne -- Steckdose -- Energie -- Wasserkraftwerk}.

\begin{textebox}[Energie in Kamerun]
In Kamerun ist \ldots{} sehr wichtig. Viele Familien wollen \ldots{} sparen, denn er ist teuer. Die \ldots{} scheint fast jeden Tag: Man kann Sonnenenergie nutzen. Am Fluss Sanaga gibt es ein großes \ldots{} . Dort macht man Strom aus Wasser. Zu Hause muss man die Geräte aus der \ldots{} ziehen, wenn man sie nicht braucht. So kann man viel Geld \ldots{} .
\end{textebox}
\end{exercice}

\begin{exercice}[Appariement : la finalité]
Relie chaque début de phrase (1--6) à la bonne fin (a--f), puis indique si la construction est \all{um ... zu} ou \all{damit}.

\begin{center}
\begin{tabularx}{\linewidth}{|X|X|}\hline
\rowcolor{popblueL} Début & Fin \\ \hline
1. Ich schalte das Licht aus, & a. damit die Kinder besser lernen können. \\ \hline
2. Wir lernen Deutsch, & b. um Strom zu sparen. \\ \hline
3. Die Schule kauft neue Lampen, & c. um in Deutschland zu arbeiten. \\ \hline
4. Er geht zum Markt, & d. damit die Rechnung nicht so hoch ist. \\ \hline
5. Die Familie verbraucht wenig Strom, & e. um frisches Gemüse zu kaufen. \\ \hline
6. Man baut ein Wasserkraftwerk, & f. damit das Dorf Elektrizität bekommt. \\ \hline
\end{tabularx}
\end{center}
\end{exercice}

\courssub{Je me prépare au Bac}

\begin{exercice}[Compréhension écrite : les éoliennes en Allemagne]
Lis le texte suivant, puis réponds aux questions.

\begin{textebox}[Windenergie in Deutschland]
In Deutschland sieht man überall Windräder: im Norden, an der Küste, aber auch auf den Feldern. Ein Windrad macht Strom aus Wind — das ist saubere Energie. Viele Deutsche finden die Windenergie gut, denn sie schadet der Umwelt nicht. Aber es gibt auch Probleme: Manche Leute sagen, die Windräder sind zu laut und zu hässlich. Außerdem braucht man viel Platz. Trotzdem will die Regierung mehr Windräder bauen, denn Deutschland muss weniger Kohle und Öl verbrauchen. Sie baut neue Windparks, damit das Land mehr sauberen Strom bekommt. Die Bürger können Geld sparen, wenn sie sauberen Strom nutzen. Experten sagen: Wir sollen die Natur schützen, und wir dürfen die Energie der Zukunft nicht vergessen. Deshalb lernen schon die Schüler in der Schule, wie man Energie spart.
\end{textebox}

\textbf{Questions de compréhension} (réponds en allemand, par des phrases complètes) :
\begin{enumerate}
\item Wo sieht man Windräder in Deutschland? (2 éléments)
\item Warum finden viele Deutsche die Windenergie gut?
\item Nenne zwei Probleme der Windräder.
\item Was will die Regierung bauen? Warum?
\item Was können die Bürger sparen?
\end{enumerate}

\textbf{Questions de langue} :
\begin{enumerate}
\item[6.] Relève dans le texte trois verbes de modalité conjugués et donne leur infinitif.
\item[7.] Trouve dans le texte une phrase exprimant la finalité et recopie-la.
\item[8.] Donne le contraire de : \all{sauber} et \all{mehr}.
\end{enumerate}
\end{exercice}

\begin{exercice}[Thème et version]
\textbf{A. Thème.} Traduis en allemand les phrases suivantes. Attention aux modaux et à la finalité.
\begin{enumerate}
\item Nous devons éteindre la lumière pour économiser l'électricité.
\item Tu peux utiliser l'énergie du soleil.
\item Je veux apprendre l'allemand pour travailler en Allemagne.
\item Les élèves doivent fermer la porte pour que la classe reste fraîche.
\item On ne doit pas gaspiller l'eau.
\end{enumerate}

\textbf{B. Version.} Traduis en français le paragraphe suivant.

\begin{textebox}[Zu Hause]
Meine Familie will weniger Strom verbrauchen. Wir müssen die Lampen ausschalten, wenn wir das Zimmer verlassen. Mein Vater hat neue Lampen gekauft, damit wir weniger Energie brauchen. Meine Mutter sagt, wir sollen auch das Handy nicht die ganze Nacht an der Steckdose lassen. So können wir jeden Monat Geld sparen.
\end{textebox}
\end{exercice}

\begin{exercice}[Production écrite : Energie sparen in Kamerun]
Sujet : \all{Wie kann man in Kamerun Energie sparen?}

Rédige un court texte de \textbf{60 à 80 mots} en allemand. Tu dois :
\begin{itemize}
\item utiliser \textbf{au moins 3 verbes de modalité} différents (\all{können, müssen, sollen, dürfen, wollen}) ;
\item utiliser \textbf{au moins 2 expressions de finalité} (\all{um ... zu} et/ou \all{damit}) ;
\item employer le vocabulaire du thème (\all{die Energie, der Strom, die Sonne, sparen, die Steckdose, ausschalten}...) ;
\item donner des exemples concrets du contexte camerounais (maison, école, marché...).
\end{itemize}
\end{exercice}

\begin{methode}[Réussir la production écrite]
\begin{enumerate}
\item \textbf{Prépare un plan} en 3 parties : introduction (le problème : l'électricité est chère), 2 ou 3 conseils concrets, conclusion (le résultat : économiser de l'argent, protéger la nature).
\item \textbf{Place tes modaux} : modal conjugué en 2\up{e} position, infinitif à la fin. Vérifie chaque phrase à voix haute.
\item \textbf{Varie la finalité} : une phrase avec \all{um ... zu} (même sujet), une avec \all{damit} (sujets différents).
\item \textbf{Relis-toi} : majuscule à tous les noms, Umlaute, virgule avant \all{um} et \all{damit}.
\end{enumerate}
\textbf{Modèle de production} (environ 75 mots) :

\all{In Kamerun ist Strom oft teuer. Deshalb müssen wir Energie sparen. Zu Hause sollen wir das Licht ausschalten, wenn wir das Zimmer verlassen. Man kann auch die Geräte aus der Steckdose ziehen, um Strom zu sparen. Die Sonne scheint jeden Tag: Wir können Sonnenenergie nutzen, damit die Rechnung nicht so hoch ist. In der Schule wollen die Schüler lernen, wie man Energie spart. So spart jede Familie Geld und schützt die Natur.}

« Au Cameroun, l'électricité est souvent chère. C'est pourquoi nous devons économiser l'énergie. À la maison, nous devrions éteindre la lumière quand nous quittons la pièce. On peut aussi débrancher les appareils de la prise pour économiser l'électricité. Le soleil brille tous les jours : nous pouvons utiliser l'énergie solaire pour que la facture ne soit pas si élevée. À l'école, les élèves veulent apprendre comment économiser l'énergie. Ainsi, chaque famille économise de l'argent et protège la nature. »
\end{methode}

\newpage

\coursec{Corrigés des exercices}

\begin{corrige}[Exercice 1 — QCM : la gestion de l'énergie]
1. b) \quad 2. a) \quad 3. c) \quad 4. b) \quad 5. b) \quad 6. b) — \all{müssen} exprime l'obligation ou la nécessité ; \all{wollen} exprime la volonté et \all{mögen} le goût.
\end{corrige}

\begin{corrige}[Exercice 2 — Vrai ou faux ?]
\begin{enumerate}
\item \all{richtig} — c'est la structure en « cadre » : \all{Ich muss Strom sparen.}
\item \all{falsch} — \all{um ... zu} s'emploie quand le sujet est le \emph{même} ; avec des sujets différents, on utilise \all{damit}.
\item \all{richtig} — \all{wollen} exprime la volonté ; attention, \all{ich will} ne signifie jamais « je voudrais » (forme polie : \all{ich möchte}).
\item \all{falsch} — \all{können} signifie « pouvoir, être capable de » ; la phrase signifie « Il peut éteindre la lumière ». « Il doit » serait \all{Er muss ...}.
\item \all{richtig} — \all{damit} est une conjonction de subordination : le verbe conjugué va à la fin.
\end{enumerate}
\end{corrige}

\begin{corrige}[Exercice 3 — Vocabulaire : articles et pluriels]
\begin{center}
\begin{tabularx}{\linewidth}{|X|X|X|}\hline
\rowcolor{popblueL} Nom (Singulier) & Pluriel & Français \\ \hline
die Energie & die Energien & l'énergie \\ \hline
der Strom & (rare : die Ströme) & l'électricité, le courant \\ \hline
die Sonne & die Sonnen & le soleil \\ \hline
das Wasserkraftwerk & die Wasserkraftwerke & la centrale hydroélectrique \\ \hline
die Steckdose & die Steckdosen & la prise de courant \\ \hline
die Lampe & die Lampen & la lampe \\ \hline
\end{tabularx}
\end{center}
Remarque : \all{der Strom} s'emploie presque toujours au singulier quand il signifie « l'électricité » ; au pluriel, \all{die Ströme} désigne les « courants » (électriques, marins) ou les « grands fleuves ».
\end{corrige}

\begin{corrige}[Exercice 4 — Conjugaison : les modaux au présent]
\begin{enumerate}
\item \all{Ich muss das Licht ausschalten.} « Je dois éteindre la lumière. »
\item \all{Kannst du mir bitte helfen ?} « Peux-tu m'aider, s'il te plaît ? »
\item \all{Wir wollen Energie sparen.} « Nous voulons économiser l'énergie. »
\item \all{Ihr dürft hier nicht rauchen.} « Vous n'avez pas le droit de fumer ici. »
\item \all{Er soll mehr Sport machen, sagt der Arzt.} « Il devrait faire plus de sport, dit le médecin. »
\item \all{Sie mögen gern Tee.} « Elles/Ils aiment le thé. » (\all{Sie} = ils/elles ou Vous formel)
\end{enumerate}
\end{corrige}

\begin{corrige}[Exercice 5 — Le bon modal]
\begin{enumerate}
\item \all{soll} — \all{Man soll die Geräte ganz ausschalten, das spart Strom.}
\item \all{dürfen} — \all{Kinder dürfen nicht an der Steckdose spielen.}
\item \all{wollen} — \all{Wir wollen ein Wasserkraftwerk besuchen.}
\item \all{Kannst} — \all{Kannst du schwimmen?}
\item \all{muss} — \all{Ich muss jetzt leider gehen, der Bus kommt.}
\item \all{mag} — \all{Mein Bruder mag keine Suppe.}
\item \all{sollt} — \all{Ihr sollt mehr Wasser trinken, sagt der Lehrer.}
\item \all{Darf} — \all{Darf ich das Fenster öffnen?}
\item \all{will} — \all{Sie will Deutsch lernen, um in Deutschland zu studieren.}
\item \all{darf} — \all{Man darf hier nicht parken.}
\end{enumerate}
\end{corrige}

\begin{corrige}[Exercice 6 — Réécriture avec un modal]
\begin{enumerate}
\item \all{Wir sollen Energie sparen.} « Nous devrions économiser l'énergie. »
\item \all{Sie will eine neue Lampe kaufen.} « Elle veut acheter une nouvelle lampe. »
\item \all{Ich kann den Computer ausschalten.} « Je peux éteindre l'ordinateur. » (verbe à particule séparable : \all{ausschalten} reste soudé à l'infinitif, à la fin de la phrase).
\item \all{Die Schüler müssen die Wörter lernen.} « Les élèves doivent apprendre les mots. »
\item \all{Ihr könnt die Sonnenenergie benutzen.} « Vous pouvez utiliser l'énergie solaire. »
\item \all{Er darf die Steckdose reparieren.} « Il a le droit de réparer la prise. »
\end{enumerate}
\end{corrige}

\begin{corrige}[Exercice 7 — Texte à trous]
\all{In Kamerun ist \textbf{Energie} sehr wichtig. Viele Familien wollen \textbf{Strom} sparen, denn er ist teuer. Die \textbf{Sonne} scheint fast jeden Tag: Man kann Sonnenenergie nutzen. Am Fluss Sanaga gibt es ein großes \textbf{Wasserkraftwerk}. Dort macht man Strom aus Wasser. Zu Hause muss man die Geräte aus der \textbf{Steckdose} ziehen, wenn man sie nicht braucht. So kann man viel Geld \textbf{sparen}.}

Traduction : « Au Cameroun, l'énergie est très importante. Beaucoup de familles veulent économiser l'électricité, car elle est chère. Le soleil brille presque tous les jours : on peut utiliser l'énergie solaire. Sur le fleuve Sanaga, il y a une grande centrale hydroélectrique. Là, on produit de l'électricité à partir de l'eau. À la maison, il faut débrancher les appareils de la prise quand on n'en a pas besoin. Ainsi, on peut économiser beaucoup d'argent. »
\end{corrige}

\begin{corrige}[Exercice 8 — Appariement : la finalité]
\begin{enumerate}
\item 1 -- b : \all{Ich schalte das Licht aus, um Strom zu sparen.} (\all{um ... zu}, même sujet : \all{ich})
\item 2 -- c : \all{Wir lernen Deutsch, um in Deutschland zu arbeiten.} (\all{um ... zu}, même sujet : \all{wir})
\item 3 -- a : \all{Die Schule kauft neue Lampen, damit die Kinder besser lernen können.} (\all{damit}, sujets différents : \all{die Schule / die Kinder})
\item 4 -- e : \all{Er geht zum Markt, um frisches Gemüse zu kaufen.} (\all{um ... zu}, même sujet : \all{er})
\item 5 -- d : \all{Die Familie verbraucht wenig Strom, damit die Rechnung nicht so hoch ist.} (\all{damit}, sujets différents)
\item 6 -- f : \all{Man baut ein Wasserkraftwerk, damit das Dorf Elektrizität bekommt.} (\all{damit}, sujets différents)
\end{enumerate}
Astuce : si tu peux remplacer la fin par « pour + infinitif » avec le même sujet, choisis \all{um ... zu} ; si la fin introduit un nouveau sujet (« pour que »), choisis \all{damit}.
\end{corrige}

\begin{corrige}[Exercice 9 — Compréhension écrite : les éoliennes en Allemagne]
\textbf{Compréhension} (réponses modèles) :
\begin{enumerate}
\item \all{Man sieht Windräder im Norden, an der Küste und auf den Feldern.} « On voit des éoliennes dans le Nord, sur la côte et dans les champs. »
\item \all{Viele Deutsche finden die Windenergie gut, denn sie schadet der Umwelt nicht.} « ... parce qu'elle ne nuit pas à l'environnement. »
\item \all{Die Windräder sind zu laut und zu hässlich, und man braucht viel Platz.} « Elles sont trop bruyantes et trop laides, et il faut beaucoup de place. »
\item \all{Die Regierung will mehr Windräder bauen, denn Deutschland muss weniger Kohle und Öl verbrauchen.} « ... parce que l'Allemagne doit consommer moins de charbon et de pétrole. »
\item \all{Die Bürger können Geld sparen.} « Les citoyens peuvent économiser de l'argent. »
\end{enumerate}
\textbf{Langue} :
\begin{enumerate}
\item[6.] Par exemple : \all{muss} (\all{müssen}), \all{können} (\all{können}), \all{sollen} (\all{sollen}) ; aussi possibles : \all{will} (\all{wollen}), \all{dürfen} (\all{dürfen}).
\item[7.] \all{Sie baut neue Windparks, damit das Land mehr sauberen Strom bekommt.} (finalité avec \all{damit} + sujet différent, verbe à la fin)
\item[8.] \all{sauber} (propre) $\rightarrow$ \all{schmutzig} (sale) ; \all{mehr} (plus) $\rightarrow$ \all{weniger} (moins).
\end{enumerate}
\end{corrige}

\begin{corrige}[Exercice 10 — Thème et version]
\textbf{A. Thème} (corrigé) :
\begin{enumerate}
\item \all{Wir müssen das Licht ausschalten, um Strom zu sparen.} (même sujet $\rightarrow$ \all{um ... zu} ; \all{ausschalten} : particule \all{aus-} soudée à l'infinitif car modal + infinitif)
\item \all{Du kannst die Sonnenenergie benutzen.} (ou : \all{... die Energie der Sonne nutzen.})
\item \all{Ich will Deutsch lernen, um in Deutschland zu arbeiten.}
\item \all{Die Schüler müssen die Tür schließen, damit die Klasse kühl bleibt.} (« pour que » + nouveau sujet $\rightarrow$ \all{damit}, verbe conjugué à la fin)
\item \all{Man darf das Wasser nicht verschwenden.} (ou : \all{Man soll das Wasser nicht verschwenden.})
\end{enumerate}
\textbf{B. Version} (corrigé) : « Ma famille veut consommer moins d'électricité. Nous devons éteindre les lampes quand nous quittons la pièce. Mon père a acheté de nouvelles lampes pour que nous ayons besoin de moins d'énergie. Ma mère dit que nous ne devrions pas non plus laisser le portable branché à la prise toute la nuit. Ainsi, nous pouvons économiser de l'argent chaque mois. »
\end{corrige}

\begin{corrige}[Exercice 11 — Production écrite : Energie sparen in Kamerun]
\textbf{Barème indicatif (sur 10 points)} :
\begin{itemize}
\item Respect de la consigne (longueur, sujet) : 2 pts
\item Vocabulaire thématique correct et varié : 2 pts
\item Au moins 3 modaux différents correctement conjugués et placés : 2 pts
\item Au moins 2 expressions de finalité (\all{um ... zu} / \all{damit}) correctes : 2 pts
\item Orthographe (majuscules noms, Umlaute, \all{ß}), ponctuation (virgule avant \all{um/damit}) : 2 pts
\end{itemize}

\textbf{Proposition de rédaction modèle} (environ 75 mots) :

\all{In Kamerun ist Strom oft teuer. Deshalb müssen wir Energie sparen. Zu Hause sollen wir das Licht ausschalten, wenn wir das Zimmer verlassen. Man kann auch die Geräte aus der Steckdose ziehen, um Strom zu sparen. Die Sonne scheint jeden Tag: Wir können Sonnenenergie nutzen, damit die Rechnung nicht so hoch ist. In der Schule wollen die Schüler lernen, wie man Energie spart. So spart jede Familie Geld und schützt die Natur.}

\textbf{Points de vigilance} :
\begin{itemize}
\item Vérifier l'ordre des mots : modal en position 2, infinitif à la fin (\all{müssen wir ... sparen}, \all{sollen wir ... ausschalten}).
\item Ne pas oublier la virgule avant \all{um} et \all{damit}.
\item \all{Strom} (masculin), \all{die Energie} (féminin), \all{die Steckdose} (féminin) — accords corrects.
\item \all{ausschalten}, \all{ziehen} (verbes à particule séparable) : à l'infinitif après un modal, \all{ausschalten} reste soudé ; \all{aus der Steckdose ziehen} $\rightarrow$ \all{ziehen} à la fin, \all{aus} reste préposition (datif : \all{aus der}).
\end{itemize}
\end{corrige}

\newpage

\coursec{Fiche bilan}

\begin{popfigure}
\begin{tikzpicture}[
  mindmap,
  grow cyclic,
  every node/.style={concept, execute at begin node=\hskip0pt},
  concept color=popblue,
  level 1/.append style={level distance=4.5cm, sibling angle=60},
  level 2/.append style={level distance=3cm, sibling angle=45},
  text=white,
  font=\small\bfseries
]
\node[concept, scale=1.2, align=center]{Energie\\sparen}
  child[concept color=poporange] {node[concept, align=center]{Text-\\verständnis}
    child[concept color=poporange!70!white] {node[concept, scale=0.7] {Global}}
    child[concept color=poporange!70!white] {node[concept, scale=0.7] {Detailliert}}
  }
  child[concept color=popgreen] {node[concept] {Wortschatz}
    child[concept color=popgreen!70!white] {node[concept, scale=0.7, align=center]{Energie-\\arten}}
    child[concept color=popgreen!70!white] {node[concept, scale=0.7, align=center]{Verben\\/Nomen}}
    child[concept color=popgreen!70!white] {node[concept, scale=0.7] {Alltag}}
  }
  child[concept color=poppurple] {node[concept] {Grammatik 1}
    child[concept color=poppurple!70!white] {node[concept, scale=0.7, align=center]{Modal-\\verben}}
    child[concept color=poppurple!70!white] {node[concept, scale=0.7, align=center]{können,\\müssen, ...}}
    child[concept color=poppurple!70!white] {node[concept, scale=0.7] {Stellung}}
  }
  child[concept color=poppink] {node[concept] {Grammatik 2}
    child[concept color=poppink!70!white] {node[concept, scale=0.7] {Finalität}}
    child[concept color=poppink!70!white] {node[concept, scale=0.7] {um ... zu}}
    child[concept color=poppink!70!white] {node[concept, scale=0.7] {damit}}
  }
  child[concept color=popgold] {node[concept] {Ausdruck}
    child[concept color=popgold!70!white] {node[concept, scale=0.7] {Fragen}}
    child[concept color=popgold!70!white] {node[concept, scale=0.7, align=center]{Zusammen-\\fassung}}
    child[concept color=popgold!70!white] {node[concept, scale=0.7] {Meinung}}
  }
  child[concept color=popblue!70!white] {node[concept] {Kontext}
    child[concept color=popblue!50!white] {node[concept, scale=0.7] {Energiewende}}
    child[concept color=popblue!50!white] {node[concept, scale=0.7] {Kamerun}}
  };
\end{tikzpicture}
\legende{Carte mentale récapitulative de la leçon ALA11 : la gestion de l'énergie}
\end{popfigure}

\begin{bilan}

\end{bilan}

\courssub{Lexique clé : die Energie und der Alltag}

\begin{center}
\begin{tabularx}{\linewidth}{|X|X|X|}\hline
\rowcolor{popblueL}
\textbf{Deutsch (Artikel, Plural)} & \textbf{Français} & \textbf{Exemple / Famille} \\ \hline
\all{die Energie, die Energien} & l'énergie & \all{erneuerbare Energien} (énergies renouvelables) \\ \hline
\all{der Strom} (sans pluriel) & l'électricité (courant) & \all{Strom sparen, der Stromverbrauch} \\ \hline
\all{die Sonne, die Sonnen} & le soleil & \all{die Sonnenenergie, die Solaranlage} \\ \hline
\all{das Wasserkraftwerk, die -kraftwerke} & la centrale hydroélectrique & \all{die Wasserkraft} \\ \hline
\all{das Kraftwerk, die -werke} & la centrale (thermique, nucléaire) & \all{das Atomkraftwerk (AKW)} \\ \hline
\all{sparen} & économiser & \all{energiesparend, die Einsparung} \\ \hline
\all{verbrauchen} & consommer & \all{der Verbrauch, der Verbraucher} \\ \hline
\all{die Steckdose, die -n} & la prise électrique & \all{den Stecker aus der Steckdose ziehen} \\ \hline
\all{der Stecker, die -} & la fiche (prise mâle) & \all{den Stecker ziehen} (débrancher) \\ \hline
\all{die Steckdosenleiste, die -n} & la multiprise & \all{eine Steckdosenleiste mit Schalter} \\ \hline
\all{ausschalten / einschalten} & éteindre / allumer & \all{das Gerät ausschalten} (séparable) \\ \hline
\all{der Standby-Modus} & le mode veille & \all{den Standby-Modus vermeiden} \\ \hline
\all{die Energiewende} & la transition énergétique & \all{die Energiewende schaffen} \\ \hline
\all{erneuerbar / fossil} & renouvelable / fossile & \all{fossile Brennstoffe} \\ \hline
\all{isolieren} & isoler & \all{die Isolierung, das Haus isolieren} \\ \hline
\all{die Heizung, die -en} & le chauffage & \all{die Heizung runterdrehen} \\ \hline
\all{die Umwelt} & l'environnement & \all{die Umwelt schützen, der Umweltschutz} \\ \hline
\end{tabularx}
\end{center}

\begin{attention}[Faux amis et pièges de vocabulaire]
\begin{itemize}
\item \all{die Energie} = l'énergie (physique, électrique) ; attention, « énergie » au sens de force personnelle se dit plutôt \all{die Kraft} ou \all{die Energie} selon le contexte.
\item \all{der Strom} désigne le courant électrique au quotidien (\all{Strom sparen} = économiser l'électricité) ; un « courant » d'eau ou d'air se dit aussi \all{der Strom} (\all{der Flussstrom}), mais un fleuve = \all{der Fluss}.
\item \all{sparen} = économiser (de l'argent, de l'énergie) ; « épargner quelqu'un » = \all{jemanden verschonen}.
\item \all{das Gerät, die Geräte} = l'appareil (électrique) ; ne pas confondre avec \all{das Gericht} (le plat, le tribunal).
\end{itemize}
\end{attention}

\courssub{Grammaire 1 : Les verbes de modalité (Modalverben)}

\begin{propriete}[Conjugaison au Présent — \cle{Modalverben}]
Les six verbes modaux modifient le sens du verbe principal, qui reste à l'\textbf{infinitif à la fin de la phrase}. Leur conjugaison est \textbf{irrégulière} : la voyelle change souvent au singulier, et les 1\iere{} et 3\ieme{} personnes du singulier sont \textbf{identiques} (pas de \textit{-t} à la 3\ieme{} personne).
\begin{center}
\begin{tabular}{|l|c|c|c|c|c|c|}\hline
\rowcolor{poppurpleL}
\textbf{Personne} & \textbf{können} & \textbf{müssen} & \textbf{dürfen} & \textbf{sollen} & \textbf{wollen} & \textbf{möchten} \\ \hline
ich & kann & muss & darf & soll & will & möchte \\ \hline
du & kannst & musst & darfst & sollst & willst & möchtest \\ \hline
er/sie/es & \textbf{kann} & \textbf{muss} & \textbf{darf} & \textbf{soll} & \textbf{will} & möchte \\ \hline
wir & können & müssen & dürfen & sollen & wollen & möchten \\ \hline
ihr & könnt & müsst & dürft & sollt & wollt & möchtet \\ \hline
sie/Sie & können & müssen & dürfen & sollen & wollen & möchten \\ \hline
\end{tabular}
\end{center}
\emph{Remarques :} \all{möchten} est le Konjunktiv II de \all{mögen} et sert de forme polie ; il se conjugue régulièrement. Au Perfekt, les modaux ont en théorie un participe (gekonnt, gemusst, gedurft, gesollt, gewollt), mais on utilise presque toujours l'\textbf{infinitif de substitution} (Ersatzinfinitiv) quand un verbe principal est présent : \all{Ich habe Energie sparen müssen.}
\end{propriete}

\begin{propriete}[Sens et emploi (niveau A2)]
\begin{itemize}
\item \all{können} : capacité, possibilité (« savoir faire », « pouvoir »). \all{Wir können viel Strom sparen.} (Nous pouvons économiser beaucoup d'électricité.)
\item \all{müssen} : obligation, nécessité (« devoir », « il faut »). \all{Wir müssen die Umwelt schützen.} (Nous devons protéger l'environnement.)
\item \all{dürfen} : permission, autorisation (« avoir le droit »). \all{Hier darf man keinen Müll wegwerfen.} (Ici, on n'a pas le droit de jeter des déchets.)
\item \all{sollen} : conseil, ordre transmis, devoir moral (« être censé »). \all{Du sollst das Licht ausschalten.} (Tu es censé éteindre la lumière — on te l'a dit.)
\item \all{wollen} : volonté, intention ferme (« vouloir »). \all{Ich will weniger Energie verbrauchen.} (Je veux consommer moins d'énergie.)
\item \all{möchten} : souhait poli (« voudrais », « aimerais »). \all{Ich möchte ein Solarpanel kaufen.} (Je voudrais acheter un panneau solaire.)
\end{itemize}
\end{propriete}

\begin{attention}[Pièges francophones — Modalverben]
\begin{itemize}
\item \textbf{Ordre des mots :} le modal conjugué occupe la 2\ieme{} position (phrase principale), l'infinitif du verbe principal va \textbf{à la toute fin}. \all{Ich muss das Licht ausschalten.} (Jamais *\all{Ich muss ausschalten das Licht}.) Avec un verbe séparable, la particule reste collée à l'infinitif : \all{ausschalten}.
\item \textbf{Perfekt :} auxiliaire \all{haben} + \textbf{deux infinitifs} à la fin (Ersatzinfinitiv). \all{Ich habe das Licht ausschalten müssen.} (J'ai dû éteindre la lumière.)
\item \textbf{\all{müssen} ou \all{sollen} ?} \all{Ich muss lernen.} (contrainte réelle : examen demain). \all{Ich soll lernen.} (conseil ou ordre de quelqu'un : mes parents, le médecin).
\item \textbf{\all{müssen nicht} ou \all{nicht dürfen} ?} \all{Ich muss nicht kommen.} = je ne suis pas obligé de venir (j'ai le choix). \all{Ich darf nicht kommen.} = il m'est interdit de venir. Le français « ne pas devoir » se traduit par \all{nicht müssen}, jamais par \all{nicht dürfen}.
\item \textbf{\all{mögen} ou \all{möchten} ?} \all{Ich mag Schokolade.} (j'aime le chocolat). \all{Ich möchte ein Glas Wasser.} (je voudrais un verre d'eau).
\item \textbf{Pas de \all{zu} :} contrairement au français « pouvoir faire », le modal allemand s'emploie \textbf{sans} \all{zu} : \all{Ich kann schwimmen.} (Je sais nager.)
\end{itemize}
\end{attention}

\courssub{Grammaire 2 : La finalité (Finalsatz)}

\begin{propriete}[Deux constructions pour exprimer le but]
\begin{center}
\begin{tabularx}{\linewidth}{|X|X|}\hline
\rowcolor{poppinkL}
\textbf{\all{um ... zu + Infinitif}} & \textbf{\all{damit + subordonnée}} \\ \hline
\emph{Sujet identique} dans les deux parties (obligatoire). & \emph{Sujets différents} possibles (ou même sujet, pour insister). \\ \hline
Verbe à l'infinitif avec \all{zu}, placé \textbf{à la fin}. & Verbe conjugué \textbf{à la fin} (structure de subordonnée). \\ \hline
\all{Ich spare Energie, um die Umwelt zu schützen.} & \all{Ich schalte das Licht aus, damit wir Strom sparen.} \\ \hline
\emph{Négation :} \all{um nicht zu ...} / \all{um kein ... zu ...} & \emph{Négation :} \all{damit ... nicht} / \all{damit kein ...} \\ \hline
\end{tabularx}
\end{center}
\emph{Règle de choix :} même sujet $\rightarrow$ \all{um ... zu} ; sujets différents $\rightarrow$ \all{damit}. En cas de doute, \all{damit} est toujours correct.
\end{propriete}

\begin{exemplebox}[Exemples contextualisés]
\all{Wir isolieren das Dach, um Heizkosten zu sparen.} (Nous isolons le toit pour économiser des frais de chauffage.) \\
\all{Die Regierung baut Windräder, damit die Energiewende gelingt.} (Le gouvernement construit des éoliennes pour que la transition énergétique réussisse — sujets différents.) \\
\all{Ich ziehe den Stecker, damit das Gerät keinen Strom mehr verbraucht.} (Je débranche la fiche pour que l'appareil ne consomme plus d'électricité.) \\
\all{Viele Familien in Yaoundé nutzen Solarlampen, um bei Stromausfällen Licht zu haben.} (Beaucoup de familles à Yaoundé utilisent des lampes solaires pour avoir de la lumière pendant les coupures de courant.)
\end{exemplebox}

\begin{attention}[Erreurs fréquentes — Finalité]
\begin{itemize}
\item \textbf{Confusion \all{um ... zu} / \all{damit} :} si le sujet change $\rightarrow$ \all{damit} obligatoire. Si le sujet est identique $\rightarrow$ \all{um ... zu} (plus élégant).
\item \textbf{Oubli de \all{zu} :} \all{Ich schalte das Licht aus, um Strom zu sparen.} (Jamais *\all{um Strom sparen}.)
\item \textbf{Verbe séparable avec \all{um ... zu} :} le \all{zu} s'insère \emph{entre} la particule et le verbe : \all{um das Licht auszuschalten} (pour éteindre la lumière).
\item \textbf{Place du verbe avec \all{damit} :} le verbe conjugué \textbf{doit} être à la fin. \all{..., damit wir die Umwelt schützen.} (Jamais *\all{damit wir schützen die Umwelt}.)
\item \textbf{Traduction de « pour » :} ne pas traduire « pour » par \all{für} devant un verbe. \all{für} + nom = but nominal (\all{für die Umwelt}, pour l'environnement) ; \all{um ... zu} = but verbal (pour + infinitif).
\end{itemize}
\end{attention}

\courssub{Méthodes express}

\begin{methode}[Commenter un texte (5 étapes)]
\begin{enumerate}
\item \textbf{Lire globalement :} titre, source, images $\rightarrow$ thème général (ici : économiser l'énergie au quotidien).
\item \textbf{Lire en détail :} surligner les mots connus, deviner les inconnus par le contexte ou la transparence (ex. \all{das Wasserkraftwerk} = Wasser + Kraft + Werk : eau + force + usine).
\item \textbf{Structurer :} identifier les paragraphes = les étapes de l'argumentation (1. contexte, 2. gestes simples, 3. solutions techniques, 4. appel à l'action).
\item \textbf{Répondre aux questions :} repérer l'information dans le texte, reformuler en allemand simple (Présent ou Perfekt), en phrases complètes.
\item \textbf{Rédiger le résumé ou l'avis :} 5 à 6 phrases avec des connecteurs (\all{zuerst, außerdem, deshalb, schließlich}) ; employer \all{man} ou \all{wir} ; intégrer au moins un modal et une construction de finalité.
\end{enumerate}
\end{methode}

\begin{methode}[Traduire Thème/Version (3 étapes)]
\begin{enumerate}
\item \textbf{Analyser :} identifier le temps, les cas, la structure (modal ? finalité ?), les pièges de vocabulaire.
\item \textbf{Transposer :} appliquer les règles allemandes (verbe en 2\ieme{} position, infinitif à la fin, cas après les prépositions). Ne jamais calquer le français mot à mot.
\item \textbf{Vérifier :} majuscules des noms, Umlaute, \all{ß}, guillemets „…“, ordre des mots (TeKaMoLo), accord sujet-verbe.
\end{enumerate}
\end{methode}



\begin{aretenir}[Checklist avant le Bac — Leçon ALA11]
\begin{itemize}
\item $\square$ Je connais le vocabulaire de base : \all{die Energie, der Strom, sparen, verbrauchen, die Steckdose, der Stecker, das Kraftwerk, erneuerbar} (genre + pluriel).
\item $\square$ Je sais conjuguer les six \cle{Modalverben} au Présent sans faute (changement de voyelle au singulier, pas de \textit{-t} à la 3\ieme{} personne).
\item $\square$ Je place l'infinitif du verbe principal \textbf{à la fin} de la phrase avec un modal : \all{Ich will Energie sparen.}
\item $\square$ Je distingue \all{müssen} (obligation réelle) de \all{sollen} (conseil/ordre transmis) et \all{dürfen} (permission) de \all{können} (capacité).
\item $\square$ Je sais que \all{nicht müssen} = ne pas être obligé, alors que \all{nicht dürfen} = avoir l'interdiction.
\item $\square$ Je forme le Perfekt des modaux avec \all{haben} + \textbf{deux infinitifs} (Ersatzinfinitiv) : \all{Ich habe Strom sparen müssen.}
\item $\square$ J'utilise \all{um ... zu + Infinitif} quand le \textbf{sujet est le même} pour exprimer le but.
\item $\square$ J'utilise \all{damit + verbe conjugué à la fin} quand les \textbf{sujets sont différents} (ou pour insister).
\item $\square$ Avec un verbe séparable, je place le \all{zu} entre la particule et le verbe : \all{um das Licht auszuschalten}.
\item $\square$ Je ne confonds pas \all{für} (préposition + accusatif, but nominal) et \all{um ... zu} (but verbal).
\item $\square$ Je sais rédiger 5 phrases cohérentes sur mes habitudes énergétiques (modal + finalité + vocabulaire du thème).
\item $\square$ Je repère les \cle{verbes à particule séparable} du thème (\all{ausschalten, einschalten, runterdrehen}) et j'applique la règle (particule à la fin au Présent).
\item $\square$ Je relis mon écrit : majuscules aux noms, \all{ß} après voyelle longue ou diphtongue, ordre TeKaMoLo, verbe en 2\ieme{} position.
\end{itemize}
\end{aretenir}

\findecours

\end{document}
